% v5.1 (2026-07-14): THE ACCEPTANCE PASS (review 5.0.01, MAJOR REVISION -- the first non-reject; all 11 verified items executed).
%   .01 the completeness-certificate proposition retitled/restated at the licensed relative form (universal claim withdrawn);
%   .02 P4 tightness WITHDRAWN -- the reviewer's 54-DOF conjugate-pair two-doublet template passes every printed filter (verified in
%   exact arithmetic); the old 63-DOF witness is Witten-odd; P4 restated as class dominance (class minimum 54 > 45); RT4 description
%   corrected; .03 F3 renamed to its content (doublet-singlet co-presence), honest role stated, chirality gloss removed (the 54-DOF
%   template is vector-like yet passes); .04 beta primer rewritten in the Weyl convention b0 = (11/3)C_A - (2/3) Sum T(r), arithmetic
%   corrected, scalar term separated; .05 Theorem gauge_unique split: nonabelian ranking forward, abelian count exits to the
%   a-posteriori closure ("No additional abelian factor is needed" deleted from the forward statement); .06 C(G)=dim(G) statement
%   line now carries the operational premise pair; .07 "all 17 compact simple Lie algebras" -> families+exceptionals with the
%   17-representative tested list; J(N_c) family-monotonicity scope stated at the formula; N_c=3 statement scoped to the SM-family
%   template class; .08 hypercharge strong sites rewritten to the licensed ratios+convention form; .09 kinetic-mixing subsection
%   REMOVED (its positive-definiteness inequality is false -- counterexample verified; the surviving point is one sentence);
%   .10 neutrino/cosmology residue EXITED to the Papers 8/41 landing fragment (Majorana/0nubb prediction register, the >5-sigma
%   nu_R argument, the I2 proposition + Planck comparison); the in-core statement is the supported one: the minimum template in
%   the declared scan contains no light sterile singlet; .11 printed run command fixed (python verify_all.py) + code-concordance
%   attribution corrected (canonical 10/5 waterfall: check_L_F6_not_from_EC in ec_inventory_reading.py; check_T_field is the
%   uniform-dimension 4/2 robustness path, no per-stage counts). Prior changelog follows.
% v5.0 (2026-07-13): THE DUAL-SUPPLEMENT SPLIT (ruling of record, E. Brooke 2026-07-13: 'make it all Paper 2 still, just with two technical supplements').
%   This document continues the reviewed lineage (v3.9 -> v4.0 -> v4.1 -> v5.0) as TECHNICAL SUPPLEMENT I -- THE CLASSIFICATION CORE.
%   The foundational chain (algebra construction, gauge identification, matter category and lift, C(G)=dim(G) program, EC/CM grounds,
%   type-III/continuum discussion, imported-proofs appendix, reconstruction-program positioning) EXITS to the NEW Technical Supplement II
%   (Paper_2_Foundational_Gauge_Program_Supplement_v1.0). The cosmology tail ((61,102) architecture, sensitivity, Planck comparisons,
%   'Five predictions' block) EXITS to staged landing fragments for Papers 8/41. Review 4.1.01 fold-in executed per the closure tracker:
%   .01 conditional-classification premise ledger C1-C10 (exactness clause deleted; one-U(1) presented as a-posteriori closure theorems
%   with non-circularity framing); .02 honest completeness language at all three sites incl. the appendix; .07 composed C(G) statement
%   consumed with its premise pair named; .09 dynamical-residue clause; .10 F2-window typed as selected filter premise + retitle + 
%   highest-weight boundary; .11 canonical F6 = FULL-SYSTEM non-degeneracy (waterfall/counts regenerated 10/5; spectator-reduction 8/4
%   and uniform-dimension 4/2 as robustness with the T66 charged-octet witness); .12 cosmology tail exited (BA rows + exit note remain).
%   Items .03/.04/.05/.06/.08-chain-side executed in Technical Supplement II. Prior changelog follows.
\documentclass[11pt]{article}
\usepackage[letterpaper,margin=1in]{geometry}
\usepackage{amsmath,amssymb,amsthm,mathtools}
\usepackage{enumitem}
\usepackage{booktabs}
\usepackage{array}
\usepackage[colorlinks=true,linkcolor=black,citecolor=black,urlcolor=black]{hyperref}
\usepackage[most]{tcolorbox}
\tcbuselibrary{skins,breakable}
\usepackage{fancyhdr}
\usepackage{caption}

% ── Theorem environments ──
\newtheorem{theorem}{Theorem}[section]
\newtheorem{lemma}[theorem]{Lemma}
\newtheorem{proposition}[theorem]{Proposition}
\newtheorem{corollary}[theorem]{Corollary}
\newtheorem{theorem_app}{Theorem}[section]
\newtheorem{lemma_app}[theorem_app]{Lemma}
\theoremstyle{definition}
\newtheorem{definition}[theorem]{Definition}
\newtheorem{remark}[theorem]{Remark}
\newtheorem{assumption}[theorem]{Assumption}

% ── Box style (matches Paper 1 / Paper 2) ──
\tcbset{
    apfbox/.style={
        enhanced,
        colback=black!3,
        frame hidden,
        borderline={0.5pt}{0pt}{black},
        arc=0pt,
        left=8pt, right=8pt, top=8pt, bottom=8pt,
        fonttitle=\bfseries\color{black},
        coltitle=black,
        detach title,
        before upper={\tcbtitle\par\smallskip},
        before skip=14pt,
        after skip=14pt,
        grow to left by=-8pt,
        grow to right by=-8pt,
        sharp corners,
        breakable,
        pad after break=4pt,
        pad before break=4pt,
        lines before break=3,
    }
}

% ── Shortcuts ──
\newcommand{\eps}{\varepsilon}
\newcommand{\End}{\mathrm{End}}
\newcommand{\SU}{\mathrm{SU}}
\newcommand{\U}{\mathrm{U}}
\newcommand{\id}{\mathrm{id}}
\DeclareMathOperator{\spn}{span}
\DeclareMathOperator{\conv}{conv}
\DeclareMathOperator{\tr}{tr}
\DeclareMathOperator{\diag}{diag}
\newcommand{\coderef}[2]{\footnote{Code reference: \texttt{#1} in \texttt{#2}.}}
\newenvironment{varlist}{\begin{itemize}[nosep,leftmargin=2em]}{\end{itemize}}

\pagestyle{fancy}
\fancyhf{}
\fancyhead[L]{\small\itshape Paper 2 Technical Supplement I --- The Classification Core}
\fancyhead[R]{\small\itshape E.\,S.\ Brooke}
\fancyfoot[C]{\thepage}
\renewcommand{\headrulewidth}{0.4pt}


% ═══════════════════════════════════════════════════════════════
\setlength{\emergencystretch}{3em}%CORPUSGUARD
\usepackage{seqsplit}%CORPUSGUARD
\usepackage{graphicx}%CORPUSGUARD
\newcommand{\codeid}[1]{\texttt{\seqsplit{#1}}}%CORPUSGUARD

%CORPUSFMT 2026-07-04: preamble normalized to the Paper 0 corpus standard
\usepackage[T1]{fontenc}
\usepackage[utf8]{inputenc}
\usepackage{lmodern}
\usepackage{parskip}
\usepackage{titlesec}
\usepackage{microtype}
\titleformat{\part}[display]{\centering\Huge\bfseries}{\partname~\thepart}{1em}{}
\titleformat{\section}{\Large\bfseries}{\thesection.}{0.5em}{}
\titleformat{\subsection}{\large\bfseries}{\thesubsection}{0.5em}{}
\titleformat{\subsubsection}{\normalsize\bfseries}{\thesubsubsection}{0.5em}{}
\setlength{\emergencystretch}{1.5em}

\newcommand{\TSII}{Technical Supplement~II (\emph{The Foundational Gauge Program})}
\newcommand{\TSI}{Technical Supplement~I (\emph{The Classification Core})}

\date{Version 5.1 --- July 2026}
\begin{document}

\begin{center}
{\LARGE\bfseries Technical Supplement I --- The Classification Core\\[4pt]
The Structure of Admissible Physics}

\bigskip
{\large E.\,S.\ Brooke}\\[2pt]
{\small ORCID: \href{https://orcid.org/0009-0001-2261-4682}{0009-0001-2261-4682}}\\[4pt]
{\small Companion to Paper~2: ``The Structure of Admissible Physics''\\
Technical Supplement I of two; the foundational program is Technical Supplement~II}

\bigskip
{\small Version 5.1 --- July 2026}
\end{center}

\bigskip

\begin{abstract}
\noindent
This supplement is the classification core of Paper~2.  Its citable
result is a conditional classification theorem
(Theorem~\ref{thm:conditional_classification}) whose full input ledger
is part of the statement: given the declared template caps and
representation tables of the scan, the fixed gauge template
$\SU(3)\times\SU(2)\times\U(1)$, the imported generation count
$N_{\mathrm{gen}}=3$, the imported one-loop Weyl beta coefficient, the
anomaly and Witten equations as filter conditions, the F3
doublet--singlet content predicate, CPT identification,
electromagnetic completeness (F6, an
independent assumption in its canonical full-system non-degeneracy
form), the Weyl-component objective, and the equal-unit channel
convention, the Standard-Model five-multiplet template at 45 Weyl
degrees of freedom is the unique minimum among the enumerated
candidates, computed exactly.  The abelian sector is closed by
a-posteriori closure theorems for the selected template: given the
scanned matter content, at least one abelian grading is necessary
(under the named reading R-EC-inv with the entry-typed inventory
I-typing), one suffices, the anomaly system fixes its direction up to
normalization and relabeling, and no second factor is admissible under
capacity minimality.  The canonical F6 predicate is full-system
non-degeneracy (10 survivors pre-CPT, 5 post-CPT); spectator reduction
(8/4) and uniform dimension (4/2) are reported as robustness checks,
with the same winner under all three.  The capacity count
$45+4+12=61$ is proved in the equal-unit channel convention; every
cosmological reading of it is a bridge axiom (BA1--BA10,
\S\ref{sec:bridge_axioms}), and the cosmology-facing material lives at
its corpus homes.  Reproducibility is contractual: the scan is a
finite enumeration in exact rational arithmetic, machine-verified
in-tree and by a standalone script.  The premise list itself --- why
these inputs and not others, and at what grade each is grounded --- is
the object of a separately presented program (\TSII); this document
consumes the inputs and proves the classification.
\end{abstract}

\medskip
\begin{tcolorbox}[apfbox, title={What Is Assumed, What Is Proved, What Is Deferred}]

\textbf{Imported from Paper~1 Supplement.}  The four structurally
necessary components of PLEC --- A1 (finite capacity), MD (positive
cost floor), A2 (argmin selection), BW (non-degeneracy) --- together
with the constitutive admissibility semantics (SP, K3, FD1--FD6) and
the finite-carrier chain ($T_{\mathrm{form}}$, $T_{\mathrm{sep}}$,
$L_{\varepsilon^*}$, $L_\Delta$, $T_{\mathrm{alg}}$,
$T_{\mathrm{GNS}}$, T2a--T2c, $L_{\mathrm{loc}}$, $L_\omega$,
$T_{\mathrm{adj}}$, $L_{\mathrm{irr}}$).  Self-contained re-proofs of
the imported results live in the appendix of \TSII.

\smallskip\noindent\textbf{Consumed from \TSII\ (proved there at the
grades stated in its grade ledger).}
\begin{enumerate}[nosep,label=\arabic*.,leftmargin=2em]
  \item $L_{\mathrm{Cauchy}}$: cost-function uniqueness, $F(n)=n$; the
  channel unit for every cost comparison here.
  \item The finite semisimple admissibility algebra and the gauge
  identification: every continuous physical gauge redundancy is inner
  ($\mathrm{Gau}_0 \subseteq \mathrm{Inn}(\mathcal{A})$; equality
  under the named full-invariance hypothesis).
  \item The exact sequence $1 \to Z(\mathrm{U}(\mathcal{A})) \to
  \mathrm{U}(\mathcal{A}) \to \mathrm{Inn}(\mathcal{A}) \to 1$ and the
  relative-phase torus $T_{\mathrm{rel}}$ as the sole abelian source.
  \item The matter category on the lifted group $\widetilde{G}$, with
  the charge representation a named input (I-typing), and the lift and
  global $\mathbb{Z}_6$ quotient.
  \item EC at [P\textsubscript{structural\_reading} $\mid$ R-EC-inv
  $+$ I-typing] and CM at [P] (A2 over the A1$+$MD admissible set).
  \item $C(G) = \dim(G)$ under the operational cost premise pair
  (operational encoding; disjoint-anchor realizability).
  \item Theorem~R's carrier requirements R1--R3, $T_{7B}$,
  $L_{\mathrm{irr,uniform}}$, and $B1'$.
\end{enumerate}

\smallskip\noindent\textbf{Core results proved in this supplement.}
\begin{enumerate}[nosep,label=\arabic*.,leftmargin=2em]
  \item $L_{\mathrm{gauge,unique}}$: gauge template uniqueness within
  the declared carrier filters (\S\ref{sec:gauge_class_supp}).
  \item $T_{\mathrm{field}}$: fermion content by complete
  1{,}680-template scan of the declared space with seven filters
  (\S\ref{sec:fermion_scan}).
  \item Theorem~\ref{thm:conditional_classification}: the conditional
  classification theorem with its full input ledger (C1--C10) --- the
  citable core.
  \item The one-U(1) a-posteriori closure theorems
  (\S\ref{sec:one_u1_closure}).
  \item $L_{\mathrm{hypercharge}}$: hypercharge uniqueness from the
  anomaly system (\S\ref{sec:hypercharge}).
  \item $L_{\mathrm{count}}$: capacity counting, $45+4+12=61$
  (\S\ref{sec:capacity_count}).
  \item The asymptotic-freedom imported-coefficient check
  (\S\ref{sec:beta_schematic}).
\end{enumerate}

\smallskip\noindent\textbf{Exited at the v4.0 core split (bridge
material).}  Lorentzian signature, $d=4$, Einstein equations, horizon
equipartition, and saturation independence live at Papers~6, 8, and 41;
this supplement retains the exit map and the bridge-axiom ledger
(\S\ref{sec:bridge_axioms}).  The cosmology tail --- the (61,\,102)
architectural reading, the sensitivity table, and the Planck
comparisons --- is staged for Papers~8 and 41 (exit note in
\S\ref{sec:capacity_count}).

\smallskip\noindent\textbf{Not proved here (honest flags).}
\begin{enumerate}[nosep,label=\arabic*.,leftmargin=2em]
  \item Quantitative $\beta$-function derivation: only the imported
  coefficient check is here (\S\ref{sec:beta_schematic}); the
  capacity-level coefficient derivation is deferred to Paper~4A.
  \item Cap completeness: the enumeration caps of the scan are
  declared inputs; that no admissible sub-45-DOF template exists
  outside them is an open proof obligation
  (\S\ref{sec:exclusion_proofs}).
  \item Cosmological identifications: every cosmological reading of
  the 61 is a bridge axiom (BA5--BA7), not a result of this document.
  \item Mass hierarchy, CKM/PMNS mixing, CP violation: Papers~4A/4B.
  \item Confinement scale $\Lambda_{\mathrm{QCD}}$ and hadronization:
  Paper~4A.
\end{enumerate}
\end{tcolorbox}

\begin{tcolorbox}[apfbox, title={Imported-ingredient hypothesis ledger (verified vs.\ assumed in the finite-cost environment)}]
\footnotesize
Every off-framework mathematical import, with its hypothesis status
\emph{in this document's finite setting}.  ``Verified'' = the theorem's
hypotheses are checked here against the finite structures it is applied
to; ``assumed'' = the hypothesis is taken as a premise and flagged where
used.

\smallskip
\begin{tabular}{@{}p{0.30\linewidth}p{0.30\linewidth}p{0.32\linewidth}@{}}
\toprule
\textbf{Import} & \textbf{Applied to} & \textbf{Hypothesis status} \\
\midrule
Cartan closed-subgroup & classification class & verified (the class is $\prod_k \mathrm{PU}(n_k)$ and closed subgroups; scope rider travels with \codeid{check\_L\_cost\_gauge}) \\
Compact simple Lie classification & the F1--F3 table & imported whole (Killing/Cartan); carrier list extended to spinors in v4.0 \\
One-loop $\beta$ coefficient & F1/F2 filters & \emph{imported, assumed}: structural sign pressure only is derived (\S\ref{sec:beta_schematic}); scan runs at the Weyl $2/3$ convention \\
$N_{\mathrm{gen}} = 3$ & filter multiplicity & \emph{imported, assumed here}: banked capacity--beta identities (\S\ref{sec:three_generation_scope}); F1/F2 conditional on it \\
Electromagnetic completeness & F6 (canonical: full-system non-degeneracy) & \emph{independent assumption}, declared at point of use (\S\ref{sec:seven_filters}); logically independent of EC --- computed \\
Compact-$\U(1)$ representation data & hypercharge normalization & \emph{premise}: the abelian factor is compact with a chosen character lattice; the anomaly system fixes charge \emph{ratios}, and the absolute unit ($Y_Q = 1/6$) is this convention (\S\ref{sec:hypercharge}) \\
HKM / Malament / Kolmogorov / Lovelock & exited bridge & assumed as bridge axioms BA1--BA4 (\S\ref{sec:bridge_axioms}); not consumed by the core \\
F2 rank window ($3 \le k \le \mathrm{rank}(G){+}1$) & the (F2) carrier filter & \emph{selected filter premise} (F2-window): a declared physical filter --- low-rank antisymmetric primitive invariants, few-body composite realizability (\S\ref{sec:R1_classification}) \\
Operational cost premise pair & $C(G) = \dim(G)$ (cost ranking) & \emph{named operational premises}: operational encoding (lower bound) $+$ disjoint-anchor realizability (upper bound); Proposition~\ref{prop:generator_separation} \\
\bottomrule
\end{tabular}
\end{tcolorbox}
\begin{tcolorbox}[apfbox, title={Terminology: four words with one fixed meaning each}]
\footnotesize
\begin{description}[nosep,leftmargin=1.6em]
  \item[faithful] of a representation or state: trivial kernel.  For
  states: $\omega(a^*a) = 0 \Rightarrow a = 0$ (GNS null ideal zero).
  Never used to mean ``physically realized.''
  \item[irreducible] of a representation: no proper nonzero invariant
  subspace.  Used only rep-theoretically; ``primitive'' is used for
  invariant tensors, never ``irreducible.''
  \item[effective] of a group action: only the identity acts trivially
  (synonym of faithful for actions); used when the group, rather than
  its representation, is the subject.
  \item[physically distinct (multiplet types)] distinct isomorphism
  classes of $\widetilde{G}$-representations
  (Definition~2.20 of \TSII{}); copies within a class
  (generations) are the same type at different multiplicity.
\end{description}
\end{tcolorbox}

\medskip
\begin{tcolorbox}[apfbox, title={Reader's map: one paper, two supplements}]
\footnotesize
Paper~2's technical material lives in two supplements.  This document
--- Technical Supplement~I, \emph{The Classification Core} --- carries
finished mathematics on declared inputs: the conditional
classification theorem with its full premise ledger, the scan and its
waterfall, anomaly and hypercharge closure, the one-U(1) a-posteriori
closure theorems, the Lie classification within the declared carrier
filters, the capacity count, and the bridge-axiom ledger.  Nothing in
this document claims a derivation from A1/MD/A2/BW.  \TSII\ carries
the foundational program --- the chain from the APF axioms toward this
document's inputs, each link at its banked grade, with its grade
ledger and open-lane register.  A reader auditing the classification
needs this document and the premise ledger; a reader auditing the
premises' grounds needs Technical Supplement~II.
\end{tcolorbox}

\begin{tcolorbox}[apfbox, title={Framing convention: ``forces'' and ``selects'' denote uniqueness, not processes}]
\footnotesize
The supplement's theorems and proofs use the verbs ``forces,'' ``selects,'' and ``rejects'' as standard mathematical shorthand for uniqueness results, $\arg\min$ properties, and inadmissibility results respectively.  Under the APF ontology, these verbs do not denote processes: nothing is forced by an agent, nothing is selected by a mechanism, nothing is rejected by a filter.  Each theorem asserts a property of the set of admissible configurations---that it has a unique member, a minimum member, or no member satisfying some condition.  For example, ``capacity minimality selects $N_c = 3$'' is shorthand for ``$N_c = 3$ is the $\arg\min$ of dim(G) over R1--R3-admissible color groups,'' and the alternative values (SU(4), SU(5), E$_6$, \ldots) are \emph{admissible} in the sense that each has cost below the capacity bound, but have higher cost than SU(3) and are therefore not what is.  The companion note \emph{What Is, Is Minimum-Cost Admissible} spells this out in full.  Readers may substitute the descriptive reading throughout the supplement without any change to the proofs.

\begin{tcolorbox}[apfbox, title={Status of A2: argmin realization is a lawlike selection rule, not a derived dynamics}]
\footnotesize
A2 is intentionally not derived in this supplement from a deeper dynamical variational principle.  It is one of the four structural inputs of the finite-regime APF package: among admissible continuations satisfying A1, MD, and BW, the realized carrier is the minimum admissibility-cost representative.  All selection claims in the supplement have the conditional form
\[
  \text{A1+MD+BW+A2+listed interface hypotheses} \;\Longrightarrow\; \text{unique minimum-cost admissible structure}.
\]
Thus the paper does not claim to prove that nature \emph{must} obey global argmin realization from conventional QFT assumptions.  It claims that, if A2 is accepted as the APF lawlike selection rule, the Standard-Model template is the unique admissible minimum under the stated finite ledger.  This is why counterfactuals that change A2, make new generators cost-free, or subtract structural gauge generators from the ledger are not small variants of the theorem; they define a different selection problem.
\end{tcolorbox}
\end{tcolorbox}

\medskip
\smallskip\noindent\textit{Self-contained Standard Model datum.}\quad
Section~\ref{sec:hypercharge} derives the unique anomaly-free hypercharge
assignments for one generation of the Standard Model fermion template
purely from the tools assembled in this supplement (gauge template from
\S\ref{sec:gauge_class_supp}, fermion content from \S\ref{sec:fermion_scan},
anomaly conditions from \S\ref{sec:hypercharge}).  The result
($Y_Q = 1/6$, $Y_u = 2/3$, $Y_d = -1/3$, $Y_L = -1/2$, $Y_e = -1$)
constitutes a nontrivial, experimentally verified prediction of the
derivation chain: the charge \emph{ratios} and their quantization are
derived, while the absolute unit is fixed by the compact-$\U(1)$
charge-lattice convention (front-matter imported-ingredient ledger).



\medskip\noindent\textbf{Dependency table.}  Each row gives a result proved in this supplement, its inputs, and the downstream results that depend on it.

\smallskip
\begin{center}\footnotesize
\renewcommand{\arraystretch}{1.15}
\begin{tabular}{@{}lp{4.9cm}p{4.3cm}@{}}
\toprule
\textbf{Theorem} & \textbf{Inputs} & \textbf{Used by} \\
\midrule
$L_{\mathrm{gauge,unique}}$ (\S\ref{sec:gauge_class_supp}) & Theorem R, $L_{\mathrm{Cauchy}}$, Lie classif.\ (all consumed from \TSII\ / imported) & $T_{\mathrm{field}}$, $L_{\mathrm{count}}$ \\
$T_{\mathrm{field}}$ (\S\ref{sec:fermion_scan}) & C1--C10 ledger (Thm.~\ref{thm:conditional_classification}) & closure theorems, $L_{\mathrm{count}}$ \\
One-U(1) closure (\S\ref{sec:one_u1_closure}) & $T_{\mathrm{field}}$ output; EC/CM at their grades & abelian-sector consistency \\
$L_{\mathrm{hypercharge}}$ (\S\ref{sec:hypercharge}) & $T_{\mathrm{field}}$, anomaly conditions & closure theorems, $L_{\mathrm{count}}$ \\
$L_{\mathrm{count}}$ (\S\ref{sec:capacity_count}) & $T_{\mathrm{field}}$, $T_{\mathrm{Higgs}}$, $L_{\mathrm{gauge,unique}}$, BA5 convention & $I2$; exited bridge readings \\
\bottomrule
\end{tabular}
\end{center}

\tableofcontents

\medskip
\noindent\textit{On operational language.}\quad
Operational verbs are local readings of an underlying admissibility-space
structure under a temporal slicing.  They are retained because they are the
working vocabulary of physics; the structural reading is always available:
``acts,'' ``forms,'' ``selects,'' and ``evolves'' mean that the corresponding
admissible configuration, quotient, minimum, or continuation relation exists.
This convention is inherited from Paper~0 and Paper~1 Supplement v8.24.


\newpage

\part{Classification Inputs and Conditional Bridges}
\label{part:core}


\section{Consumed inputs: the foundational chain at a glance}
\label{sec:consumed_inputs}

This document proves the classification and closure results and
consumes the foundational chain as cited premises.  The chain itself
--- statements, proofs, grades, and named premises --- is \TSII;
nothing here claims a derivation from the APF axioms A1/MD/A2/BW.  The
box below states what is consumed and at what grade.

\begin{tcolorbox}[apfbox, title={Chain outputs consumed by this document}]
\footnotesize
\begin{enumerate}[nosep,leftmargin=1.8em]
  \item \textbf{The cost unit.}  $L_{\mathrm{Cauchy}}$: $F(n) = n$ on
  the discrete channel domain; every cost comparison in this document
  is a channel count.  Grade [P].
  \item \textbf{The algebra and the gauge identification.}  A finite
  semisimple admissibility algebra $\mathcal{A} \cong \bigoplus_k
  M_{n_k}(\mathbb{C})$ (standard $C^*$ route, given the IJC occupancy,
  which is empirical), with every continuous physical gauge redundancy
  inner: $\mathrm{Gau}_0 \subseteq \mathrm{Inn}(\mathcal{A})$, and
  equality under the named full-invariance hypothesis.  The SM
  specialization $\{n_k\} = \{3,2,1\}$ enters this document as
  hypothesis H1 of the closure theorems, not as a theorem of this
  document.
  \item \textbf{The exact sequence and the abelian source.}
  $1 \to Z(\mathrm{U}(\mathcal{A})) \to \mathrm{U}(\mathcal{A}) \to
  \mathrm{Inn}(\mathcal{A}) \to 1$; continuous abelian gauge freedom
  lives in the relative-phase torus $T_{\mathrm{rel}}$, physical
  exactly where charged matter sees it.
  \item \textbf{EC and CM at their grades.}  Enforcement completeness
  at [P\textsubscript{structural\_reading} $\mid$ R-EC-inv $+$
  I-typing]\coderef{check\_L\_EC\_inventory\_reading}{ec\_inventory\_reading.py};
  capacity minimality at [P] (A2 over the A1$+$MD admissible set,
  stated at channel-configuration level).
  \item \textbf{The gauge-cost equality.}  $C(G) = \dim(G)$ under the
  operational cost premise pair --- operational encoding for the lower
  bound, disjoint-anchor realizability for the upper bound --- realized
  in Proposition~\ref{prop:generator_separation}\coderef{check\_L\_cost\_gauge}{core.py}.
  \item \textbf{Carrier requirements and the matter category.}
  Theorem~R's carriers (complex of dimension $\ge 3$ with a low-rank
  antisymmetric primitive invariant; pseudoreal of dimension 2; a
  single abelian grading) and the matter category on the lifted group
  $\widetilde{G}$, with the charge representation a named input
  (I-typing).
\end{enumerate}
The Sep/IJC branch material is consumed the same way: the bridge from
superadditivity to a noncommutative algebra is derived given the IJC
occupancy, and the occupancy is empirical --- read off the record, not
derived from A1.  The grade-ledger table near the front of \TSII\ is
the one-page map of all of the above.
\end{tcolorbox}

\section{Asymptotic freedom: structural sign heuristic and imported coefficient check}
\label{sec:beta_schematic}
% ══════════════════════════════════════════════════════════════

Paper~2 \S3.6--3.7 derives the structural necessity of asymptotic freedom
from the non-abelian self-interaction of admissibility channels: gauge
bosons carry charge, producing anti-screening that overwhelms matter
screening.  The full quantitative derivation (one-loop $\beta$-coefficient
from Casimir arithmetic) is in Paper~4A ($L_{\mathrm{beta,capacity}}$).
This section provides the schematic sign calculation using only the
matter content derived in this paper.

\textit{Proof boundary.}\quad The one-loop $\beta$-function formula
is a structural result of gauge theory.  Paper~4A derives it from
capacity arguments ($L_{\mathrm{beta,capacity}}$); here we use it as
an imported formula (verifiable in any QFT textbook) and evaluate it
with the derived matter content.

The one-loop $\beta$-function coefficient for a gauge group factor
$G_i$ with adjoint Casimir $C_A$ is stated once, in the Weyl
convention used everywhere in this document (the scan's F1--F2
filters included):
\begin{equation}
  b_0 \;=\; \frac{11}{3}\,C_A \;-\; \frac{2}{3}\sum_{r}\,T(r),
  \label{eq:beta_one_loop}
\end{equation}
where the sum runs over the left-handed Weyl fermion species charged
under $G_i$ and $T(r)$ is the Dynkin index of the representation each
carries.  Complex scalars contribute a separate term
$-\tfrac{1}{3}\sum_{s} T(s)$; we keep it out of the fermion displays
below and account for it once, at the end.  The textbook Dirac form
$b_0 = (11/3)C_A - (4/3)\,T_R\,N_f$, with $N_f$ counting Dirac
flavors, is the same number repackaged (one Dirac flavor is two Weyl
species); mixing the two conventions --- the Dirac coefficient on a
Weyl count --- is exactly the mispricing the F1 filter statement
warns against (\S\ref{sec:seven_filters}).
Asymptotic freedom holds when $b_0 > 0$: the gauge self-interaction
(the $C_A$ term, from non-abelian structure) overwhelms the matter
screening (the $\sum T$ term).

\begin{proposition}[Sign of $b_0$ for the derived matter content]
\label{prop:AF_sign}
With the SM fermion content derived in $T_{\mathrm{field}}$
(\S\ref{sec:fermion_scan}), both non-abelian factors are asymptotically
free:
\begin{align}
  b_0^{\SU(3)} &= \frac{11}{3} \cdot 3 - \frac{2}{3} \cdot 6
    = 11 - 4 = 7 > 0, \label{eq:b0_su3} \\
  b_0^{\SU(2)} &= \frac{11}{3} \cdot 2 - \frac{2}{3} \cdot 6
    = \frac{22}{3} - 4 = \frac{10}{3} > 0, \label{eq:b0_su2}
\end{align}
with $\sum_r T(r) = 6$ for each factor, computed in the proof.
\end{proposition}

\begin{proof}
For $\SU(3)$: $C_A = 3$.  Each generation carries four colored Weyl
species --- the two $\SU(2)$ components of $Q$ (each a color
fundamental) plus $u^c$ and $d^c$ --- so $\sum_r T(r)$ per generation
is $4 \times \tfrac{1}{2} = 2$, and over three generations
$\sum_r T(r) = 6$.  Therefore
$b_0^{\SU(3)} = 11 - (2/3)(6) = 11 - 4 = 7$.

For $\SU(2)$: $C_A = 2$.  Each generation carries four Weyl doublets
--- $Q$ in three colors plus $L$ --- so $\sum_r T(r)$ per generation
is $4 \times \tfrac{1}{2} = 2$, and over three generations
$\sum_r T(r) = 6$: the matter contribution is $(2/3)(1/2)(12) = 4$.
Therefore $b_0^{\SU(2)} = 22/3 - 4 = 10/3$.

The scalar term is separate and does not change either sign: the
single complex Higgs doublet contributes $-\tfrac{1}{3} \cdot
\tfrac{1}{2} = -\tfrac{1}{6}$ to $b_0^{\SU(2)}$ (giving $19/6 > 0$
with the fermions) and nothing to $b_0^{\SU(3)}$.

Both coefficients are positive.  The signs depend only on the matter
content (derived in $T_{\mathrm{field}}$) and the group theory
(Casimirs of $\SU(3)$ and $\SU(2)$), not on any free parameter.
\end{proof}

\noindent\textit{What this establishes and what it does not.}\quad
The sign of $b_0$ determines the qualitative running: positive $b_0$
means the coupling decreases at high energy (asymptotic freedom) and
grows at low energy (confinement regime for $\SU(3)$).  This section
does \emph{not} derive the formula~(\ref{eq:beta_one_loop}) from
capacity arguments --- that derivation is in Paper~4A
($L_{\mathrm{beta,capacity}}$).  Here the formula is used as an
imported result (standard in any QFT textbook~\cite{PeskinSchroeder1995})
and evaluated with the derived matter content.  The structural claim
of Paper~2 --- that non-abelian self-interaction generically produces
asymptotic freedom when the matter content is sufficiently small ---
is confirmed by this explicit calculation.
\coderef{check\_T\_confinement}{gauge.py}


\subsection{Comparison with established $\beta$-function derivations}
\label{sec:AF_comparison}

The one-loop $\beta$-function coefficient
$b_0 = (11/3)C_A - (4/3)T_R N_f$ has been derived by at least four
independent methods in the literature.  This subsection identifies each,
states what the APF reuses vs.\ reinterprets, and clarifies the
boundary between established results and the APF's contribution.

\smallskip\noindent\textbf{1.\ Feynman-diagram computation
(Gross--Wilczek, Politzer, 1973).}\quad
The original derivation computes the one-loop vacuum polarization of
the gluon propagator via three diagrams: the fermion loop (screening,
$-4/3\, T_R N_f$), the ghost loop (antiscreening, $+1/3\, C_A$), and
the gluon self-energy with a three-gluon vertex (antiscreening,
$+10/3\, C_A$).  The sum gives $b_0 = (11/3)C_A - (4/3)T_R N_f$.

\textit{What the APF reuses:} the numerical result
$b_0 = (11/3)C_A - (4/3)T_R N_f$.  This is imported as an established
formula (Proposition~\ref{prop:AF_sign}) and evaluated with the
derived matter content.  The APF does not re-derive this formula in
the present paper.

\textit{What the APF reinterprets:} the physical mechanism.  In the
standard derivation, the $11/3$ arises from the sum of gluon and ghost
loop contributions; the ghost loop is a gauge-fixing artifact with no
direct physical interpretation.  In the APF's capacity language
(\S4.1 of \TSII{}), the antiscreening is attributed to the
\emph{superadditive cost of co-located non-abelian channels}
($L_\Delta > 0$), which is a physical statement about admissibility
economics, not a gauge-fixing artifact.  The ghost contribution is
absorbed into the capacity self-cost.

\textit{What is genuinely new:} the sign of $b_0$ is evaluated with
no free choice left in the matter input.  In standard QCD, the sign
of $b_0$ depends on the number of quark flavors $N_f$, which is a
free input.  Here $N_f$ is fixed within the ledger
($T_{\mathrm{field}}$: 6 flavors from the scan, under the C1--C10
premises), and the sign follows by evaluating the imported one-loop
formula with those ledger-fixed inputs --- sign pressure with imported
weights, exactly the banked check's register.  The novelty is not in
the formula but in the fact that its inputs are the ledger's, not
free choices.

\smallskip\noindent\textbf{2.\ Background-field method
(Abbott, 1981; 't~Hooft, 1976).}\quad
The background-field method computes the $\beta$-function by expanding
around a classical background $\bar{A}_\mu$ and integrating out
quantum fluctuations.  The result is gauge-invariant at each step,
making the antiscreening origin more transparent: the three
polarization states of a massive spin-1 particle contribute
$+(11/3)C_A$ (two physical transverse polarizations contribute
$+2C_A$, and the longitudinal / temporal sector contributes
$+(5/3)C_A$ after ghost subtraction).

\textit{APF relationship:} the background-field decomposition has a
direct capacity analogue.  The ``background field'' corresponds to
the admissibility structure at a fixed interface (the cost kernel $g$);
the ``quantum fluctuations'' are perturbations of the distinction
structure.  The $(11/3)C_A$ counts the capacity cost of maintaining $p$
self-interacting channels in $d = 4$ spacetime dimensions: two
physical polarizations (from the transversality condition, which is
the admissibility analogue of gauge fixing) plus the longitudinal
capacity overhead.  Paper~4A makes this correspondence precise.

\smallskip\noindent\textbf{3.\ Index-theoretic derivation
(Atiyah--Singer, heat-kernel methods).}\quad
The $\beta$-function can be extracted from the heat-kernel expansion
of the functional determinant $\det(-D^2)$, where $D$ is the gauge
covariant derivative.  The Seeley--DeWitt coefficient $a_2(D^2)$
gives the one-loop divergence, and its group-theoretic content
reproduces $b_0$.  This approach makes no reference to individual
Feynman diagrams; the result follows from the spectral geometry of
the gauge-covariant Laplacian.

\textit{APF relationship:} the index-theoretic derivation is the
closest mathematical relative of the APF's capacity approach.  Both
compute a spectral invariant of the gauge connection: the
Seeley--DeWitt coefficient is a spectral invariant of $-D^2$, while
the APF's capacity cost is a spectral invariant of the admissibility
algebra's cost kernel $B$ (via the eigenvalue decomposition in
Lemma~2.22 of \TSII{}).  The capacity inequality
(eq.~6 of \TSII{}), $p \cdot C_A > T_R \cdot N_f^{\mathrm{total}}$,
is a statement about the spectral content of the gauge sector
vs.\ the matter sector, analogous to the comparison of
$a_2$ contributions in the heat-kernel expansion.  The precise
coefficient $11/3$ corresponds to the $a_2$ coefficient of the
spin-1 Laplacian in $d = 4$, which Paper~4A derives from
the polarization count (the same $d = 4$ that the APF derives in
the exited $d{=}4$ argument, Paper~6; \S\ref{sec:bridge_axioms}).

\smallskip\noindent\textbf{4.\ Twistor-space methods
(Penrose, 1967; Witten, 2004).}\quad
Witten's reformulation of perturbative gauge theory in twistor space
computes scattering amplitudes without Feynman diagrams, using
holomorphic curves in $\mathbb{CP}^3$.  The one-loop $\beta$-function
can be extracted from the holomorphic anomaly of the twistor string.
The result agrees with the standard computation but organizes the
contributions differently: the antiscreening is attributed to the
self-dual sector of the gauge field, and the screening to the
matter multiplets' contribution to the holomorphic anomaly.

\textit{APF relationship:} the twistor approach is structurally
different from the APF (twistor space is infinite-dimensional and
uses complex geometry; the APF works with finite-dimensional
admissibility algebras).  However, both share a feature: they derive
the $\beta$-function sign from \emph{structural} properties of the
gauge field (self-duality in the twistor case, self-interaction cost
in the APF case) rather than from individual diagram counting.  The
APF does not use twistor methods; the comparison is noted to show
that the impulse to derive $b_0$ from structural principles rather
than diagram-by-diagram computation has precedent.

\smallskip\noindent\textbf{Summary: new vs.\ reinterpreted.}\quad

\begin{center}\small
\renewcommand{\arraystretch}{1.2}
\begin{tabular}{@{}p{4.5cm}p{1.8cm}p{6.5cm}@{}}
\toprule
\textbf{Element} & \textbf{Status} & \textbf{Detail} \\
\midrule
Formula $b_0 = (11/3)C_A - (4/3)T_R N_f$ &
  Imported &
  Standard one-loop result (Gross--Wilczek--Politzer, 1973);
  not rederived here. \\[3pt]
Sign $b_0 > 0$ for SM content &
  Imported + derived inputs &
  Formula imported; matter content $N_f$ derived from
  $T_{\mathrm{field}}$.  The sign is therefore a \emph{prediction}
  (given A1), not a parameter choice. \\[3pt]
Capacity inequality $p C_A > T_R N_f^{\mathrm{total}}$ &
  New &
  Derived from $L_\Delta$ (superadditivity) and
  Lemma~2.22 of \TSII{} (generator--channel bridge).
  Provides a sufficient condition for $b_0 > 0$ without the
  formula. \\[3pt]
Antiscreening from $L_\Delta$ &
  New interpretation &
  The superadditive cost of co-located non-abelian channels
  provides a physical mechanism for the $C_A$ term;
  reinterprets the ghost + gluon-loop sum as admissibility
  self-cost. \\[3pt]
Precise coefficients $11/3$, $4/3$ &
  Deferred &
  Paper~4A ($L_{\mathrm{beta,capacity}}$) derives these from
  polarization counting in $d = 4$ and the spin-statistics
  connection.  Not available in this supplement. \\[3pt]
Index-theoretic parallel &
  Noted &
  The capacity cost is a spectral invariant of $B$, analogous
  to the Seeley--DeWitt $a_2$ coefficient.  Not developed
  further here. \\
\bottomrule
\end{tabular}
\end{center}

\smallskip\noindent The honest assessment is: the APF's
contribution to the $\beta$-function story in this supplement is
\emph{modest}.  The formula is imported; the sign is verified with
derived inputs; the capacity inequality provides a new sufficient
condition but does not reproduce the precise coefficients.  The
genuine novelty is that $N_f$ (and hence the sign of $b_0$) is
derived rather than assumed.  The reinterpretation of antiscreening
as admissibility self-cost is a new physical picture but does not yet
produce quantitative predictions beyond what the standard formula
gives.  Paper~4A is where the capacity-based derivation must stand
or fall.


\newpage

% ══════════════════════════════════════════════════════════════


\noindent\textit{Paper 4A coefficient program.}\quad
The deferred coefficient derivation has three explicit steps: (i) construct
the capacity-level analogue of a Wilsonian shell by resolving which
continuation channels become active under refinement; (ii) prove that the
adjoint self-correlation contribution is weighted by the same channel-count
invariant as $C_A$, while a matter channel is weighted by the finite-index
embedding invariant corresponding to $T_R$; and (iii) show regulator
independence by proving that refinements with the same representation index
produce the same shell increment.  Already testable intermediate results are
therefore the signs, the $C_A/T_R$ weighting pattern, and the fact that
adding enough charged matter flips the sign exactly as the standard
coefficient predicts.  This supplement imports the coefficient but records
these as the Paper~4A proof obligations.


% ══════════════════════════════════════════════════════════════
%  v4.0 CORE SPLIT: the spacetime/cosmology/structural-prediction
%  material exited to its corpus homes (Papers 6, 8, 41, 45, 46).
%  This section is the exit map: the conditional-bridge theorem,
%  the bridge-axiom ledger, and one pointer per exited territory.
% ══════════════════════════════════════════════════════════════
\section{Conditional bridges: the exit map and bridge axioms}
\label{sec:bridge_axioms}

This supplement's claim surface ends with the gauge/matter template and its
capacity accounting.  Everything downstream of that surface --- causal order,
continuum spacetime, Lorentzian signature, Einstein dynamics, horizon
thermodynamics, cosmological density fractions, and the structural
predictions (vacuum stability, proton stability, strong CP) --- is a
\emph{bridge}: a conditional extension whose hypotheses are not theorems of
this document.  Earlier versions carried the bridge derivations in full;
they now live with their corpus homes (Papers~6, 8, 41, 45, and 46), and
what remains here is the boundary statement and the explicit list of bridge
axioms, each with its named import and its failure mode.

\begin{theorem}[Conditional spacetime bridge, not an input to gauge recovery]
\label{thm:conditional_spacetime_bridge}
The finite gauge/matter results of this supplement do not depend on the
Lorentzian, $d=4$, Einstein-equation, horizon, or cosmological-density
claims.  Those claims form a downstream bridge: if the positive cost kernel,
irreversible ordering, locality, and the imported HKM/Malament/Lovelock
hypotheses are all satisfied, then the cost-geometry interpretation follows.
If any bridge hypothesis fails, the gauge/matter derivation remains intact
and only the spacetime/cosmology interpretation is downgraded.
\end{theorem}

\noindent\textit{Why this is one-way.}\quad The gauge/matter derivation uses
only the finite admissibility ledger, Sep/IJC routing, the finite
$*$-algebra carrier, inner automorphisms, the classification of compact Lie
algebras, and the finite template scan.  None of those steps invokes a
continuum manifold, Lorentzian signature, Lovelock dynamics, or horizon
thermodynamics.  The bridge consumes the derived cost kernel and
irreversibility order as inputs; nothing flows back.

\subsection{The bridge-axiom ledger}
\label{sec:bridge_axiom_ledger}

Each axiom below is a premise the bridge \emph{assumes}.  None is derived
from A1, MD, A2, or BW in this document.  A reader who rejects a bridge
axiom loses exactly the row's conclusion and nothing upstream.

\begin{enumerate}[label=\textbf{BA\arabic*},leftmargin=3em]
  \item \textbf{Causal order.}  The commitment relation $\prec$ is
  irreflexive and transitive.  \emph{Import:} order axioms --- irreversible
  positive commitment alone does not prove irreflexivity (applying the
  identity is not an undoing operation), and transitivity needs a
  composition-monotonicity premise.  \emph{Failure mode:} $\prec$ admits
  cycles or fails to compose; no causal set.  \emph{Home:} Paper~6.
  \item \textbf{Continuum structure.}  The event set carries a $C^2$
  differentiable manifold structure compatible with the projective-limit
  measure.  \emph{Import:} the manifold itself --- Kolmogorov extension
  yields a measure on a product-limit space, not charts; tangent spaces,
  exponential maps, and $C^2$ regularity are assumed with it.
  \emph{Failure mode:} the continuum description has no differentiable
  refinement; the metric chain stops at the order.  \emph{Home:} Paper~6.
  \item \textbf{Lorentzian reconstruction hypotheses.}  $(M, \prec)$
  satisfies the Hawking--King--McCarthy/Malament distinguishing conditions.
  \emph{Import:} HKM 1976 / Malament 1977 hypotheses.  \emph{Failure mode:}
  the order is not realizable as the causal order of any smooth Lorentzian
  manifold.  \emph{Home:} Paper~6.
  \item \textbf{Second-order dynamics.}  Locality, diffeomorphism
  covariance, conservation, and second-order metric dependence --- the
  Lovelock premises.  \emph{Import:} Lovelock 1971 hypotheses; a finite
  information bound does not by itself limit a continuum response law to
  second derivatives.  \emph{Failure mode:} higher-derivative or nonlocal
  gravitational dynamics.  \emph{Home:} Paper~6.
  \item \textbf{Common capacity unit.}  Each Weyl fermion component, each
  real Higgs component, and each gauge generator occupies exactly one
  capacity slot of equal weight, so that $45 + 4 + 12 = 61$ is a count in
  a single currency.  \emph{Import:} the unit-equivalence premise.
  \emph{Failure mode:} sector-dependent slot weights; the 61 loses its
  privileged status.  \emph{Home:} Paper~8.
  \item \textbf{Cosmological identification.}  Capacity-slot fractions are
  Friedmann density parameters ($\Omega_\Lambda = 42/61$ and kin).
  \emph{Import:} a phenomenological identification (a reading, not a
  theorem); maximizing entropy over 61 equal slots gives a uniform
  distribution over \emph{slots}, and the step from slot fractions to
  cosmological densities is an ansatz.  \emph{Failure mode:} the fractions
  are bookkeeping only.  \emph{Home:} Paper~8 (dark-sector reading:
  Paper~35).
  \item \textbf{Horizon equipartition conditions.}  The maximum-entropy
  equal-priors conditions under which the horizon ledger equidistributes.
  \emph{Import:} the equipartition assumption set of the exited horizon
  section.  \emph{Failure mode:} weighted priors; the equipartition
  fractions shift.  \emph{Home:} Paper~41.
  \item \textbf{Saturation conditionality (baryon number).}  The
  proton-stability conclusion is conditional on full Bekenstein saturation
  \emph{and} on the completeness of the derived particle content: within
  that content no mediator generates the $\Delta B = 1$ dimension-six
  operators, but those operators are gauge-invariant under
  $\mathrm{SU}(3){\times}\mathrm{SU}(2){\times}\mathrm{U}(1)$ as EFT
  structures, and ultraviolet degrees of freedom outside the derived
  content could source them.  The honest claim is \emph{no generation
  within the derived content}, not non-existence of the operators.
  \emph{Home:} bank checks \codeid{check\_T\_proton},
  \codeid{check\_L\_B\_operator\_exclusion} (\codeid{gauge.py}).
  \item \textbf{Vacuum bridge.}  The exited vacuum-stability argument
  shows uniqueness of the minimum of the APF potential under its own
  assumptions; it does not by itself constrain the renormalization-group-improved
  Standard Model effective potential.  \emph{Home:} Paper~45.
  \item \textbf{Strong CP mechanism.}  $\theta_{\mathrm{QCD}} = 0$ is a
  proposed mechanism carried at [P-heuristic]; the admissibility principles
  stated here do not derive the absence or irrelevance of the
  $\theta$-term.  \emph{Home:} Paper~46.
\end{enumerate}

\subsection{Exit map}
\label{sec:exit_map}

\begin{center}
\small
\begin{tabular}{@{}llll@{}}
\toprule
\textbf{Territory} & \textbf{Was (v3.9)} & \textbf{Home} & \textbf{Bridge axioms} \\
\midrule
Causal order $\to$ Lorentzian $\to$ $d{=}4$ $\to$ Einstein & \S9 & Paper~6 + Supp. & BA1--BA4 \\
MECE partition ($61 = 42 + 3 + 16$) & \S17 & Paper~8 & BA5 \\
Cosmological fractions ($\Omega_\Lambda = 42/61$) & \S17 & Paper~8 / 35 & BA5, BA6 \\
Horizon equipartition ($L_{\mathrm{equip}}$) & \S18 & Paper~41 & BA7 \\
Saturation independence ($L_{\mathrm{sat,part}}$) & \S19 & Paper~8 & BA5 \\
Vacuum stability & \S20 & Paper~45 & BA9 \\
Proton stability & \S20 & bank (\codeid{gauge.py}) & BA8 \\
Strong CP & \S20 & Paper~46 & BA10 \\
\bottomrule
\end{tabular}
\end{center}

\smallskip\noindent The bank checks for the exited results
(\codeid{check\_Delta\_signature}, \codeid{check\_T8},
\codeid{check\_T9\_grav}, \codeid{check\_P\_exhaust},
\codeid{check\_L\_equip}, \codeid{check\_L\_saturation\_partition},
\codeid{check\_T\_vacuum\_stability}, \codeid{check\_T\_theta\_QCD},
\codeid{check\_T\_proton}) remain banked and machine-verified; this
supplement simply no longer presents their derivations as its own claim
surface.

\part{Gauge Template and Matter Content}\label{part:gauge_matter}
\label{part:gauge}


% ══════════════════════════════════════════════════════════════
\section{Gauge template uniqueness: Lie classification within the declared carrier filters}
\label{sec:gauge_class_supp}
% ══════════════════════════════════════════════════════════════

Paper~2 \S5.1 states $L_{\mathrm{gauge,unique}}$ in compact form.
This section provides the full classification with explicit exclusion
reasons across the four infinite families and five exceptional
algebras (the bank check tests a 17-representative list; the families
are handled by closed-form monotonicity), quantified capacity costs,
and closed-form proofs for every elimination.

The capacity cost of a gauge group $G$ is $C(G) = \dim(G) \cdot F(1)$
where $F$ is the unique cost function from $L_{\mathrm{Cauchy}}$
(Theorem~1.3 of \TSII{}).  Since $F(1) = 1$ (unit normalization),
$C(G) = \dim(G)$ --- an equality resting on the two named operational
premises (operational encoding for the lower bound; disjoint-anchor
realizability for the upper bound), composed in
Proposition~\ref{prop:generator_separation}.  This is the quantity
minimized in the classification below.

\begin{proposition}[Gauge generators are separated admissibility interfaces]\label{prop:generator_separation}
Under the operational-encoding and disjoint-anchor-realizability
premises (the operational cost premise pair of the front-matter
ledger), for any compact Lie group $G$ with Lie algebra $\mathfrak{g}$
of dimension $n$, the $n$ independent generators $\{T^a\}_{a=1}^n$
constitute $n$ \emph{separated} admissibility interfaces in the sense
of $L_{\mathrm{Cauchy}}$'s additivity premise (C1): for any pair
$a \ne b$, committing admissibility to $T^a$ does not reshape the
channel landscape seen by $T^b$, so the superadditive surplus
$\Delta(T^a, T^b) = 0$.  Therefore the total cost is additive:
$C(G) = \sum_{a=1}^n F(1) = \dim(G)$.
\end{proposition}

\begin{proof}
The exact count is the banked gauge clause: $n(G) = \dim(G)$, hence
$C(G) = \dim(G) \cdot \varepsilon$ in $F(1) = 1$ units
\coderef{check\_L\_cost\_gauge}{core.py}.  The lower bound
$C(G) \ge \dim(G)$ is Lemma~2.22 of \TSII{} (resolution rank:
invariance of domain on the $\dim(G)$-dimensional group manifold ---
basis-independent, commutativity-independent), and it consumes that
lemma's named operational-encoding premise: a continuous injective
operational encoding of the local group action is what invariance of
domain applies to.  The upper bound (no
superadditive surplus \emph{across generator channels beyond the
count}) is the C1-additivity clause: what must be checked is that the
$\dim(G)$ resolution channels are separated interfaces in
$L_{\mathrm{Cauchy}}$'s sense, i.e.\ that committing one does not
reshape the landscape of another ($\Delta = 0$).  By $T_M$ (interface
monogamy) and $L_{\mathrm{loc}}$, channels with disjoint anchor
supports are separated.  That the $\dim(G)$ resolution distinctions
can be anchored on disjoint supports is a named \emph{construction
premise} (disjoint-anchor realizability), not a consequence of
coordinate independence: the group coordinates are independent by
construction, but independence of coordinates does not by itself
deliver physical disjointness of anchor supports, so the realizability
of a disjointly anchored channel family is stated as the premise the
upper bound consumes.  Earlier versions argued the separation
from Killing-orthogonality of the generators; that argument is
demoted --- Killing orthogonality is a Lie-algebra inner-product
statement, it is not by itself orthogonality of the admissibility
projections under the anchor form $B$, and the Killing form is
degenerate on abelian factors.  The separation claim is carried by
the disjoint-anchor realizability premise (and, at the bank level, by
the two proof routes of \codeid{check\_L\_cost\_gauge}, which carry
the same premise structure: the check verifies the stated dimensions
and arithmetic, not the encoding or realizability premises); Killing
orthogonality is retained only as the natural bookkeeping basis.  Therefore the cost
is linear in $\dim(G)$ for every group in the classification class,
including $E_6$ at 78.
\end{proof}

\smallskip\noindent\textit{Remark.}\quad This proposition closes an otherwise
tacit step: without it, the cost formula $C(G) = \dim(G)$ would implicitly
assume that gauge generators are C1-admissible without derivation.
The capacity-minimality comparisons throughout this section
($E_6 \colon 78 \text{ vs.\ } \SU(3) \colon 8$, etc.) all rely on
Proposition~\ref{prop:generator_separation}; it is the bridge from
$L_{\mathrm{Cauchy}}$'s $\mathbb{N}$-domain additivity to the linearity
of cost in $\dim(G)$.

% ======================================================================
% Paper 2 Technical Supplement — §5
% Gauge Template Uniqueness: Exhaustive Lie Algebra Classification
% ======================================================================
%
% PURPOSE.  Paper 2 §5.1 states L_gauge,unique in a compact table
% covering the classification (families + exceptionals; 17 tested
% representatives).  The reviewer asks (Q10):
% "Can you provide a theorem excluding other groups with complex 3D
% irreps or alternative minimal faithful carriers (e.g., embeddings
% into SU(N>3)) on 'capacity minimality' grounds, with a quantified
% cost function?"  This section provides the full classification with
% explicit exclusion reasons, quantified capacity costs, and
% closed-form proofs for every elimination.
%
% UPSTREAM: Theorem R (§3), L_Cauchy (§1), Lie algebra classification
% (imported mathematics: Killing 1888, Cartan 1894).
% ======================================================================


\subsection{Setup: the classification problem}
\label{sec:gauge_setup}

Theorem~R (\S3.5 of \TSII{}) provides three abstract carrier
requirements:
\begin{enumerate}[label=\textbf{R\arabic*},nosep,leftmargin=3em]
  \item A faithful complex carrier of dimension $\ge 3$ with irreducible
        totally antisymmetric invariant of rank $k \ge 3$.
  \item A faithful pseudoreal carrier of dimension exactly~$2$.
  \item A single abelian factor ($\U(1)$) with distinct charge
        assignments for all physical multiplets.
\end{enumerate}
R1--R3 are carried in from the foundational chain (Theorem~R,
\S3.5 of \TSII{}), where their derivations and grades live; this
section consumes them as the classification's carrier requirements.
The problem is now purely mathematical: \emph{which compact Lie group(s)
can host all three carriers simultaneously, at minimum total capacity cost?}

The capacity cost of a gauge group $G$ is $\dim(G) \cdot \varepsilon^*$
($L_{\mathrm{Cauchy}}$, Theorem~1.3 of \TSII{}, under the
operational cost premise pair of
Proposition~\ref{prop:generator_separation}): the number of
independent gauge generators, each costing one admissibility channel.
This cost functional is unique (no free parameter) and is the quantity
to be minimized.


\subsection{Classification of the confining factor over the declared carriers (R1)}
\label{sec:R1_classification}

\noindent\textit{Input.}  The classification of compact simple Lie algebras
is a completed theorem of pure mathematics.  There are exactly four infinite
families ($A_n$, $B_n$, $C_n$, $D_n$ for $n \ge 1, 2, 3, 4$ respectively)
and five exceptional algebras ($G_2$, $F_4$, $E_6$, $E_7$, $E_8$).
For each, the fundamental (lowest-dimensional faithful) representation
has a known dimension and reality type (real, pseudoreal, or complex).
The scan below does \emph{not} stop at fundamentals: for the
orthogonal series the chiral spinor representations are faithful
carriers of the spin groups and must be tested in their own right ---
type $D_{2k+1}$ has \emph{complex} chiral spinors
($\mathrm{Spin}(10)$'s $\mathbf{16}$ is the physical case), type
$D_{2k}$ real or pseudoreal ones, and types $B_n$, $C_n$ contribute no
complex spinors.  Earlier versions tested only the vector-type
fundamentals of the $D$ series; the spinor rows are now in the table.
This data is imported from standard references
\cite{Humphreys1972,FultonHarris1991} and is not framework-dependent.

One boundary is stated once, honestly: no highest-weight
exhaustiveness theorem over all faithful complex irreducible carriers
is claimed.  The examined carrier set is the fundamentals of the four
families and the five exceptional algebras, plus the stated spinor
rows of the orthogonal series; candidates outside that set are
eliminated by the cost ranking (under the operational premise pair of
Proposition~\ref{prop:generator_separation}), not by a
representation-theoretic exclusion theorem.

\noindent\textit{Filters.}  R1 imposes three necessary conditions on the
confining factor $G_{\mathrm{conf}}$:
\begin{enumerate}[label=(F\arabic*),nosep]
  \item The fundamental representation must be \emph{complex}
        ($V \not\cong V^*$ as $G$-modules).  This is forced by $B1'$
        (Theorem~3.6 of \TSII{}): real and pseudoreal carriers
        cannot sustain oriented composites.
  \item The representation must admit a totally antisymmetric
        \emph{primitive} invariant of rank $k \ge 3$ --- primitive
        meaning an invariant that supports a $k$-body neutral composite
        of the carrier at small $k$, exemplified by the rank-$N$
        $\varepsilon$-tensor of the $\SU(N)$ fundamental.  Stated
        precisely, once and for all: (F2) requires a totally
        antisymmetric \emph{primitive} invariant in $V^{\otimes k}$ at
        rank $3 \le k \le \mathrm{rank}(G) + 1$, where primitive means
        not a power or product of lower-rank invariants.  For the
        $\SU(N)$ fundamental the $\varepsilon$-tensor sits at
        $k = N = \mathrm{rank}(G) + 1$ and is primitive; the top
        exterior power of a generic complex carrier sits at
        $k = \dim V \gg \mathrm{rank}(G) + 1$ and fails the rank
        bound.  Two
        precisions, both corrections to earlier versions: (i) the
        invariant must be totally \emph{antisymmetric} --- a symmetric
        invariant does not produce a fermionic composite of identical
        carriers ($E_6$'s cubic on the $\mathbf{27}$ is symmetric and
        does NOT satisfy this filter; the v3.9 table marked it as a
        rank-3 pass, which was wrong); (ii) a bare existence condition
        at unrestricted rank is vacuous for complex carriers, because
        the top exterior power $\Lambda^{\dim V} V$ of any
        representation of a connected semisimple group is invariant
        (trivial determinant character).  The filter's selective
        content is therefore antisymmetry at \emph{low} rank: for the
        $\SU(N)$ fundamental the $\varepsilon$-invariant sits at
        $k = N$, a physical few-body composite for the small $N$ the
        cost ranking keeps alive; candidates whose only antisymmetric
        invariant is a large-rank volume form gain nothing physical
        from it, and the cost ranking eliminates them independently.
        Bilinear ($k=2$) invariants are excluded because they produce
        self-conjugate composites ($B = \bar{B}$).  The rank window
        itself is a \emph{selected filter premise} --- named here
        \textbf{F2-window} and carried in the imported-ingredient
        ledger: requiring the primitive antisymmetric invariant to sit
        at $3 \le k \le \mathrm{rank}(G)+1$ is a declared physical
        filter (few-body neutral composites of the carrier must exist
        at low rank), motivated by the composite-realizability reading
        of R1, and it is a premise of the classification, not a
        consequence of Lie theory.
  \item The representation must be \emph{faithful}: every non-identity
        group element acts non-trivially.  This is forced by
        admissibility completeness at its grade (EC,
        [P\textsubscript{structural\_reading} $\mid$ R-EC-inv $+$
        I-typing]): if a gauge transformation acts trivially on all
        matter, it is physically undetectable and should not be
        counted.
\end{enumerate}

\begin{theorem}[$L_{\mathrm{gauge,unique}}$ --- Step 2: Confining Factor Classification]\label{thm:R1_class}
The only compact simple Lie algebra whose fundamental representation
satisfies (F1)--(F3) is $A_n = \mathfrak{su}(n+1)$ with $n \ge 2$
(equivalently, $\SU(N_c)$ with $N_c \ge 3$).
\end{theorem}

\begin{proof}
We test each family and each exceptional algebra against (F1)--(F3).
The full data is:

\smallskip
\begin{center}
\small
\renewcommand{\arraystretch}{1.15}
\begin{tabular}{@{}llccccl@{}}
\toprule
\textbf{Algebra} & \textbf{Group} & $\dim_{\mathrm{fund}}$ &
\textbf{Type} & \textbf{(F1)}  & \textbf{(F2)} & \textbf{Verdict} \\
\midrule
\multicolumn{7}{@{}l}{\textit{Family $A_n = \mathfrak{su}(n+1)$}} \\
$A_1$ & $\SU(2)$  &  2 & Pseudoreal & Fail & ---  & Excluded (F1) \\
$A_2$ & $\SU(3)$  &  3 & Complex    & Pass & $k{=}3$: Pass & \textbf{Pass} \\
$A_3$ & $\SU(4)$  &  4 & Complex    & Pass & $k{=}4$: Pass & Pass$^*$ \\
$A_4$ & $\SU(5)$  &  5 & Complex    & Pass & $k{=}5$: Pass & Pass$^*$ \\
$A_{n\ge5}$ & $\SU(n{+}1)$ & $n{+}1$ & Complex & Pass & $k{=}n{+}1$: Pass & Pass$^*$ \\[4pt]
\multicolumn{7}{@{}l}{\textit{Family $B_n = \mathfrak{so}(2n+1)$, $n \ge 2$}} \\
$B_2$ & $\mathrm{SO}(5)$  &  5 & Real & Fail & --- & Excluded (F1) \\
$B_3$ & $\mathrm{SO}(7)$  &  7 & Real & Fail & --- & Excluded (F1) \\
$B_{n\ge2}$ & $\mathrm{SO}(2n{+}1)$ & $2n{+}1$ & Real & Fail & --- & Excluded (F1) \\[4pt]
\multicolumn{7}{@{}l}{\textit{Family $C_n = \mathfrak{sp}(2n)$, $n \ge 3$}} \\
$C_3$ & $\mathrm{Sp}(6)$  &  6 & Pseudoreal & Fail & --- & Excluded (F1) \\
$C_{n\ge3}$ & $\mathrm{Sp}(2n)$ & $2n$ & Pseudoreal & Fail & --- & Excluded (F1) \\[4pt]
\multicolumn{7}{@{}l}{\textit{Family $D_n = \mathfrak{so}(2n)$, $n \ge 4$ (vector and spinor carriers)}} \\
$D_4$ & $\mathrm{SO}(8)$  &  8 & Real & Fail & --- & Excluded (F1) \\
$D_{n\ge4}$ & $\mathrm{SO}(2n)$ & $2n$ & Real & Fail & --- & Excluded (F1) \\
$D_4$ spinor & $\mathrm{Spin}(8)$ & 8 & Real & Fail & --- & Excluded (F1) \\
$D_{2k}$ spinor & $\mathrm{Spin}(4k)$ & $2^{2k-1}$ & Real/psr. & Fail & --- & Excluded (F1) \\
$D_5$ spinor & $\mathrm{Spin}(10)$ & 16 & Complex & Pass & no cubic inv. & Excluded (F2)$^\ddagger$ \\
$D_{2k+1}$ spinor & $\mathrm{Spin}(4k{+}2)$ & $2^{2k}$ & Complex & Pass & $^\ddagger$ & Excluded (F2)$^\ddagger$ \\[4pt]
\multicolumn{7}{@{}l}{\textit{Exceptional algebras}} \\
$G_2$ &          &  7 & Real       & Fail & --- & Excluded (F1) \\
$F_4$ &          & 26 & Real       & Fail & --- & Excluded (F1) \\
$E_6$ &          & 27 & Complex    & Pass & $k{=}3$: sym., Fail & Excluded$^\dagger$ \\
$E_7$ &          & 56 & Pseudoreal & Fail & --- & Excluded (F1) \\
$E_8$ &          &248 & Real       & Fail & --- & Excluded (F1) \\
\bottomrule
\end{tabular}
\end{center}

\smallskip\noindent ${}^*$Pass$^*$: satisfies (F1)--(F3) but is
non-minimal (higher $\dim(G)$ than $\SU(3)$; eliminated in Step~3
below).

\smallskip\noindent ${}^\dagger$$E_6$ satisfies (F1), but its
celebrated rank-3 invariant on the $\mathbf{27}$ is the \emph{symmetric}
cubic form, not a totally antisymmetric 3-tensor --- so it does not
satisfy (F2) as now correctly stated (the v3.9 table's ``$k{=}3$:
Pass'' conflated the two symmetry types).  Independently of (F2), $E_6$
is excluded by capacity minimality: $\dim(E_6) = 78$ vs.\
$\dim(\SU(3)) = 8$; the cost cut is retained as the exclusion of
record so that nothing rests on the invariant-type correction alone.

\smallskip\noindent ${}^\ddagger$The $\mathrm{Spin}(10)$ spinor
$\mathbf{16}$ is a faithful complex carrier and passes (F1); it fails
(F2) at low rank because it admits \emph{no} cubic invariant at all:
$\mathbf{16} \otimes \mathbf{16} = \mathbf{10} \oplus \mathbf{120}
\oplus \mathbf{126}$ contains no $\overline{\mathbf{16}}$, so no
invariant exists in $\mathbf{16}^{\otimes 3}$ of any symmetry type.
Higher-rank antisymmetric singlets, where they exist, are large-body
volume-form-type composites with no small-$k$ screening role, and the
cost ranking ($\dim \mathrm{Spin}(10) = 45$ vs.\ $8$) eliminates the
candidate independently.  The same two-step exclusion (no low-rank
antisymmetric invariant; cost) covers the general $D_{2k+1}$ spinor.

\medskip\noindent\textit{Mechanism of exclusion.}\quad
The $B_n$, $C_n$ families, the $D_n$ vector carriers, the $D_{2k}$
spinors, and the exceptionals $G_2$, $F_4$, $E_7$, $E_8$ are all
excluded by (F1) alone: those carriers are real or pseudoreal, and
$B1'$ (Theorem~3.6 of \TSII{}) proves that real/pseudoreal carriers
cannot sustain oriented composites.  The complex carriers among
compact simple Lie algebras are the $A_n$ ($n \ge 2$) fundamentals
(and their conjugates --- the conjugate fundamental is a second
faithful $n{+}1$-dimensional irrep, equivalent for this scan under the
CPT/conjugation quotient), the $E_6$ $\mathbf{27}$, and the
$D_{2k+1}$ chiral spinors; the latter two fall to (F2) and cost as
recorded in the table notes.

$E_6$ passes (F1), fails (F2) (its rank-3 invariant on the
$\mathbf{27}$ is symmetric, per the table), and is independently
excluded on capacity grounds:
$\dim(E_6) = 78$ requires 78 gauge generators, each costing one
admissibility channel.  The $\SU(3)$ alternative requires 8 generators.
By $L_{\mathrm{Cauchy}}$ (Theorem~1.3 of \TSII{}), the cost
functional is linear in generator count, so $E_6$ costs
$78/8 = 9.75$ times as much as $\SU(3)$.  Minimality selects the
cheapest adequate structure --- and minimality is A1 $+$ MD $+$ A2, not
A1 alone: A1 bounds the budget, MD (via $L_{\varepsilon^*}$) prices the
floor, A2 selects the argmin (Lemma~2.24 of \TSII{}).

The only surviving family is $A_n = \SU(N_c)$ with $N_c \ge 3$.
\end{proof}

\smallskip\noindent\textit{Where the $E_6$ cut actually happens.}\quad
$E_6$ satisfies (F1) --- it has a complex faithful carrier --- but not
(F2): its rank-3 invariant on the $\mathbf{27}$ is symmetric, not
antisymmetric.  The point of this paragraph survives that verdict:
$E_6$'s exclusion does not \emph{need} (F2).  Independently of any
representation-theoretic filter, the difference between ``$\SU(N_c)$
survives'' and
``$E_6$ is rejected'' is a
capacity-minimality result under the minimality package (A1 $+$ MD $+$
A2).  A reader who doubts the
invariant-type verdict on the cubic --- or who accepts the carrier
requirements but doubts the capacity-minimality clause --- still
faces the cost cut: the
minimality clause closes the gap on its own, consumed at its grade
(CM, Lemma~2.24 of \TSII{}) through the admissibility-cost ledger.  The $E_6$
exclusion is therefore a direct consequence of the minimum-cost
ontology: $E_6$ is admissible (its cost fits below the capacity bound)
but is not minimum-cost; $\SU(3)$ is.  ``$E_6$ is rejected'' is
shorthand for ``$E_6$ has higher cost than $\SU(3)$ and both are
admissible,'' not for an act of rejection.


\subsection{Capacity optimization: $N_c = 3$}
\label{sec:Nc_optimization}

Within the surviving $\SU(N_c)$ family, all $N_c\ge3$ pass the purely
local color filters.  The remaining question is not whether larger color
groups are mathematically possible; they are.  The APF question is which
member of the admissible family minimizes the finite admissibility ledger
under the same representation and anomaly constraints.

\begin{definition}[Color optimization problem]
\label{def:color_optimization}
The color optimization problem is the quadruple
$(\mathcal{N},J,\mathcal C,\preceq)$ where:
\begin{enumerate}[nosep,label=(\roman*),leftmargin=1.8em]
  \item $\mathcal{N}=\{N_c\in\mathbb N: N_c\ge3\}$ is the family surviving
  oriented-composite robustness ($B1'$).
  \item $\mathcal C(N_c)$ is the set of chiral templates with one
  color-doublet family, one lepton-doublet family, anomaly cancellation,
  asymptotic freedom, and the F1--F7 scan constraints generalized from
  $3$ to $N_c$ colors.
  \item The objective is the channel ledger
  \[
    J(N_c)=(4N_c+3)N_{\mathrm{gen}}+4+\bigl((N_c^2-1)+3+1\bigr),
    \qquad N_{\mathrm{gen}}=3,
  \]
  where the three terms are fermions, Higgs, and gauge generators.
  \item $N_c\preceq N_c'$ iff $J(N_c)\le J(N_c')$ after both candidates
  pass the same template filters.
\end{enumerate}
The optimization is invariant under positive affine rescalings of the
ledger because $L_{\mathrm{Cauchy}}$ fixes all channel costs to the same
integer unit; no continuous weight can reverse the ordering.
The scope of $J$ is clause~(ii)'s template class: $J(N_c)$ proves
family-monotonicity within SM-family-shaped templates (one
color-doublet family, one lepton-doublet family, filters generalized
to $N_c$ colors); it is not a joint optimization over template shape.
The joint statement is the two-phase reading below
(Proposition~\ref{prop:witten} and the consistency note following it).
\end{definition}

\begin{remark}[What the color cost does and does not count]
\label{rem:color_cost_functional_v35}
The objective $J$ is a structural channel ledger, not an on-shell degree-of-freedom count and not a quadratic-Casimir energy functional.  It counts three positive-cost resources: matter channels, Higgs channels, and gauge-generator routing channels.  Casimir weights, beta-function coefficients, and propagating polarizations may refine dynamics after the template is fixed; they are not subtracted from the structural existence cost.  A cost functional that subtracts gauge constraints, treats new generators as free, or cancels matter channels against ghosts is therefore not a small deformation of $J$ but a different ontology: it counts physical excitations after gauge fixing rather than admissibility infrastructure before gauge fixing.  Proposition~\ref{prop:nc3_robustness_v34} states the allowed robustness cone for genuine positive-cost deformations.
\end{remark}

Two further filters apply:

\begin{proposition}[Witten anomaly excludes even $N_c$]\label{prop:witten}
For gauge group $\SU(N_c) \times \SU(2) \times \U(1)$, the number of
$\SU(2)$ doublets per generation is $N_c + 1$ (one quark doublet
contributing $N_c$ colors, plus one lepton doublet).  The Witten anomaly
\cite{Witten1982} requires this number to be even.  Therefore $N_c$ must
be odd.
\end{proposition}

\noindent\textit{Joint consistency with $T_{\mathrm{field}}$.}\quad
The statement ``one quark doublet plus one lepton doublet per generation''
is the SM-like template structure that emerges from $T_{\mathrm{field}}$
(\S\ref{sec:fermion_scan}) rather than an input independent of it.  The
Witten-anomaly cut applied in the current section therefore functions as
a \emph{joint consistency} check with the fermion scan: for any odd
$N_c \ge 3$, the minimum-DOF anomaly-free template has the (quark-doublet,
lepton-doublet) structure, and the Witten filter is satisfied; for any
even $N_c$, no admissible template of this structure survives.  The
logical ordering is not ``N$_c$ is fixed in \S5 then templates are scanned
in \S6.''  Rather, the admissible N$_c$'s and admissible fermion templates
are a joint optimum, and the supplement exhibits that joint optimum by
first narrowing N$_c$ to odd values under the generic template ansatz
(this section) and then closing N$_c = 3$ by the full scan (\S\ref{sec:fermion_scan}).
Readers who prefer a single-step derivation may take Propositions~\ref{prop:witten}
and \ref{prop:Nc3} together with Theorem~\ref{thm:Tfield} as a simultaneous
result: $(N_c = 3, \text{SM template})$ is the unique joint
minimum-cost admissible pair.

\begin{proposition}[Capacity minimality selects $N_c = 3$]\label{prop:Nc3}
Under Definition~\ref{def:color_optimization}, within the SM-family
template class of its clause~(ii), $N_c=3$ is the unique minimum-cost
admissible color rank.
\end{proposition}

\begin{proof}
The Witten anomaly removes even $N_c$.  For odd $N_c\ge3$ the objective
is
\[
  J(N_c)=(4N_c+3)\cdot3+4+(N_c^2+3)=N_c^2+12N_c+16.
\]
For $N_c' > N_c \ge 3$,
\[
  J(N_c')-J(N_c)=(N_c'-N_c)(N_c'+N_c+12)>0.
\]
Therefore $J$ is strictly increasing on the admissible odd ranks, and
$N_c=3$ is the unique minimum.  The first alternatives make the ordering
visible:
\begin{center}
\scriptsize
\renewcommand{\arraystretch}{1.12}
\begin{tabular}{@{}p{1.5cm}p{3.2cm}p{2.5cm}p{1.5cm}p{3.4cm}@{}}
\toprule
\textbf{Candidate} & \textbf{Local color filter} & \textbf{Witten filter} & \textbf{$J(N_c)$} & \textbf{Verdict} \\
\midrule
$\SU(2)$ & Fails $B1'$ / baryon-orientation rank & -- & -- & Not in $\mathcal N$ \\
$\SU(3)$ & Pass & Pass ($4$ doublets) & 61 & \textbf{Unique minimum} \\
$\SU(4)$ & Pass local color & Fails ($5$ doublets) & -- & Global anomaly \\
$\SU(5)$ & Pass & Pass ($6$ doublets) & 101 & Admissible but non-minimal \\
$\SU(7)$ & Pass & Pass ($8$ doublets) & 149 & Admissible but non-minimal \\
$E_6$ & Simple alternative & No product chirality ledger & $>61$ & Non-minimal / wrong template \\
\bottomrule
\end{tabular}
\end{center}
The result is therefore not a hidden phenomenological choice of three
colors.  It is the first admissible rank after the oriented-composite
lower bound and the global anomaly cut, and the ledger is strictly
monotone thereafter.
\end{proof}




\begin{proposition}[Robustness cone for the $N_c=3$ argmin]
\label{prop:nc3_robustness_v34}
The $N_c=3$ result is stable under every positive cost deformation that keeps
fermion channels and gauge generators as independently costly admissibility
resources.  More precisely, replace Definition~\ref{def:color_optimization}'s
unit ledger by
\begin{equation}
  J_{\bf w}(N_c)=w_g\bigl(N_c^2+3\bigr)
  +w_f\bigl(3(4N_c+3)\bigr)+4w_H+\rho(N_c),
\end{equation}
with $w_g,w_f,w_H>0$.  If the perturbation satisfies
\begin{equation}
  |\rho(N_c')-\rho(N_c)|<
  (N_c'-N_c)\bigl[w_g(N_c'+N_c)+12w_f\bigr]
\end{equation}
for all odd $N_c'>N_c\ge 3$, then $N_c=3$ remains the unique minimum
among admissible odd ranks.
\end{proposition}

\begin{proof}
For any odd $N_c'>N_c\ge3$,
\begin{equation}
  J_{\bf w}(N_c')-J_{\bf w}(N_c)
  =(N_c'-N_c)\bigl[w_g(N_c'+N_c)+12w_f\bigr]
  +\rho(N_c')-\rho(N_c).
\end{equation}
The stated bound makes this difference strictly positive.  Even ranks are
removed by Proposition~\ref{prop:witten}; $N_c=2$ is removed by the
oriented-composite carrier filter; exceptional/simple alternatives add
independent generators or fail the same low-energy product-template filters.
Thus the first admissible odd rank is the unique optimum throughout the
positive robustness cone.  The unit-cost ledger used in
Definition~\ref{def:color_optimization} is the special case
$w_g=w_f=w_H=1$ and $\rho=0$.
\end{proof}

\begin{remark}[Counterfactuals that would flip the color optimum]
The theorem also identifies what would have to fail.  A larger color rank
could win only if one allowed negative or zero marginal cost for new gauge
generators, treated additional colored fermion channels as cost-free, or
changed the admissibility filters so that the $N_c=3$ SM template was no
longer admissible while a higher-rank template remained admissible.  Those are
not small perturbations of the APF ledger; they are changes to the resource
semantics or to the physical constraints R1--R3.
\end{remark}

\subsection{The uniqueness core: abelian factor, center quotient, color rank}
\label{sec:uniqueness_core_v33}

The preceding propositions establish the four pieces that reviewers most
often conflate: algebra-level inner automorphisms, the linear cover acting
on matter, the single relative abelian direction, and the color-rank optimum.
For reference:
\begin{enumerate}[label=\textbf{U\arabic*},nosep]
  \item Algebraically, $\mathrm{Inn}(\mathcal A)=\prod_k\mathrm{PU}(n_k)$.
  \item On charged matter modules, the relevant projective factors are
  represented by their simply connected covers, giving the familiar
  $\SU(n)$ action where fundamental matter exists.
  \item The center-action table in Proposition~2.26 of \TSII{} shows
  that the common kernel on the derived multiplets is exactly $\mathbb Z_6$.
  \item The color-rank optimum is the joint optimum over rank and template,
  summarized in Table~\ref{tab:joint_Nc_template_optimum}.
\end{enumerate}

\begin{center}
\scriptsize
\renewcommand{\arraystretch}{1.1}
\begin{tabular}{@{}p{2.6cm}p{3.0cm}p{3.0cm}p{4.2cm}@{}}
\toprule
\textbf{Rank} & \textbf{Best local template at that rank} & \textbf{Minimum ledger} & \textbf{Reason not selected / selected} \\
\midrule
$N_c=2$ & two-color pseudo-real template & below SM-like count but fails baryon-orientation and chiral color separation & Excluded before optimization: no oriented three-color baryon carrier and no faithful QCD-like confinement ledger. \\
$N_c=3$ & SM template: $Q,u^c,d^c,L,e^c$ & $J=61$ & Unique joint survivor: Witten-even doublet count, anomaly closure, minimum active channels. \\
$N_c=4$ & nearest SM-like extension & -- & Witten anomaly: $N_c+1=5$ weak doublets. \\
$N_c=5$ & SM-like extension with five colors & $J=101$ & Anomaly-compatible but strictly higher active-channel cost than $N_c=3$. \\
$N_c\ge7$ odd & SM-like extension & $J=N_c^2+12N_c+16>101$ & Strictly monotone cost above the $N_c=3$ optimum. \\
Exceptional/simple unification & $E_6$, $\SU(5)$-like simple carriers & $>61$ or wrong low-energy factorization & Adds unneeded generators/exotics or imports high-energy unification structure not forced by finite-interface admissibility. \\
\bottomrule
\end{tabular}
\captionof{table}{Per-$N_c$ joint optimum audit.  The table distinguishes the earlier marginal SM-like comparison from the stronger joint statement: for each allowed color-rank family, the best admissible template either fails a structural filter or has larger active-channel ledger than $(N_c=3,\mathrm{SM})$.}
\label{tab:joint_Nc_template_optimum}
\end{center}

\begin{lemma}[Outer-permutation locality criterion]
\label{lem:outer_kinematic}
Let $\mathcal A=\bigoplus_k M_{n_k}(\mathbb C)$ be the finite admissibility
algebra.  A block permutation is a continuous gauge symmetry only if it
preserves every interface's local-support assignment and lies in the identity
component of the codespace-preserving automorphism group.  Outer permutations
between distinct blocks fail this criterion: they move matter modules between
support blocks, or, when the blocks are cost-degenerate, form a discrete
permutation component separated from $\mathrm{Inn}(\mathcal A)$.
Consequently they are kinematic/flavor relabelings, not Lie-algebra gauge
generators.
\end{lemma}

\begin{proof}
Inner automorphisms act inside a fixed simple summand and preserve the local
support projector of each interface.  A nontrivial outer permutation changes
the summand label.  If the summands are not identical as cost-equipped local
codespaces, support preservation fails immediately.  If they are identical,
the permutation is an element of a finite symmetric group of degenerate
labels; it has no infinitesimal generator and no continuous path to the
identity inside $\mathrm{Aut}(\mathcal A)$.  Gauge bosons arise from the
Lie algebra of continuous local symmetries, so such discrete permutations
cannot be gauge directions.
\end{proof}

\smallskip\noindent\textit{Edge case: family-non-universal $\U(1)$ directions.}\quad
Anomaly-free combinations such as $L_\mu-L_\tau$ require generation labels
to be independently gauge-addressable distinctions.  The minimal template
used in this supplement treats generations as repeated copies of the same
one-generation module; it does not pay an additional admissibility channel for
``muon family'' versus ``tau family'' as a primitive gauge label.  Therefore
family-non-universal abelian directions are excluded within the minimal-template
theorem by capacity minimality, not by anomaly algebra alone.  A theory that
postulates generation-resolving structure as input lies outside the minimal
APF theorem rather than contradicting it.

\subsection{Chiral factor classification (R2)}
\label{sec:R2_classification}

\begin{proposition}[$\SU(2)$ is the unique chiral factor]\label{prop:R2_unique}
R2 requires a compact simple Lie group with a faithful pseudoreal
representation of dimension exactly~$2$.  Among all compact simple Lie
algebras:
\begin{enumerate}[label=(\roman*),nosep]
  \item $\mathfrak{su}(2)$ has a 2-dimensional irreducible representation
        (the fundamental doublet) which is pseudoreal
        ($J = i\sigma_2 K$, $J^2 = -1$).
  \item All other compact simple Lie algebras have minimum faithful
        dimension $\ge 3$: $\mathfrak{su}(n+1)$ for $n \ge 2$ has
        $\dim_{\mathrm{fund}} = n + 1 \ge 3$; $\mathfrak{so}(2n+1)$
        has $\dim_{\mathrm{fund}} \ge 5$; $\mathfrak{sp}(2n)$ for $n \ge 2$
        has $\dim_{\mathrm{fund}} \ge 4$; all exceptionals have
        $\dim_{\mathrm{fund}} \ge 7$.
\end{enumerate}
Therefore $\SU(2)$ is the unique compact simple Lie group satisfying R2.

At the group level, the gauge factor must be $\SU(2)$ rather than its
quotient $\mathrm{SO}(3) = \SU(2)/\mathbb{Z}_2$: the doublet
representation does not descend to $\mathrm{SO}(3)$ (it is not
single-valued under $2\pi$ rotations), so faithfulness requires
the simply connected cover $\SU(2)$.
\end{proposition}


\subsection{Abelian factor and simple-group alternatives (R3)}
\label{sec:R3_simple}

\begin{proposition}[$\U(1)$ is the abelian carrier]\label{prop:R3_unique}
R3 requires at least one abelian grading direction, and
$\U(1) \cong \mathbb{R}/\mathbb{Z}$ is the unique connected compact
1-dimensional abelian Lie group, so each abelian grading direction is
a $\U(1)$ (Theorem~R, R3 derivation in \S3 of \TSII{}).  That
\emph{exactly one} such factor closes the template --- that a single
$\U(1)$ with distinct hypercharge assignments resolves all multiplet
degeneracies of the scanned content, and that any additional factor
adds a gauge generator with no admissibility gain --- is the
a-posteriori closure of \S\ref{sec:one_u1_closure}, consumed after
the scan (Theorem~\ref{thm:gauge_unique}, clause~(ii)).
\end{proposition}

\begin{proposition}[Simple-group alternatives are excluded]\label{prop:no_simple}
Any simple compact Lie group $G_{\mathrm{simple}}$ containing
$\SU(3) \times \SU(2) \times \U(1)$ as a subgroup has
$\dim(G_{\mathrm{simple}}) \ge 24$:

\smallskip
\begin{center}
\small
\begin{tabular}{@{}lccc@{}}
\toprule
\textbf{Group} & $\dim$ & \textbf{Cost ratio} & \textbf{Contains SM?} \\
\midrule
$\SU(3) \times \SU(2) \times \U(1)$ & 12 & $1.0\times$ & (Product) \\
$\SU(5)$   & 24 & $2.0\times$ & Yes (Georgi--Glashow) \\
$\mathrm{SO}(10)$ & 45 & $3.75\times$ & Yes \\
$E_6$ & 78 & $6.5\times$ & Yes \\
\bottomrule
\end{tabular}
\end{center}

\smallskip\noindent Every simple envelope costs at least $2\times$ the
product structure.  Minimality (A1 $+$ MD $+$ A2,
Lemma~2.24 of \TSII{}) selects the product.
The earlier versions added a second exclusion clause: that a simple
group has a single coupling constant and therefore cannot yield
independently running confining and chiral factors.  That clause is
withdrawn --- it is not a valid exclusion, because symmetry breaking
and threshold matching give the unbroken subgroup factors distinct
infrared runnings from a single ultraviolet coupling (the standard
grand-unification pattern).  What survives, and suffices, is the cost
comparison above: every simple envelope prices at least $2\times$ the
product structure at the same enforcement content, and minimality
(A1 $+$ MD $+$ A2) selects the product.  The requirement that the
chiral factor not confine at accessible scales is dynamical input, not
a group-theoretic exclusion, and is not used here.
\end{proposition}


\subsection{Consolidation: unique gauge template}
\label{sec:gauge_consolidation}

Combining Steps~1--6:

\begin{theorem}[$L_{\mathrm{gauge,unique}}$: Gauge Template Uniqueness]\label{thm:gauge_unique}
Two clauses, at different logical stations.

\smallskip\noindent\textit{(i) Nonabelian carrier ranking (forward).}\quad
The unique nonabelian gauge structure satisfying the carrier
requirements R1--R2 at minimum capacity cost is
\[
  G_{\mathrm{na}} \;=\; \SU(3) \times \SU(2),
\]
and R3 requires at least one abelian grading direction alongside it.
No alternative compact nonabelian structure satisfies both carrier
requirements.  No simple-group envelope is capacity-competitive.

\smallskip\noindent\textit{(ii) Abelian closure (a-posteriori pointer).}\quad
That exactly one $\U(1)$ closes the template --- necessity,
sufficiency, directional uniqueness, and the inadmissibility of any
second factor --- is not part of the forward statement: it is proved
a-posteriori for the selected matter content by the closure theorems
of \S\ref{sec:one_u1_closure}, consumed after the scan.  The composite
of (i) with that closure is the full gauge template
\[
  G \;=\; \SU(3) \times \SU(2) \times \U(1),
  \qquad \dim(G) = 12.
\]

\smallskip\noindent\textit{Derivation chain:}\quad
$\text{A1} \to \text{Theorem R (R1+R2+R3)} \to
\text{Lie classification} \to
\SU(N_c) \times \SU(2)\ (+\text{ one abelian grading, R3}) \to
\text{Witten} + \text{minimality} \to N_c = 3 \to
\text{one-U(1) closure (\S\ref{sec:one_u1_closure})}$.

\smallskip\noindent\textit{Imported mathematics:}\quad
The classification of compact simple Lie algebras (Killing 1888, Cartan 1894)
and the Witten anomaly theorem (Witten 1982) are used as external
mathematical inputs.  Both are established results with independent
proofs; neither is framework-specific.  All physical inputs (R1--R3)
are Theorem~R's carrier requirements, proved at their grades in
\S3 of \TSII{}.
\end{theorem}
\coderef{check\_L\_gauge\_template\_uniqueness}{gauge.py}




\newpage


% ══════════════════════════════════════════════════════════════
\section{Fermion content: complete scan of the declared space}
\label{sec:fermion_scan}
% ══════════════════════════════════════════════════════════════

\begin{tcolorbox}[apfbox, title={Reproducibility contract for the 1{,}680-template scan}]
\footnotesize
The scan is fully specified by four inputs: (i) the derived gauge template
$\SU(3)\times\SU(2)\times\U(1)$; (ii) the bounded representation list
proved in \S\ref{sec:exclusion_proofs}; (iii) the generation count and
multiplicity bounds stated in the search-space construction; and (iv) the
seven exact filters F1--F7 below.  The output is the set of templates
surviving all seven filters, quotiented by CPT.  No floating-point
thresholds, fitted coefficients, or data-dependent choices enter.  A
reproduction must report the initial 1{,}680 candidates, the waterfall
counts in \S\ref{sec:waterfall} under the canonical F6 predicate
(full-system non-degeneracy: 10 survivors pre-CPT, 5 post-CPT), the
robustness counts under the two variant predicates (spectator
reduction 8/4; uniform dimension 4/2; the predicate statement in
\S\ref{sec:seven_filters} states all three and designates the
canonical one), and the unique minimum-capacity 45-Weyl-fermion
survivor.
The standalone script and repository link in \cite{BrookeFermionScan} are
therefore not supplemental evidence for a different claim; they are the
machine-readable form of the finite enumeration specified here.

\smallskip\noindent\textbf{Audit files required for independent reproduction.}\quad
A complete release consists of: \texttt{scan\_inputs.json} (representation
bounds and equivalence relation), \texttt{scan\_outputs.csv} (all 1{,}680
canonical templates), \texttt{filter\_waterfall.csv} (F1--F7 survivor counts),
\texttt{near\_misses.csv} (templates failing only the last nontrivial filter),
and \texttt{certificate.sha256} (hash of the ordered bundle).  If any of these
files is absent from a public release, the paper's scan claim should be read as
a mathematically specified enumeration rather than as independently reproduced
software evidence.  The theorem claim itself is the finite enumeration proved
below; the repository is the machine-checkable certificate layer.
\end{tcolorbox}

\begin{proposition}[Boundedness of the declared enumeration]
\label{prop:scan_completeness_v35}
The 1{,}680-template scan is complete relative to the declared
representation list and multiplicity caps of
\S\ref{sec:exclusion_proofs}: P1--P2 exclude the enumerated
representation enlargements unconditionally, P3--P5 certify dominance
against the named extension directions, and the 1{,}680 candidates
exhaust the declared space.  The caps themselves are declared scan
inputs; that no admissible sub-45-DOF template exists outside them is
an open proof obligation, not a claim of this proposition.
\end{proposition}

\begin{proof}[Proof sketch]
P1--P2 reduce the representation content unconditionally: a single
field in an $\SU(3)$ representation of dimension $\ge 10$, or a
colored field in an $\SU(2)$ representation of dimension $\ge 3$,
exceeds the AF budget regardless of the rest of the template.  P3 and
P5 are dominance certificates against the named enlargement
directions (colorless triplets; field types beyond the anomaly-closed
core), and P4 is the two-colored-doublet class-dominance result
(class minimum 54, Proposition~\ref{prop:P4}).  Within the declared
caps, the 1{,}680 candidates are the exact finite quotient of the
declared space (the closed-form count in
\S\ref{sec:exclusion_proofs}).  What this establishes is boundedness
and dominance, not universal completeness: the caps are declared
inputs, and their completeness is the named open obligation of the
reproducibility contract.  The public scan artifact is the
machine-readable certificate of this bounded enumeration.
\end{proof}

\smallskip\noindent\textit{Descriptive reading of the scan.}\quad
The enumeration in this section is \emph{not} a physical process by
which nature ``chooses'' the Standard Model fermion content from a
space of candidates.  The admissibility conditions (five exclusion
proofs bounding the search space, plus seven sequential filters
encoding anomaly cancellation, asymptotic freedom, and
doublet--singlet content)
define a predicate
$\mathrm{Adm}(\mathcal{T})$ on chiral fermion templates.  The scan
is an \emph{algorithmic verification} that the set
$\{\mathcal{T} : \mathrm{Adm}(\mathcal{T})\}$ has a unique member up
to the stated equivalences --- namely the 45-Weyl-fermion SM
template.  The 1{,}680 distinct templates (4{,}680 ordered
iterations) are artefacts of the verification procedure, not
candidates in a dynamical competition.  The claim of this section is
therefore descriptive: ``the admissibility predicate is satisfied by
exactly one template,'' not ``the framework selects one template
from among 1{,}680.''

This section specifies the 1{,}680-template search space, defines the
seven sequential filters, provides the survivor count after each
filter, and states the boundedness result at its licensed form:
P1--P2 exclude the enumerated representation enlargements
unconditionally, P3--P5 certify dominance against the named extension
directions, and the caps themselves are declared scan inputs whose
completeness is an open proof obligation.

% ── §6: Fermion content: complete scan of the declared space ──

\subsection{Template definition}
\label{sec:template_def}

A \emph{chiral fermion template} is a finite multiset
$\mathcal{T} = \{(r_i^{(3)}, r_i^{(2)})\}_{i=1}^{|\mathcal{T}|}$
where $r_i^{(3)}$ is an irreducible representation of $\SU(3)$ and
$r_i^{(2)}$ is an irreducible representation of $\SU(2)$.  Each entry
represents a left-handed Weyl fermion field.  The total number of Weyl
degrees of freedom per generation is
$\mathrm{DOF} = \sum_i \dim(r_i^{(3)}) \cdot \dim(r_i^{(2)})$.
The template describes one generation; the full content is
$N_{\mathrm{gen}} \times \mathrm{DOF}$ with $N_{\mathrm{gen}} = 3$
(derived in Paper~2, \S9).

The search space is bounded by five closed-form exclusion proofs that
restrict both the allowed representations and the number of field types.


\subsection{Search space bounds: five exclusion proofs}
\label{sec:exclusion_proofs}

\begin{proposition}[P1: Large $\SU(3)$ representations are AF-excluded]
\label{prop:P1}
Any $\SU(3)$ representation with dimension $\ge 10$ (i.e.,
$\mathbf{10}$, $\mathbf{15}$, $\ldots$) exceeds the $\SU(3)$
asymptotic freedom budget as a single field.  For the $\mathbf{10}$:
the Dynkin index is $T(\mathbf{10}) = 15/2$; under the all-left-handed
Weyl convention of the scan (coefficient $2/3$ per Weyl field --- the
$4/3$ form is the Dirac-pair coefficient of \S\ref{sec:beta_schematic}
and does not apply to Weyl entries), the AF contribution of a single
field in the $\mathbf{10}$ with $\SU(2)$ singlet and 3 generations is
$(2/3)(15/2)(1)(3) = 15 > 11$.
Therefore no template containing any $\SU(3)$ rep of dimension
$\ge 10$ can be asymptotically free.  The allowed $\SU(3)$
representations are: $\mathbf{1}, \mathbf{3}, \bar{\mathbf{3}},
\mathbf{6}, \bar{\mathbf{6}}, \mathbf{8}$.
\end{proposition}

\begin{proposition}[P2: Colored $\SU(2)$ triplets and higher are AF-excluded]
\label{prop:P2}
Any colored field ($\dim(r^{(3)}) \ge 3$) in an $\SU(2)$
representation of dimension $\ge 3$ exceeds the $\SU(2)$ AF budget.
For a $(\mathbf{3}, \mathbf{3})$ field: $T(\mathbf{3}_{\SU(2)}) = 2$,
$\dim(\mathbf{3}_{\SU(3)}) = 3$, Weyl-convention contribution
$= (2/3)(2)(3)(3) = 12 > 22/3 \approx 7.33$.  Therefore colored fields are
restricted to $\SU(2)$ representations $\{\mathbf{1}, \mathbf{2}\}$.
\end{proposition}

\noindent\emph{Remark (what P2 does and does not exclude).}  P2's inequality
carries the factor $N_{\mathrm{gen}} = 3$: it excludes colored weak-triplets
\emph{at the forced multiplicity}.  A single unreplicated copy of such a field
is AF-clean ($ (2/3)(2)(3)(1) = 4 < 22/3$), and contents built on it can be
chiral and anomaly-free.  Their exclusion is the multiplicity theorems of
\S\ref{sec:three_generation_scope} (the capacity--beta identities and the
rival kill), not this proposition alone.  Example~4 of the worked-example
section exhibits the sharpest such rival and its kill map.

\begin{proposition}[P3: Colorless $\SU(2)$ triplets cannot lower the minimum]
\label{prop:P3}
This is a \emph{dominance} bound, not an a-priori search-space
theorem, and its logical form matters.  Phase one: the scan, run on
the declared template space (which contains no colorless triplets),
exhibits an admissible survivor at 45 DOF.  Phase two: any template
extending or modifying a survivor by a colorless $\SU(2)$ triplet
carries that triplet's own $3 \times 3 = 9$ DOF (or, replacing a
doublet, a net $+3$), so no colorless-triplet template beats the
exhibited 45 without removing content that the filters
(doublet--singlet content, anomaly closure) require.  The proposition therefore certifies that
\emph{enlarging} the declared space by colorless triplets cannot lower
the minimum; it is consumed after the scan's phase-one result, never
before it.  (The v3.9 wording cited the SM minimum inside a
search-space bound --- circular as stated; the two-phase form above is
the intended and valid argument.)
\end{proposition}

\begin{proposition}[P4: Multiple colored doublets are DOF-dominated]
\label{prop:P4}
Every template with two colored $\SU(2)$ doublets that passes the
printed filters carries at least $54$ Weyl degrees of freedom over
three generations --- strictly above the 45-DOF Standard-Model
template.  The class minimum is attained by the conjugate-pair
template
\[
(\mathbf{3},\mathbf{2})_{Y_1} + (\bar{\mathbf{3}},\mathbf{2})_{-Y_1} +
(\mathbf{3},\mathbf{1})_{Y_2} + (\bar{\mathbf{3}},\mathbf{1})_{-Y_2},
\qquad \mathrm{DOF} = (6+6+3+3) \times 3 = 54,
\]
which passes every printed filter: $b_0^{(3)} = 5 > 0$ (F1),
$b_0^{(2)} = 4/3 > 0$ (F2), the F3 content predicate (both types
present), exact $[\SU(3)]^3$ cancellation (F4), six doublets counted
with color multiplicity --- even (F5), and the full anomaly system,
which vanishes identically for \emph{any} rational $(Y_1, Y_2)$ with
the conjugate structure, non-degenerately for $Y_1 \neq 0$ (F6).  The
class bound is a direct enumeration of the two-colored-doublet class
under the printed filters, in exact rational arithmetic, carried in
the release's audit layer.  Two earlier claims are withdrawn as
false: the premise that two colored doublets force lepton content
(the conjugate-pair template is anomaly-free with no leptons at all),
and the 63-DOF ``tightest template'' (its seven doublet components,
counted with color multiplicity, are odd --- it fails the Witten
filter F5 and was never admissible).  What P4 contributes to the
search-space bound is dominance: $54 > 45$, so no two-doublet
template can beat the Standard-Model minimum.
\end{proposition}

\begin{proposition}[P5: Extra field types cannot lower the minimum]
\label{prop:P5}
Same dominance form as P3.  The a-priori per-type floor ($\ge 3$ DOF:
a colorless singlet across 3 generations) does \emph{not} by itself
bound many-type templates near 45 --- six minimal types floor at only
18 DOF, and the v3.9 wording (``$\ge 45 + 3$'') presupposed the
five-type baseline it was meant to establish.  The valid statement:
within the declared space, every F6 survivor containing types beyond
the five SM roles carries them as additions to an anomaly-closed core
(the scan's survivor list exhibits exactly this --- the 6- and 7-entry
survivors sit at 48--72 DOF, each strictly above the 45-DOF core it
extends), so extra types never lower the minimum.  Completeness of the
declared space itself is a separate obligation, stated with the
enumeration below.
\end{proposition}

\noindent P1--P2 are theorems bounding the representation content
(AF-forced): $\SU(3)$ reps $\in \{3, \bar{3}, 6, \bar{6}, 8\}$
(excluding singlets, which are colorless) and, for colored fields,
$\SU(2)$ reps $\in \{1, 2\}$.  The remaining bounds are \emph{declared
enumeration caps}, stated here as scan inputs rather than theorems:
one colored $\SU(2)$ doublet (P4 certifies dominance against the
two-doublet enlargement direction), at most 3 colored singlets, at most one lepton
doublet, at most 2 lepton singlets --- hence template size
$|\mathcal{T}| \le 7$ (the v3.9 text said $|\mathcal{T}| \le 5$, which
contradicted its own enumeration: the survivor list contains 6- and
7-entry templates).  Completeness of these caps --- that no admissible
sub-45-DOF template exists outside them --- is carried by P3/P5-style
dominance for the enumerated directions and is otherwise an explicit
proof obligation of the reproducibility contract, not a closed
theorem.
Combined with multiplicity choices (0--3 colored singlets drawn
with replacement, 0--2 lepton singlets, doublet present or absent),
this yields 1{,}680 distinct candidate templates.\footnote{Colored
singlets are unordered within a template, so the correct enumeration
uses combinations with replacement: $5 \times (1 + 5 + 15 + 35)
\times 2 \times 3 = 1{,}680$.  A standalone verification script
with exact rational arithmetic and four built-in red-team challenges
is archived at \url{https://doi.org/10.5281/zenodo.19154197}; an
interactive Colab walkthrough is at
\url{https://colab.research.google.com/drive/1otO1Qwp_Lr3OsK6ykkeMRATanecPkKg6};
source code is at
\url{https://github.com/Ethan-Brooke/APF-Paper-2-The-Structure-of-Admissible-Physics}.}


\subsection{Seven filters}
\label{sec:seven_filters}

Each filter is an exact declared predicate, stated at its ledger type
in Theorem~\ref{thm:conditional_classification}: F1--F2 consume the
imported one-loop Weyl coefficient and the imported generation count
(C3--C4); F3 is the doublet--singlet content predicate (C6); F4--F6 are the imported
anomaly and Witten equations consumed as filter conditions (C5), with
F6 additionally carrying the independent electromagnetic-completeness
assumption in its canonical full-system form (C8); F7 is the CPT
convention (C7).  Their APF motivations live where cited --- the
asymptotic-freedom discussion of \S\ref{sec:beta_schematic} and the
foundational program of \TSII{} --- and the classification consumes
them as stated inputs.  The filters are applied sequentially; each
eliminates candidates that passed all previous filters.

\begin{enumerate}[label=\textbf{F\arabic*},nosep,leftmargin=3em]
  \item \textbf{$\SU(3)$ asymptotic freedom.}\quad
  $b_0^{(3)} = 11 - (2/3) \sum_i T(r_i^{(3)}) \dim(r_i^{(2)}) \cdot N_{\mathrm{gen}} > 0$,
  with the sum over the template's left-handed Weyl entries ($2/3$ is
  the Weyl coefficient; the verifier has always used it --- the $4/3$
  printed here in v3.9 was the Dirac-pair form, a display error that,
  with Weyl entries, mis-prices the SM at $b_0 = 3$ instead of $7$).
  $N_{\mathrm{gen}} = 3$ is imported (\S\ref{sec:three_generation_scope});
  F1--F2 verdicts are conditional on it.
  Derived from non-abelian self-interaction dominance (\S\ref{sec:beta_schematic}).

  \item \textbf{$\SU(2)$ asymptotic freedom.}\quad
  $b_0^{(2)} = 22/3 - (2/3) \sum_i T(r_i^{(2)}) \dim(r_i^{(3)}) \cdot N_{\mathrm{gen}} > 0$
  (same Weyl convention and $N_{\mathrm{gen}}$ conditionality as F1).

  \item \textbf{Doublet--singlet content.}\quad
  The template must contain both colored doublets and colored
  singlets.  This is a content predicate, not a chirality test, and
  we name it at its content: co-presence of the two types removes
  doublet-only and singlet-only templates, but it does not by itself
  exclude vector-like content --- the conjugate-pair two-doublet
  template of Proposition~\ref{prop:P4} contains both types and is
  exactly vector-like (every field acquires a gauge-invariant
  conjugate mass partner).  Within the scanned space, which carries
  exactly one colored doublet by construction, no conjugate
  colored-doublet partner can occur, and the chirality of the
  surviving assignment is settled downstream, at the charge level, by
  F6's non-degeneracy ($Y_Q \neq 0$).  A true chirality predicate ---
  the absence of any gauge-invariant rep-conjugate pairing with
  opposite hypercharge --- is charge-dependent and evaluable only
  after hypercharges are assigned; it is reported separately in the
  release's audit layer rather than consumed as a scan filter.

  \item \textbf{$[\SU(3)]^3$ anomaly cancellation.}\quad
  $\sum_i A(r_i^{(3)}) \cdot \dim(r_i^{(2)}) = 0$, where $A$ is the
  cubic anomaly coefficient ($A(\mathbf{3}) = 1/2$,
  $A(\bar{\mathbf{3}}) = -1/2$, $A(\mathbf{6}) = 7/2$,
  $A(\bar{\mathbf{6}}) = -7/2$, $A(\mathbf{8}) = 0$).

  \item \textbf{Witten anomaly freedom.}\quad
  The total number of $\SU(2)$ doublets (counted with $\SU(3)$ multiplicity)
  must be even: $\sum_i \dim(r_i^{(3)}) \cdot [r_i^{(2)} = \mathbf{2}]
  \equiv 0 \pmod{2}$.

  \item \textbf{Full anomaly system solvable.}\quad
  The combined system of $[\SU(2)]^2[\U(1)]$, $[\SU(3)]^2[\U(1)]$,
  $[\mathrm{grav}]^2[\U(1)]$, and $[\U(1)]^3$ anomaly cancellation
  conditions must admit rational hypercharge solutions.  Templates
  containing colored singlets with $A(R) = 0$ (e.g.\ the $\mathbf{8}$)
  are handled by spectator reduction: the spectator-reduced branch
  sets $Y = 0$ on such fields, decoupling them from the anomaly
  conditions, and the reduced template is then tested analytically.
  (The \emph{full} anomaly system, solved without spectator reduction,
  additionally admits charged-octet solutions with
  $Y_8 = -\tfrac{3}{5}\,Y_Q$; they are eliminated at the minimality
  cut --- see the predicate statement below.)  For the spectator-free core, the system reduces
  to a quadratic in $z = Y_u/Y_Q$ having rational roots
  (\S\ref{sec:hypercharge}).  This reduction presupposes a hypercharged
  quark doublet ($Y_Q \neq 0$), and that non-degeneracy condition is
  load-bearing rather than incidental.  Read weakly --- as the bare
  existence of \emph{some} anomaly-free assignment --- F6 would admit a
  lighter spurious survivor: the Standard-Model generation stripped of its
  charged lepton singlet, $\{Q(\mathbf{3},\mathbf{2}),\,
  \bar\psi_a(\bar{\mathbf{3}},\mathbf{1}),\,
  \bar\psi_b(\bar{\mathbf{3}},\mathbf{1}),\, L(\mathbf{1},\mathbf{2})\}$
  at $42 < 45$ Weyl degrees of freedom, which is chiral and passes
  F1--F5, and is anomaly-free under the degenerate assignment
  $Y_Q = Y_L = 0$, $Y = \pm t$ on the two colored singlets.  The framework
  excludes it by \textbf{electromagnetic completeness} --- every coloured
  multiplet must carry a separating charge, so $Y_Q \neq 0$, and the
  spectator-free anomaly system is therefore required to be
  \emph{non-degenerate}, not merely solvable.  Electromagnetic
  completeness is hereby declared an \emph{independent assumption} of
  the scan at exactly this point of use: it is motivated by the
  admissibility semantics (EC's separating-label requirement applied
  to the unbroken electromagnetic charge), but it is not a consequence
  of anomaly cancellation, and it does the work of removing the 42-DOF
  near-miss.  A reader who rejects it gets the stripped generation as
  a second survivor below the SM.  The stripped template is recorded
  as the sharpest near-miss in \texttt{near\_misses.csv}.
  \emph{Predicate statement for re-implementers:} the
  \textbf{canonical} F6 predicate --- the convention generating this
  document's waterfall table and every count in the body --- is
  \emph{full-system non-degeneracy}: solve the complete anomaly
  system, spectators included, and require a non-degenerate rational
  solution ($Y_Q \neq 0$) --- 10 survivors pre-CPT, 5 post-CPT.  The
  full system admits charged-octet solutions, with
  $Y_8 = -\tfrac{3}{5}\,Y_Q$ forced linearly on every solution of the
  DOF-66 template (e.g.\ $(Y_Q, Y_u, Y_d, Y_8, Y_L) =
  (5, 13, -5, -3, -15)$), and SM-embedding solutions on the 69- and
  72-DOF templates; all are eliminated at the minimality cut.  Two
  variant predicates are in circulation across the artifacts and are
  reported as robustness checks with their counts
  (\S\ref{sec:waterfall}): (i)~\emph{spectator reduction}
  (implemented by the shipped standalone scan, v2, Zenodo
  \texttt{10.5281/zenodo.19154197}): strip anomaly-free spectators
  ($A(R) = 0$ colored singlets, set to $Y = 0$ by the reduced system)
  and require a non-degenerate rational solution on the remainder ---
  8 pre-CPT, 4 post-CPT; (ii)~\emph{uniform dimension} (the in-tree
  code path: \texttt{check\_T\_field} in \texttt{gauge.py} requires
  every colored singlet to have $\SU(3)$ dimension $N_c = 3$) --- 4
  pre-CPT, 2 post-CPT.  The winner (minimum DOF 45) is identical under
  all three predicates; only the intermediate counts
  differ\coderef{check\_L\_F6\_not\_from\_EC}{ec\_inventory\_reading.py}.
  The predicate inequivalence is witnessed by the DOF-66 charged
  octet: it survives the canonical full-system predicate and is absent
  under spectator reduction.
  One further computed fact sharpens the independence declaration
  above: EC-label-injectivity does \emph{not} remove the 42-DOF
  near-miss --- its unique nontrivial solution $(0, t, -t, 0)$ is
  label-injective and per-field faithful --- so F6's independence from
  EC is a computed theorem over this scan space, not only a
  declaration.

  \item \textbf{CPT quotient.}\quad
  Templates related by charge conjugation (replacing every representation
  with its conjugate) describe the same physics.  The CPT equivalence
  class is identified and only the canonical representative is counted.
\end{enumerate}


\subsection{Waterfall table}
\label{sec:waterfall}

\begin{center}
\small
\resizebox{\textwidth}{!}{\begin{tabular}{@{}lrrl@{}}
\toprule
\textbf{Filter} & \textbf{Survivors} & \textbf{Eliminated} & \textbf{Mechanism} \\
\midrule
Search space (P1--P5) & 1{,}680 & --- & Bounded by exclusion proofs \\
F1+F2: Asymptotic freedom & 348 & 1{,}332 & $b_0^{(3)} > 0$ and $b_0^{(2)} > 0$ \\
F3: Doublet--singlet content & 324 & 24 & Colored doublets + singlets required \\
F4: $[\SU(3)]^3$ anomaly & 24 & 300 & $\sum A_3 d_2 = 0$ \\
F5: Witten anomaly & 12 & 12 & $\#$doublets even \\
F6: Full anomaly system & 10 & 2 & Full-system non-degenerate rational hypercharges ($Y_Q \neq 0$) \\
F7: CPT quotient & 5 & 5 & Conjugate pairs identified \\
Minimality & \textbf{1} & 4 & Min DOF = 45 selected \\
\bottomrule
\end{tabular}}
\end{center}

\smallskip\noindent\textit{Counting convention (narrative vs.\ enumeration).}\quad
The counts in the table above are the number of \emph{distinct}
templates surviving each filter, with unordered colored-singlet
positions, under the canonical full-system F6 predicate.  Both
implementations (\texttt{check\_T\_field} in \texttt{gauge.py} and
\texttt{fermion\_scan\_standalone.py}) use an ordered-tuple
enumeration (\texttt{itertools.product}) and therefore test each
distinct template multiple times, producing a raw iteration count of
$4{,}680$ rather than $1{,}680$.  Their F6 conventions are the two
robustness variants: the standalone scan's spectator reduction gives
$8 \to 4$, and \texttt{check\_T\_field}'s uniform-dimension
convention gives $4 \to 2$ after quotienting the ordering redundancy
(the predicate statement in \S\ref{sec:seven_filters} states all
three conventions and designates the canonical one).  The minimality
winner is identical under either enumeration and every stated F6
convention: the Standard Model at 45 DOF.  A reviewer running the code
will see the enumeration multiplicities and the robustness-variant
counts; the table reports the unordered canonical count because that
is what the five exclusion proofs (P1--P5) actually bound.

\smallskip\noindent\textit{Robustness: the two variant F6 predicates.}\quad
Templates \#5--\#10 in the post-F6 listing below include the adjoint
$\mathbf{8}$ as a colored $\SU(2)$ singlet.  Under the canonical
full-system predicate these templates survive F6 --- the DOF-66 pair
via its charged-octet solutions ($Y_8 = -\tfrac{3}{5}\,Y_Q$), the 69-
and 72-DOF pairs via their SM-embedding solutions --- and fail only at
the minimality cut.  Two variant predicates are reported as robustness
checks.  \emph{Spectator reduction}, implemented by the shipped
standalone scan (\texttt{fermion\_scan\_standalone.py}), strips the
anomaly-free spectator ($A(\mathbf 8)=0$, set to $Y = 0$ by the
reduced system) and tests the remainder: 8 survivors pre-CPT, 4
post-CPT --- the DOF-66 pair drops, because its spectator-reduced
remainder is the degenerate 42-DOF stripped generation; the charged
octet is exactly the predicate-inequivalence witness.  \emph{Uniform
dimension}, the in-tree code path (\texttt{check\_T\_field} in
\texttt{gauge.py}), requires every colored singlet at $\SU(3)$
dimension $N_c = 3$, eliminating octet templates at F6: 4 pre-CPT, 2
post-CPT.  All three predicates reproduce the same winner --- the SM
at DOF 45 --- at the minimality cut; the intermediate counts are
10/5 (canonical), 8/4, and 4/2.

\smallskip\noindent The final winner is exact and reproduced by
\texttt{check\_T\_field} in \texttt{gauge.py} and independently by
\texttt{fermion\_scan\_standalone.py}~\cite{BrookeFermionScan}.

\begin{proposition}[Filter-order independence]\label{prop:filter_order}
The survivor set $\mathcal{S} = \{\mathcal{T} : \mathcal{T} \text{
passes F1--F7}\}$ is independent of the order in which the seven
filters are applied.
\end{proposition}

\begin{proof}
Each filter $F_k$ defines a predicate $P_k: \mathcal{T} \to
\{0, 1\}$ on templates.  A template survives all filters iff it
satisfies the conjunction $P_1 \wedge P_2 \wedge \cdots \wedge P_7$.
Logical conjunction is commutative and associative:
$P_{\sigma(1)} \wedge \cdots \wedge P_{\sigma(7)} = P_1 \wedge
\cdots \wedge P_7$ for any permutation $\sigma \in S_7$.  Therefore
the survivor set $\mathcal{S} = \{\mathcal{T} : \bigwedge_k
P_k(\mathcal{T}) = 1\}$ is identical regardless of evaluation order.

What changes with filter order is the \emph{waterfall count} (number
of templates eliminated at each stage), because a template eliminated
by $F_k$ in the standard order might instead be eliminated by
$F_j$ with $j > k$ if the order is reversed.  The final survivor
set, however, is order-invariant.  The standalone verification
script confirms this under its spectator-reduction F6 predicate:
running all $7! = 5{,}040$ permutations on the 1{,}680-template set
yields the same 4 post-CPT survivors (and the same unique DOF-45
winner) in every case.  The conjunction argument is
predicate-independent, so the invariance holds equally for the
canonical full-system predicate's 5 post-CPT survivors.
\end{proof}

\smallskip\noindent\textit{Search-space bias and completeness.}\quad
A distinct concern is whether the search space itself (the 1{,}680
templates bounded by P1--P5) introduces bias by excluding candidates
that might survive the filters.  P1--P2 are unconditional theorems: a
single field in an $\SU(3)$ representation of dimension $\ge 10$, or a
colored field in an $\SU(2)$ representation of dimension $\ge 3$,
exceeds the AF budget regardless of what else the template contains,
and the AF and DOF bounds are monotone (adding fields only increases
DOF and shrinks the AF margin).  The remaining bounds are declared
enumeration caps (\S\ref{sec:exclusion_proofs}): P3 and P5 are
dominance results certifying that enlarging the declared space in the
enumerated directions cannot lower the minimum, and P4 is the
two-colored-doublet class-dominance enumeration (class minimum 54,
computed over the two-doublet class --- outside the caps, which admit
one colored doublet).  The licensed claim is
therefore minimality \emph{among the enumerated candidates}, together
with dominance against the named enlargement directions.  Completeness
of the caps --- that no admissible sub-45-DOF template exists outside
them --- is an explicit proof obligation of the reproducibility
contract, not a closed theorem; the proof obligations named there are
the open program.

\smallskip\noindent\textit{Explicit listing of all post-F6 survivors.}\quad
The 10 templates surviving F1--F6 under the canonical full-system
predicate (before CPT quotient) are:

\smallskip
\begin{center}\small
\renewcommand{\arraystretch}{1.15}
\begin{tabular}{@{}clcl@{}}
\toprule
\# & Template & DOF & Status \\
\midrule
1 & $(3,2) + (\bar{3},1) + (\bar{3},1) + (1,2) + (1,1)$ & 45 & SM \\
2 & $(\bar{3},2) + (3,1) + (3,1) + (1,2) + (1,1)$ & 45 & CPT of \#1 \\
3 & $(3,2) + (\bar{3},1) + (\bar{3},1) + (1,2) + (1,1) + (1,1)$ & 48 & SM + extra singlet \\
4 & $(\bar{3},2) + (3,1) + (3,1) + (1,2) + (1,1) + (1,1)$ & 48 & CPT of \#3 \\
5 & $(3,2) + (\bar{3},1) + (\bar{3},1) + (8,1) + (1,2)$ & 66 & Charged octet ($Y_8 = -\tfrac{3}{5}Y_Q$) \\
6 & $(\bar{3},2) + (3,1) + (3,1) + (8,1) + (1,2)$ & 66 & CPT of \#5 \\
7 & $(3,2) + (\bar{3},1) + (\bar{3},1) + (8,1) + (1,2) + (1,1)$ & 69 & SM + octet (SM-embedding) \\
8 & $(\bar{3},2) + (3,1) + (3,1) + (8,1) + (1,2) + (1,1)$ & 69 & CPT of \#7 \\
9 & $(3,2) + (\bar{3},1) + (\bar{3},1) + (8,1) + (1,2) + (1,1) + (1,1)$ & 72 & SM + octet + singlet \\
10 & $(\bar{3},2) + (3,1) + (3,1) + (8,1) + (1,2) + (1,1) + (1,1)$ & 72 & CPT of \#9 \\
\bottomrule
\end{tabular}
\end{center}

\smallskip\noindent Templates \#5--\#10 contain the adjoint
$\mathbf{8}$ as a colored $\SU(2)$ singlet.  Under the canonical
full-system predicate they carry the anomaly solutions recorded in
the predicate statement (the DOF-66 pair with the octet charged at
$Y_8 = -\tfrac{3}{5}\,Y_Q$; the 69- and 72-DOF pairs with
SM-embedding solutions).  Their DOF cost ($8 \times 3 = 24$ per
generation for the octet alone) is far above the minimum.

\smallskip\noindent After the CPT quotient (F7), templates \#2, \#4,
\#6, \#8, and \#10 are identified with \#1, \#3, \#5, \#7, and \#9
respectively, leaving 5 CPT-inequivalent survivors.  Minimality selects DOF $= 45$,
giving exactly one winner: the Standard Model --- under F6
(electromagnetic completeness, an independent assumption,
\S\ref{sec:seven_filters}); the F6-free argmin is the 42-DOF stripped
generation recorded there as the sharpest near-miss.  The second-lightest
survivor (\#3, DOF $= 48$) differs from the SM by one additional
right-handed neutrino-like singlet per generation;
its exclusion is by minimality ($48 > 45$); a bridge-conditional
observational cross-check is recorded at its corpus home
(\S\ref{sec:neutrinos}).


\subsection{The unique survivor}
\label{sec:unique_survivor}

\begin{theorem}[$T_{\mathrm{field}}$: SM Fermion Content]\label{thm:Tfield}
The unique minimal chiral fermion template consistent with
$\SU(3) \times \SU(2) \times \U(1)$, three generations,
electromagnetic completeness (F6 --- an independent assumption,
declared at its point of use in \S\ref{sec:seven_filters}, in its
canonical full-system non-degeneracy form), and all
remaining admissibility constraints is:
\[
  \{Q(\mathbf{3},\mathbf{2}),\;
    u^c(\bar{\mathbf{3}},\mathbf{1}),\;
    d^c(\bar{\mathbf{3}},\mathbf{1}),\;
    L(\mathbf{1},\mathbf{2}),\;
    e^c(\mathbf{1},\mathbf{1})\}
  \times 3 \text{ generations},
\]
with $\mathrm{DOF} = (6+3+3+2+1) \times 3 = 45$ Weyl fermions.
\end{theorem}

\noindent No other template among the 1{,}680 candidates has
$\mathrm{DOF} \le 45$ and passes all seven filters.  P1--P2 exclude
the enumerated representation enlargements unconditionally; P3--P5
certify dominance for the enumerated extension directions; the caps
themselves are declared scan inputs whose completeness is an open
proof obligation (\S\ref{sec:exclusion_proofs}).  The licensed
statement: the SM fermion content is the unique minimum among the
enumerated candidates, with dominance against the named enlargements.
\coderef{check\_T\_field}{gauge.py}; standalone verification: \texttt{fermion\_scan\_standalone.py}~\cite{BrookeFermionScan}

The scan result and its premise ledger assemble into the
supplement's citable core, stated as one theorem with the full input
list typed --- theorem, imported result, declared input, reading, or
convention.

\begin{theorem}[Conditional classification of the fermion template]
\label{thm:conditional_classification}
Fix the following premises --- the full input ledger of this theorem.
\begin{enumerate}[label=(C\arabic*),nosep]
  \item \textbf{Declared scan inputs.}\quad The template caps and
  representation tables of \S\ref{sec:exclusion_proofs} (P1--P2 as
  theorems; the enumeration caps as declared inputs).
  \item \textbf{The gauge template.}\quad The fixed template
  $\SU(3)\times\SU(2)\times\U(1)$, consumed through
  $T_{\mathrm{field}}$'s GIVEN (Theorem~\ref{thm:Tfield}).
  \item \textbf{Generation count.}\quad $N_{\mathrm{gen}} = 3$,
  imported (\S\ref{sec:three_generation_scope}) and consumed by the
  F1--F2 filters.
  \item \textbf{One-loop Weyl beta coefficient.}\quad The
  coefficient $2/3$ per left-handed Weyl entry in the F1--F2
  asymptotic-freedom filters, imported from standard gauge theory
  (\S\ref{sec:beta_schematic}).
  \item \textbf{The anomaly system.}\quad The four gauge and
  gravitational anomaly equations and the Witten parity condition
  (F4--F6 and F5), consumed as filter conditions; their solutions are
  computed exactly over the rationals, not assumed.
  \item \textbf{The F3 content predicate.}\quad F3's
  doublet--singlet co-presence condition: colored doublets and
  colored singlets both present.  A content predicate, not a
  chirality test (\S\ref{sec:seven_filters} states its honest role;
  chirality of the winner is settled at F6's non-degeneracy).
  \item \textbf{CPT identification.}\quad F7's convention that
  conjugate templates describe the same physics.
  \item \textbf{Electromagnetic completeness (F6).}\quad An
  independent assumption, declared at its point of use
  (\S\ref{sec:seven_filters}), in its canonical full-system
  non-degeneracy form ($Y_Q \neq 0$ on the complete anomaly system);
  without it the argmin moves to the 42-DOF stripped generation.  F6
  and EC-label-injectivity are logically independent --- computed, in
  both directions\coderef{check\_L\_F6\_not\_from\_EC}{ec\_inventory\_reading.py}.
  \item \textbf{The objective.}\quad Weyl-component count at the
  left-handed Weyl convention: the quantity minimized is total Weyl
  degrees of freedom.
  \item \textbf{Cost convention.}\quad Equal unit cost per counted
  channel --- one capacity slot of equal weight per Weyl component,
  Higgs component, and gauge generator (BA5,
  \S\ref{sec:bridge_axioms}; $L_{\mathrm{Cauchy}}$'s unit
  normalization $F(n)=n$) --- with the gauge-cost equality
  $C(G) = \dim(G)$ resting on the named operational premise pair
  (operational encoding; disjoint-anchor realizability;
  Proposition~\ref{prop:generator_separation}).
\end{enumerate}
Under (C1)--(C10), the Standard-Model five-multiplet template at 45
Weyl degrees of freedom is the unique minimum among the enumerated
candidates, computed
exactly\coderef{check\_T\_field}{gauge.py}.

Two remarks fix the premises' status honestly.  First, within APF the
premises (C2) and (C3) have banked grounds --- the gauge-uniqueness
chain (\S\ref{sec:gauge_class_supp}, resting on the foundational
program's carrier theorems) and the capacity--beta identities with
their coderefs (\S\ref{sec:three_generation_scope}) --- but this
theorem consumes them as \emph{inputs}: its truth does not certify
them, and their grounds live where cited.  Second, the abelian-factor
count is likewise not certified here.  The one-U(1) results
(\S\ref{sec:one_u1_closure}) are a-posteriori closure theorems for
the selected template: they consume the template and the scanned
matter content and verify consistency and minimal-label uniqueness of
the abelian sector, under the reading R-EC-inv with the entry-typed
inventory
I-typing\coderef{check\_L\_EC\_inventory\_reading}{ec\_inventory\_reading.py};
neither this theorem nor those derive the other's input.  Of the
premises: (C8) is an independent assumption; (C1)'s caps, (C7), (C9),
and (C10) are declared conventions of the enumeration and the ledger;
(C4) and (C5) are imported standard results; (C2) and (C3) are
imported APF results consumed as inputs.
\end{theorem}

\noindent The theorem statements upstream and downstream of the scan
(Theorem~\ref{thm:Tfield}, the waterfall conclusion of
\S\ref{sec:waterfall}, Theorem~\ref{thm:unique_U1_consolidated}) carry
their conditionality clauses individually; this consolidated statement
is the one to cite.


\newpage


% ══════════════════════════════════════════════════════════════
\section{Anomaly cancellation and hypercharge uniqueness}
\label{sec:hypercharge}
% ══════════════════════════════════════════════════════════════


\begin{tcolorbox}[apfbox, title={Novelty boundary: standard anomaly algebra versus APF input selection}]
\footnotesize
The anomaly equations used in this section are standard Standard-Model
consistency equations.  The APF novelty is not the algebraic fact that,
given the SM chiral template, anomaly cancellation fixes hypercharge up to
normalization.  The novelty claim is upstream: the finite admissibility
argument selects the gauge template and the anomaly-tested fermion template
before the standard equations are applied.  Thus the anomaly calculation is
a closure check on APF-selected inputs, not an independent new anomaly
formalism.
\end{tcolorbox}

\noindent\textit{What is standard and what is new.}\quad
The anomaly equations in this section are standard four-dimensional gauge
anomaly conditions applied to a chiral fermion template.  APF does not
claim novelty for the algebraic fact that, once the SM representation
pattern is assumed, anomaly cancellation fixes a rational hypercharge ray
(up to normalization and charge conjugation).  The APF contribution is
upstream: the gauge template, the representation pattern, the absence of
exotic fermions, the single abelian factor, the generation count entering
the ledger, and the minimum-capacity choice of the 45-Weyl survivor are
all fixed before the anomaly system is solved.  Hypercharge uniqueness is
therefore a sanity check and closure theorem for the derived template, not
a disguised input of the template.

% ── §7: Anomaly cancellation and hypercharge uniqueness ──

\subsection{The four anomaly conditions}
\label{sec:anomaly_conditions}

The quantum consistency of a chiral gauge theory requires the
cancellation of all gauge anomalies.  For
$\SU(3) \times \SU(2) \times \U(1)$ with the surviving template, there
are four independent conditions on the five hypercharges
$\{Y_Q, Y_u, Y_d, Y_L, Y_e\}$.  Three are linear:
\begin{align}
  [\SU(2)]^2[\U(1)] &= 0: &
    N_c\, Y_Q + Y_L &= 0, \label{eq:anom1} \\
  [\SU(3)]^2[\U(1)] &= 0: &
    2Y_Q - Y_u - Y_d &= 0, \label{eq:anom2} \\
  [\mathrm{grav}]^2[\U(1)] &= 0: &
    2N_c\, Y_Q + 2Y_L - N_c\, Y_u - N_c\, Y_d - Y_e &= 0. \label{eq:anom3}
\end{align}
The fourth is the cubic anomaly:
\begin{equation}
  [\U(1)]^3 = 0: \quad
  2N_c\, Y_Q^3 + 2Y_L^3 - N_c\, Y_u^3 - N_c\, Y_d^3 - Y_e^3 = 0.
  \label{eq:anom4}
\end{equation}

\subsection{Reduction to a quadratic}
\label{sec:reduction_quadratic}

Solving (\ref{eq:anom1})--(\ref{eq:anom3}) gives
$Y_L = -N_c\, Y_Q$, $Y_d = 2Y_Q - Y_u$, $Y_e = -2N_c\, Y_Q$.
Substituting into (\ref{eq:anom4}) and dividing by $Y_Q^3 \ne 0$
(non-trivial grading), define $z := Y_u / Y_Q$:
\begin{equation}
  z^2 - 2z - (N_c^2 - 1) = 0.
  \label{eq:z_quadratic}
\end{equation}

\noindent\textit{Why the roots are rational.}\quad
Equation~(\ref{eq:z_quadratic}) is a \emph{monic} polynomial (leading
coefficient $= 1$) with integer coefficients ($-2$ and $-(N_c^2-1)$,
where $N_c \in \mathbb{Z}$).  By the rational root theorem for monic
integer polynomials: if $z = p/q$ in lowest terms is a root, then
$q \mid 1$, so $z \in \mathbb{Z}$.  The discriminant
$\Delta = 4 + 4(N_c^2 - 1) = 4N_c^2$ is a perfect square, confirming
integer roots:
\begin{equation}
  z = 1 \pm N_c.
\end{equation}
Since $Y_L = -N_c Y_Q$, $Y_d = 2Y_Q - Y_u = (2-z)Y_Q$, and
$Y_e = -2N_c Y_Q$ are all integer multiples of $Y_Q$, and $Y_Q$ itself
is a free normalization (any nonzero rational value is admissible, since
it sets the overall scale of the $\U(1)$ generator; the
compact-$\U(1)$ representation-data premise of the front-matter
ledger), choosing
$Y_Q = 1/6$ gives $Y_i \in \mathbb{Q}$ for all five multiplets.
What the anomaly system proves is that all hypercharge \emph{ratios}
are rational --- the solution is a unique rational ray, up to the
$u \leftrightarrow d$ relabeling.  The absolute unit is fixed by the
compactness of the $\U(1)$ and the character-lattice convention
($Y_Q = 1/6$; the compact-$\U(1)$ representation-data premise of the
front-matter ledger), not by anomalies alone.
For $N_c = 3$: $z \in \{4, -2\}$.  These are related by
$z \to 2 - z$ (the $u \leftrightarrow d$ relabeling), so they describe
the same physics.  Choosing $z = 4$ and normalizing $Y_Q = 1/6$:

\begin{theorem}[$L_{\mathrm{hypercharge}}$: Unique Hypercharge Assignment]
\label{thm:hypercharge}
\[
  Y_Q = \tfrac{1}{6}, \quad
  Y_u = \tfrac{2}{3}, \quad
  Y_d = -\tfrac{1}{3}, \quad
  Y_L = -\tfrac{1}{2}, \quad
  Y_e = -1.
\]
These are the unique hypercharges (up to normalization and $u \leftrightarrow d$)
consistent with anomaly cancellation for the derived template with $N_c = 3$.
\end{theorem}

\noindent\textit{Consequences.}\quad Electric charges $Q = T_3 + Y$ give
$Q(u) = 2/3$, $Q(d) = -1/3$, $Q(e) = -1$, $Q(\nu) = 0$.  Charge
\emph{ratios} are fixed by the anomaly structure; with the
compact-$\U(1)$ unit convention $Y_Q = 1/6$, all electric charges
land on multiples of $1/3$.  The ratios are derived; the unit is the
convention (\S\ref{sec:one_u1_closure}).  The quark--lepton relation
$Y_L = -N_c\, Y_Q$ ties lepton hypercharge to the number of colors.
\coderef{check\_T\_field (hypercharge section)}{gauge.py}


\newpage


% ══════════════════════════════════════════════════════════════


\section{The one-U(1) a-posteriori closure theorems}
\label{sec:one_u1_closure}

The abelian sector's uniqueness results are stated here as
\emph{a-posteriori closure theorems} for the selected template.  The
logical order is: gauge template (consumed at C2 of
Theorem~\ref{thm:conditional_classification}) $\to$
$T_{\mathrm{field}}$ (the scan, \S\ref{sec:fermion_scan}) $\to$ the
closure theorems below $\to$ $L_{\mathrm{hypercharge}}$
(\S\ref{sec:hypercharge}).  Proposition~\ref{prop:unique_U1} is
template-independent: for \emph{any} fixed set of non-abelian
multiplets $\mathcal{M}$ satisfying EC and CM, the number of
admissible $\U(1)$ factors is determined by the multiplet separation
problem.  Theorem~\ref{thm:unique_U1_consolidated} is its SM
specialization.

There is no circularity, and the claim structure is worth stating
plainly.  The scan takes the template --- with one $\U(1)$ available
to grade --- as an input: F6 solves hypercharge assignments on that
template, so the scan consumes the abelian slot, not only
$\SU(3)\times\SU(2)$ quantum numbers.  The closure theorems then
verify, for the scanned matter content, that the abelian sector is
consistent and that its minimal labeling is unique: at least one
abelian grading is necessary (under R-EC-inv $+$ I-typing), one
suffices, and no second is admissible under CM.  Neither result
claims to derive the other's input: the scan does not derive the
template, and the closure theorems do not derive the matter content.
A mutual-consistency reading is available as the honest strong form
--- the (template, matter content, single-abelian-grading) triple is a
simultaneous solution of the full constraint set --- and whether it is
the \emph{unique} simultaneous solution is the fixed-point program,
the head of the open-lane register of \TSII.

\smallskip\noindent The unique abelian factor result is split into two
levels.  \textbf{Theorem~A} (Proposition~\ref{prop:unique_U1}) is
template-independent: it shows that for \emph{any} multiplet set
$\mathcal{M}$ and \emph{any} block decomposition, EC and CM determine
the number of admissible abelian factors from the separation deficiency
$\delta$ alone.  \textbf{Theorem~B}
(Theorem~\ref{thm:unique_U1_consolidated}) is the SM specialization:
it plugs in $\{n_k\} = \{3, 2, 1\}$ and $\mathcal{M} =
\mathcal{M}_{\mathrm{SM}}$ and outputs ``exactly one $\U(1)$,
uniquely hypercharge.''

\begin{proposition}[Unique minimal abelian grading factor ---
structural (Theorem~A)]\label{prop:unique_U1}
Let
$\mathcal{A} \cong \bigoplus_{\alpha=1}^{m} M_{n_\alpha}(\mathbb{C})$
be the admissibility algebra at a fixed interface, with non-abelian gauge
factor
$G_{\mathrm{na}} = \prod_{\alpha=1}^{m} \SU(n_\alpha)$
and relative-phase torus
$T_{\mathrm{rel}} := Z(\mathrm{U}(\mathcal{A})) / \mathrm{U}(1)_{\mathrm{diag}}
\cong \mathrm{U}(1)^{m-1}$.
Let $\mathcal{M} = \{r_1, \ldots, r_N\}$ denote the set of physically
distinct matter multiplet types occurring in the derived template,
regarded as irreducible carriers of $G_{\mathrm{na}}$.

Assume that $G_{\mathrm{phys}}$ is an admissible gauge structure
(Definition~2.20 of \TSII{}; EC and CM hold by
Lemma~2.24 of \TSII{}).

\begin{lemma}[Abelian factors arise only from $T_{\mathrm{rel}}$]
\label{lem:abelian_source}
Every continuous abelian gauge factor of $G_{\mathrm{phys}}$ is a
one-dimensional subgroup of $T_{\mathrm{rel}}$.  No abelian gauge
factor can arise from any other source.
\end{lemma}

\begin{proof}
By the exact sequence (Proposition~2.19 of \TSII{}):
\[
  1 \;\to\; Z(\mathrm{U}(\mathcal{A}))
  \;\to\; \mathrm{U}(\mathcal{A})
  \;\xrightarrow{\;\pi\;}\;
  \mathrm{Inn}(\mathcal{A}) \;\to\; 1.
\]
$\mathrm{Inn}(\mathcal{A}) = \prod_\alpha \mathrm{PU}(n_\alpha)$ is
semisimple (each $\mathrm{PU}(n)$ is semisimple for $n \ge 2$, and
trivial for $n = 1$): its Lie algebra has no nonzero abelian
\emph{ideal}, equivalently the group has no connected abelian
\emph{normal direct factor}.  (It certainly contains continuous
abelian subgroups --- every maximal torus is one; what semisimplicity
excludes is an abelian factor that could serve as an independent gauge
direction with its own coupling.)  The abelian continuous subgroups of
$\mathrm{U}(\mathcal{A}) = \prod_\alpha \mathrm{U}(n_\alpha)$ that
commute with the whole group --- the candidates for an independent
abelian gauge factor --- are contained in
$\ker(\pi) = Z(\mathrm{U}(\mathcal{A}))$.
Quotienting by the physically unobservable diagonal
$\mathrm{U}(1)_{\mathrm{diag}}$ gives
$T_{\mathrm{rel}} = Z / \mathrm{U}(1)_{\mathrm{diag}}$ as the full
group of continuous abelian gauge transformations.  Any abelian gauge
factor is therefore a one-dimensional subgroup of $T_{\mathrm{rel}}$.

The argument excludes the following loopholes:
\begin{enumerate}[label=(\alph*),nosep]
  \item \textit{Abelian subgroups of $\mathrm{Inn}(\mathcal{A})$?}\quad
  $\mathrm{PU}(n)$ has a maximal torus of dimension $n - 1$ --- an
  abelian subgroup, but not an abelian \emph{normal factor}: its
  elements are non-central for $n \ge 2$, it is conjugated
  non-trivially, and its generators are already counted in
  $G_{\mathrm{na}}$ as Cartan directions with the non-abelian factor's
  single coupling.  A torus inside a simple factor is not an
  independent abelian gauge direction; only a normal abelian direct
  factor would be, and semisimplicity excludes exactly that.
  \item \textit{Discrete abelian factors?}\quad Discrete subgroups
  (e.g., $\mathbb{Z}_N$) can survive as global symmetries but do not
  contribute continuous gauge generators.  CM counts continuous
  admissibility channels; discrete symmetries have no per-generator
  capacity cost.
  \item \textit{External abelian factors?}\quad Any abelian factor not
  contained in $\mathrm{U}(\mathcal{A})$ would act non-trivially on
  $\mathcal{H}_\omega$ but not arise from the admissibility algebra.
  This contradicts the gauge identification theorem
  (Theorem~2.11 of \TSII{}): the gauge group is
  $\mathrm{Inn}(\mathcal{A})$ (plus the central extensions from the
  exact sequence), and no external group acts.
\end{enumerate}
\end{proof}

\begin{definition}[Separation deficiency]\label{def:sep_deficiency}
The \emph{separation deficiency} of $G_{\mathrm{na}}$ on $\mathcal{M}$
is the number of pairs of physically distinct multiplet types that
share identical non-abelian quantum numbers:
\[
  \delta(G_{\mathrm{na}}, \mathcal{M}) := \bigl|\{(r_i, r_j) : i < j,\;
  [r_i]_{G_{\mathrm{na}}} = [r_j]_{G_{\mathrm{na}}}\}\bigr|.
\]
EC is satisfied by $G_{\mathrm{na}}$ alone iff $\delta = 0$.  If
$\delta > 0$, additional abelian grading is required to resolve the
remaining degeneracies.
\end{definition}

\begin{lemma}[Minimum number of abelian factors from separation deficiency]
\label{lem:min_abelian}
Let $\mathcal{P} = \{P_1, \ldots, P_s\}$ be the partition of
$\mathcal{M}$ into equivalence classes under the non-abelian label
map (i.e., $P_k = \{r_i : [r_i]_{G_{\mathrm{na}}} = \ell_k\}$
for each distinct non-abelian label $\ell_k$).  Let
$m_{\max} := \max_k |P_k|$ be the size of the largest equivalence
class.  Then:
\begin{enumerate}[label=(\roman*),nosep]
  \item At least $\lceil \log_2 m_{\max} \rceil$ binary distinctions
  are needed to separate the elements of the largest class.
  \item A single $\U(1)$ factor assigns a real-valued charge $Y_i$ to
  each $r_i \in P_k$.  If the $Y_i$ are pairwise distinct within
  each class $P_k$, then one $\U(1)$ suffices to separate all
  multiplets.
  \item Therefore, the minimum number of abelian factors is $0$ if
  $\delta = 0$, and $1$ if $\delta > 0$ and there exists a charge
  assignment $Y$ that separates all classes.
\end{enumerate}
\end{lemma}

\begin{proof}
(i) is standard: $m_{\max}$ elements require $\lceil \log_2 m_{\max}
\rceil$ binary attributes to distinguish.  (ii): a single continuous
$\U(1)$ charge $Y \in \mathbb{R}$ has cardinality $\mathfrak{c}$, so
it can separate any finite number of elements --- the condition is that
the charge assignment $Y: P_k \to \mathbb{R}$ is injective on each
class.  (iii) combines (i) and (ii): if $m_{\max} \ge 2$, at least
one abelian factor is needed; if an injective $Y$ exists, one
suffices.
\end{proof}

\begin{lemma}[Diagonal rational grading from Schur + anomaly cancellation]
\label{lem:A3_derived}
Each $\U(1) \subset T_{\mathrm{rel}}$ assigns a rational charge to each
$r_i$, constant on that multiplet type.  Both properties are derived:
\begin{enumerate}[label=(\alph*),nosep]
  \item \textbf{Diagonal action (Schur).}\quad Each $r_i$ is an
  irreducible representation of $G_{\mathrm{na}} = \prod_\alpha
  \SU(n_\alpha)$.  The relative-phase torus
  $T_{\mathrm{rel}} \subset Z(\mathrm{U}(\mathcal{A}))$ lies in the
  center of the unitary group.  By Schur's lemma, any central element
  acts on an irreducible representation by a scalar.  Therefore each
  $\U(1) \subset T_{\mathrm{rel}}$ acts on $r_i$ by a phase
  $e^{i Y_i \theta}$, where $Y_i \in \mathbb{R}$ depends only on the
  multiplet type, not on the state within the multiplet.  This is
  diagonal action.
  \item \textbf{Rational charges (anomaly cancellation).}\quad The
  charges $Y_i$ are constrained by gauge anomaly cancellation
  (\S\ref{sec:hypercharge}): the four conditions
  $[\SU(2)]^2[\U(1)]$, $[\SU(3)]^2[\U(1)]$,
  $[\mathrm{grav}]^2[\U(1)]$, $[\U(1)]^3$ are polynomial equations
  in $\{Y_i\}$ with integer coefficients.  Solving the first three
  linearly and substituting into the fourth yields
  (Theorem~\ref{thm:hypercharge}):
  $z^2 - 2z - (N_c^2 - 1) = 0$, where $z := Y_u / Y_Q$.
  This polynomial is \emph{monic} (leading coefficient $= 1$) with
  integer coefficients.  By the rational root theorem for monic
  integer polynomials, any rational root $p/q$ (in lowest terms)
  must have $q \mid 1$, i.e., $z \in \mathbb{Z}$.  Indeed,
  $z = 1 \pm N_c$ are integers.  Since all other $Y_i$ are rational
  multiples of $Y_Q$ (from the linear conditions), and the standard
  normalization $Y_Q = 1/6$ (a \emph{convention}: the anomaly system
  is homogeneous and determines only the ray of charges, up to the
  discrete $u \leftrightarrow d$ exchange), all charge \emph{ratios}
  are rational.  What the anomaly system proves is rationality of
  ratios; the absolute charge unit is fixed by the compactness of the
  $\U(1)$ and the choice of character lattice made in the global-group
  ($\mathbb{Z}_6$-quotient) analysis, and choosing $Y_Q = 1/6$ is that
  convention, not a derivation of rationality from anomalies alone
  (the compact-$\U(1)$ representation-data premise; front-matter
  imported-ingredient ledger).
\end{enumerate}
\end{lemma}

\noindent Both properties were previously listed as assumption (A3).
The derivation shows they follow from Schur's lemma (standard
representation theory, no additional physical input) and anomaly
cancellation (already derived in this supplement).  No assumption
beyond EC, CM, the selected gauge template, and the scanned matter
content is needed.

Then:
\begin{enumerate}[label=(\roman*),nosep]
  \item If $\delta(G_{\mathrm{na}}, \mathcal{M}) > 0$, then at
  least one nontrivial $\U(1) \subset T_{\mathrm{rel}}$ is necessary
  (from EC: unresolved degeneracies violate admissibility completeness).
  \item If there exists a charge assignment $Y: \mathcal{M} \to \mathbb{Q}$
  such that the augmented label map
  $r_i \mapsto ([r_i]_{G_{\mathrm{na}}},\, Y(r_i))$
  is injective, then one abelian grading factor is sufficient
  (Lemma~\ref{lem:min_abelian}(ii)).
  \item Any additional independent abelian
  factor $Y'$ is inadmissible unless it changes the induced partition of
  $\mathcal{M}$; if $Y$ already separates all types, every further
  $\U(1)$ is redundant and excluded by CM
  (Lemma~2.22 of \TSII{}: each redundant generator wastes
  $\varepsilon^*$ capacity).
  \item Therefore, whenever $\delta > 0$ and one abelian grading
  separates all types, there exists exactly one physically admissible
  abelian gauge factor, unique up to overall normalization and change
  of basis in $T_{\mathrm{rel}}$.
\end{enumerate}
\emph{Scope.}  Uniqueness here is uniqueness \emph{relative to the
EC/CM criterion}: the claim is that once one abelian grading separates
all multiplet types, any further independent $\U(1)$ current adds a
generator (cost $\ge \varepsilon^*$, Lemma~2.22 of \TSII{})
while changing no gauge label assignment, and A2 eliminates it.  The
criterion by which a second $\U(1)$ could earn its keep is label
separation; a hypothetical second current that is ``physically
independent'' in some respect not expressible as a label distinction
is outside EC's vocabulary and is excluded by CM's cost accounting,
not by a separation argument.  This is a premise-level commitment of
the admissibility semantics (Definition~2.20 of \TSII{}),
named here so the reader can see exactly what carries the exclusion.
One further clause is owned in the same register: a second abelian
current distinct from hypercharge only \emph{dynamically} --- identical
kinematic labels on every multiplet, different couplings or breaking
pattern --- is excluded by the EC/CM vocabulary itself (kinematic
type-distinctness exhausted by counted labels, R-EC-inv), not by
finite capacity; dynamically distinguished currents are outside the
reading's scope by construction, and admitting them would change the
count.
\end{proposition}

\begin{proof}
\textit{Step 1: Necessity (at least one abelian factor).}\quad
If $\delta(G_{\mathrm{na}}, \mathcal{M}) > 0$, there exist
distinct multiplet types $r_i \ne r_j$ with
$[r_i]_{G_{\mathrm{na}}} = [r_j]_{G_{\mathrm{na}}}$.  By EC
(Lemma~2.24 of \TSII{}), these two physical
types must be distinguishable by the gauge structure.
Since $G_{\mathrm{na}}$ does not separate them, the distinction must
come from the remaining gauge freedom.  By
Lemma~\ref{lem:abelian_source}, the only continuous abelian gauge
freedom is $T_{\mathrm{rel}}$.  Hence at least one nontrivial
$\U(1) \subset T_{\mathrm{rel}}$ is necessary.  This gives (i).

\textit{Step 2: Sufficiency of one factor (explicit separation matrix).}\quad
We construct the separation matrix $S$ for
$\mathcal{M}_{\mathrm{SM}}$.  The non-abelian equivalence classes are:
\begin{center}\small
\begin{tabular}{@{}lll@{}}
\toprule
\textbf{Class} & \textbf{Non-abelian label} & \textbf{Members} \\
\midrule
$P_1$ & $(3, 2)$ & $\{Q\}$ (singleton) \\
$P_2$ & $(\bar{3}, 1)$ & $\{u^c, d^c\}$ \\
$P_3$ & $(1, 2)$ & $\{L\}$ (singleton) \\
$P_4$ & $(1, 1)$ & $\{e^c\}$ (singleton) \\
\bottomrule
\end{tabular}
\end{center}
The separation deficiency is $\delta = |P_2|(|P_2| - 1)/2 = 1$:
a single unresolved pair $(u^c, d^c)$.  The maximum class size is
$m_{\max} = 2$.  By Lemma~\ref{lem:min_abelian}, one $\U(1)$ factor
suffices iff it assigns distinct charges to $u^c$ and $d^c$.

\smallskip\noindent\fbox{\parbox{0.95\linewidth}{\textbf{Chirality/sign convention (fixed corpus-wide in v4.0).}\quad
Template entries are \emph{left-handed Weyl fields}: $u^c, d^c, e^c$
name the LH conjugates in representations $(\bar{\mathbf{3}},\mathbf{1})$,
$(\bar{\mathbf{3}},\mathbf{1})$, $(\mathbf{1},\mathbf{1})$, with
hypercharges $Y(u^c) = -2/3$, $Y(d^c) = +1/3$, $Y(e^c) = +1$.
Numerical charge \emph{values} in this document are quoted for the
right-handed partner fields, $Y_u := Y(u_R) = -Y(u^c) = +2/3$,
$Y_d := Y(d_R) = -1/3$, $Y_e := Y(e_R) = -1$, because those are the
conventional tabulated values.  Representation conjugation and
$\U(1)$ sign always flip together; v3.9 paired the $^c$ symbols with
the right-handed values, which mixed the two conventions --- the
notation below now uses $Y_u, Y_d, Y_e$ (right-handed values)
throughout, alongside the LH template reps.  The anomaly system is
invariant under the joint flip, so no computed quantity changes.}}
\smallskip

The hypercharge assignment $Y_u = 2/3 \ne -1/3 = Y_d$
resolves this degeneracy.  The full charge vector
$(Y_Q, Y_u, Y_d, Y_L, Y_e) = (1/6, 2/3, -1/3, -1/2, -1)$
is pairwise distinct across all five multiplet types.  The augmented
label map
\[
  r_i \mapsto \bigl([r_i]_{\SU(3) \times \SU(2)},\; Y(r_i)\bigr)
\]
is therefore injective on $\mathcal{M}_{\mathrm{SM}}$: no two
multiplets share both their non-abelian quantum numbers and their
hypercharge.  One $\U(1)$ suffices.  This gives (ii).

\textit{Step 3: No second abelian factor needed or admissible.}\quad
Suppose a second independent abelian factor with charge assignment
$Y': \mathcal{M} \to \mathbb{Q}$ is added.  The augmented map
becomes $r_i \mapsto ([r_i]_{G_{\mathrm{na}}}, Y(r_i), Y'(r_i))$.
Since $Y$ already makes this map injective (Step 2), the additional
charge $Y'$ does not refine the partition of $\mathcal{M}$: the
partition $\{\{Q\}, \{u^c\}, \{d^c\}, \{L\}, \{e^c\}\}$ is already
the finest possible (each class is a singleton).  By CM
(Lemma~2.24 of \TSII{}), the generator of $Y'$ consumes
$\varepsilon^* > 0$ capacity (Lemma~2.22 of \TSII{}) with zero
admissibility gain.  Under finite capacity (A1), this is inadmissible.
This gives (iii).

\textit{Step 4: Uniqueness up to normalization.}\quad
The surviving abelian factor is a one-dimensional subgroup of
$T_{\mathrm{rel}} \cong \mathrm{U}(1)^{m-1}$.  In the SM case,
$m = 3$ (blocks $M_3, M_2, M_1$), so $T_{\mathrm{rel}} \cong
\mathrm{U}(1)^2$.  The requirement that the charge assignment
resolves the $(u^c, d^c)$ degeneracy and satisfies all four anomaly
conditions (\S\ref{sec:hypercharge}) determines a \emph{unique
direction} in $T_{\mathrm{rel}}$, up to normalization
($Y \mapsto \lambda Y$, $\lambda \in \mathbb{Q}^\times$).  This
uniqueness is proved in Theorem~\ref{thm:hypercharge}: the anomaly
system has exactly two solutions $z = 1 \pm N_c$ (with $N_c = 3$:
$z = 4$ and $z = -2$), which correspond to the same direction
under the $u \leftrightarrow d$ relabeling symmetry.  Therefore
the abelian factor is unique.  This gives (iv).
\end{proof}

\begin{corollary}[Application to the derived Standard Model template]
\label{cor:unique_U1_SM}
For the derived multiplet set
$\mathcal{M}_{\mathrm{SM}} = \{Q,\, u^c,\, d^c,\, L,\, e^c\}$,
the non-abelian labels under $\SU(3) \times \SU(2)$ are
$Q \sim (3,2)$, $L \sim (1,2)$, $u^c \sim (\bar{3},1)$,
$d^c \sim (\bar{3},1)$, $e^c \sim (1,1)$.
Hence the non-abelian label map is not injective (it conflates $u^c$
with $d^c$): $\delta = 1$.  Hypercharge resolves this degeneracy with
pairwise distinct charges.  By Proposition~\ref{prop:unique_U1},
exactly one abelian gauge factor is necessary and sufficient.
\end{corollary}

The preceding chain of lemmas, proposition, and corollary assembles the
result from reusable components.  The reviewer asks for a single
self-contained theorem that takes the exact-sequence data as input and
outputs ``exactly one $\U(1)$'' without requiring the reader to trace
through the intermediate results.  We provide this.

\begin{theorem}[Unique abelian factor from exact-sequence data ---
SM specialization (Theorem~B)]%
\label{thm:unique_U1_consolidated}
\textbf{Hypotheses:}
\begin{enumerate}[label=(H\arabic*),nosep]
  \item \textbf{Admissibility algebra.}\quad
  $\mathcal{A} = M_3(\mathbb{C}) \oplus M_2(\mathbb{C}) \oplus
  \mathbb{C}$ (Wedderburn blocks $\{n_k\} = \{3, 2, 1\}$, from
  T2a + Theorem~R).
  \item \textbf{Exact sequence.}\quad
  $1 \to Z(\mathrm{U}(\mathcal{A})) \to \mathrm{U}(\mathcal{A})
  \xrightarrow{\pi} \mathrm{Inn}(\mathcal{A}) \to 1$, with
  $\mathrm{Inn}(\mathcal{A}) = \mathrm{PU}(3) \times \mathrm{PU}(2)$,
  $\ker(\pi) = \mathrm{U}(1)^3$,
  $T_{\mathrm{rel}} = \ker(\pi) / \mathrm{U}(1)_{\mathrm{diag}}
  \cong \mathrm{U}(1)^2$
  (Proposition~2.19 of \TSII{}).
  \item \textbf{Matter content.}\quad
  $\mathcal{M} = \{Q(3,2),\; u^c(\bar{3},1),\; d^c(\bar{3},1),\;
  L(1,2),\; e^c(1,1)\}$ (from $T_{\mathrm{field}}$,
  \S\ref{sec:fermion_scan}).  This hypothesis is the consumption site
  of the entry-typed inventory I-typing: the five listed multiplets
  are \emph{antecedently} distinct types --- $u^c \ne d^c$ prior to
  any abelian label is an input here, not a consequence of the labels
  --- and Definition~2.20 of \TSII{}'s copy clause is
  scoped to the full realized $\widetilde{G}$ precisely so that this
  inventory is not recounted under restriction to proper subgroups.
  \item \textbf{EC and CM.}\quad The gauge structure satisfies
  admissibility completeness (EC: the label map is injective) and
  capacity minimality (CM: no generator can be removed while
  preserving EC); EC holds under the kinematic type-inventory reading
  R-EC-inv with the entry-typed inventory I-typing (the
  Definition~2.20 of \TSII{} scoping premise and
  Lemma~2.24 of \TSII{}'s premise list); CM from PLEC's A2
  component over the A1 $+$ MD-defined admissible set
  (Lemma~2.24 of \TSII{}).
\end{enumerate}

\smallskip\noindent\textbf{Conclusion:}\quad Under (H1)--(H4) --- in
particular under EC as read through R-EC-inv with I-typing (H3--H4)
--- the physical gauge group
contains exactly one continuous abelian factor, which is a uniquely
determined one-dimensional subgroup of $T_{\mathrm{rel}} \cong
\mathrm{U}(1)^2$.  This factor is hypercharge $\U(1)_Y$, with charge
assignments $(Y_Q, Y_u, Y_d, Y_L, Y_e) =
(1/6, 2/3, -1/3, -1/2, -1)$ up to normalization (the compact-$\U(1)$
charge-lattice convention; front-matter ledger) and the
$u \leftrightarrow d$ relabeling.

\smallskip\noindent\textbf{Grades:}\quad the necessity leg ($\ge 1$
abelian factor) carries
[P\textsubscript{structural\_reading} $\mid$ R-EC-inv $+$ I-typing];
the sufficiency and uniqueness-of-direction legs are unconditional
given the anomaly system; the CM leg (no second factor) is [P]
(A2 over the A1 $+$ MD admissible set).
\end{theorem}

\begin{proof}
The proof proceeds in five steps: source exclusivity, necessity,
sufficiency, uniqueness of direction, and redundancy of any second
factor.

\textit{Step 1 (Source exclusivity): every abelian factor lives in
$T_{\mathrm{rel}}$.}\quad
By (H2), $\mathrm{Inn}(\mathcal{A}) = \mathrm{PU}(3) \times
\mathrm{PU}(2)$ is semisimple (both factors are semisimple for
$n \ge 2$).  A semisimple group has no continuous abelian
quotients: every connected abelian subgroup of $\mathrm{PU}(n)$ is
a maximal torus (the image of the Cartan subalgebra of
$\mathfrak{su}(n)$), which is part of the non-abelian structure.
The only continuous abelian subgroups of $\mathrm{U}(\mathcal{A})$
that project to the identity under $\pi$ are in $\ker(\pi) =
Z(\mathrm{U}(\mathcal{A}))$.  Quotienting by the unobservable
$\mathrm{U}(1)_{\mathrm{diag}}$ gives $T_{\mathrm{rel}} \cong
\mathrm{U}(1)^2$.  Any abelian gauge factor is a one-dimensional
subgroup of this torus.

\textit{Step 2 (Necessity): at least one $\U(1)$ is required.}\quad
By (H3), the non-abelian label map
$r_i \mapsto [r_i]_{\SU(3) \times \SU(2)}$ is:
\begin{center}\small
\begin{tabular}{@{}lccccc@{}}
\toprule
$r_i$ & $Q$ & $u^c$ & $d^c$ & $L$ & $e^c$ \\
\midrule
$[r_i]$ & $(3,2)$ & $(\bar{3},1)$ & $(\bar{3},1)$ & $(1,2)$ & $(1,1)$ \\
\bottomrule
\end{tabular}
\end{center}
The map sends both $u^c$ and $d^c$ to $(\bar{3}, 1)$.  Therefore
the label map is not injective: $\delta(\SU(3) \times \SU(2),
\mathcal{M}) = 1$.  By EC (H4), all five types must be
distinguished by the gauge structure.  Since the non-abelian
labels fail to separate $u^c$ from $d^c$, at least one abelian
factor from $T_{\mathrm{rel}}$ is required.

\textit{Step 3 (Sufficiency): one $\U(1)$ resolves all
degeneracies.}\quad
A one-dimensional subgroup $\U(1) \subset T_{\mathrm{rel}}$ assigns a
rational charge $Y_i$ to each $r_i$ (by Schur's lemma: central
elements act by scalars on irreducible representations; by the anomaly
system: the charges are rational, \S\ref{sec:hypercharge}).  The
augmented label map $r_i \mapsto ([r_i]_{\SU(3) \times \SU(2)},
Y_i)$ is injective iff $Y_u \ne Y_d$ (this is the only
unresolved pair from Step~2; the other three types already have
distinct non-abelian labels).  The anomaly system
(Theorem~\ref{thm:hypercharge}) determines:
\begin{align*}
  Y_u &= z \cdot Y_Q, \qquad Y_d = (2 - z) \cdot Y_Q,
\end{align*}
where $z$ satisfies $z^2 - 2z - (N_c^2 - 1) = 0$ with $N_c = 3$,
giving $z \in \{4, -2\}$.  For either root:
$Y_u - Y_d = (2z - 2) Y_Q \ne 0$ (since $z \ne 1$).
Therefore one $\U(1)$ suffices.

\textit{Step 4 (Uniqueness of direction): the anomaly system
admits a unique direction in $T_{\mathrm{rel}}$.}\quad
$T_{\mathrm{rel}} \cong \mathrm{U}(1)^2$ is a two-dimensional torus.
A one-dimensional subgroup is specified by a direction
$\mathbf{q} = (q_1, q_2) \in \mathbb{Z}^2 \setminus \{0\}$ (up to
positive rescaling).  The charge of multiplet $r_i$ is
$Y_i = q_1 \cdot c_1(r_i) + q_2 \cdot c_2(r_i)$, where
$c_1, c_2$ are the canonical $\mathrm{U}(1)$ charges from the
center $Z = \mathrm{U}(1)_3 \times \mathrm{U}(1)_2 \times
\mathrm{U}(1)_1$.  (In the canonical basis:
$c_1(r_i) = n_3$-ality of the $\SU(3)$ rep, $c_2(r_i)$ =
$\SU(2)$ isospin parity.)

The four anomaly conditions (\S\ref{sec:hypercharge}) are four
polynomial equations in $(q_1, q_2)$.  Three are linear; the
fourth is cubic.  The linear conditions determine the charge
ratios $Y_u/Y_Q$, $Y_d/Y_Q$, $Y_L/Y_Q$, $Y_e/Y_Q$
as functions of a single parameter $z = Y_u/Y_Q$.  The cubic
condition reduces to the quadratic $z^2 - 2z - 8 = 0$ with roots
$z = 4$ and $z = -2$.  These two roots correspond to the same
direction in $T_{\mathrm{rel}}$ (they are related by the
$u \leftrightarrow d$ relabeling symmetry of the fermion template).
Therefore the anomaly system admits a \emph{unique direction}
$\mathbf{q}$ in $T_{\mathrm{rel}}$, up to normalization and the
discrete relabeling.  (``Admits'' rather than ``selects'': no
mechanism ranges over candidate directions; the admissibility
conditions --- four anomaly equations plus EC, CM --- have a unique
non-trivial common solution up to the stated equivalences, and
that solution is the realised admissible hypercharge --- admissibility
space's argmin (A2) on the Y-assignment alternatives.)

\textit{Step 5 (No second factor): CM excludes any additional
$\U(1)$.}\quad
After Step~3, the augmented label map
$r_i \mapsto ([r_i]_{\SU(3) \times \SU(2)}, Y_i)$ is injective on
all five multiplet types (the full charge vector
$(1/6, 2/3, -1/3, -1/2, -1)$ is pairwise distinct).  The partition of
$\mathcal{M}$ into equivalence classes under this map is
$\{\{Q\}, \{u^c\}, \{d^c\}, \{L\}, \{e^c\}\}$ --- the discrete
partition (all singletons).

A second independent $\U(1)' \subset T_{\mathrm{rel}}$ would assign
a second charge $Y'_i$ to each multiplet.  The augmented map
$r_i \mapsto ([r_i], Y_i, Y'_i)$ induces the same discrete
partition (singletons cannot be further refined).  Therefore $Y'$
provides zero additional admissibility information.  By CM (H4), the
generator of $\U(1)'$ occupies one admissibility channel
(Lemma~2.22 of \TSII{}) at cost $\varepsilon^* > 0$
($L_{\varepsilon^*}$) with zero admissibility gain.  Under finite
capacity (A1), this is inadmissible.

\textit{Conclusion.}\quad Steps~1--5 establish: exactly one
continuous abelian factor exists in the physical gauge group; it is a
one-dimensional subgroup of $T_{\mathrm{rel}} \cong \mathrm{U}(1)^2$;
the anomaly system selects a unique direction; and no second factor is
admissible.  The physical gauge group is therefore
$[\SU(3) \times \SU(2) \times \U(1)_Y] / \mathbb{Z}_6$
(Proposition~2.26 of \TSII{}).
\end{proof}

\begin{tcolorbox}[apfbox, title={Unique $\U(1)$ proof in conventional language}]
\footnotesize
Starting from $\mathcal A=M_3(\mathbb C)\oplus M_2(\mathbb C)\oplus\mathbb C$, the nonabelian inner automorphism group contributes $\mathrm{PU}(3)\times\mathrm{PU}(2)$ and no central $\U(1)$.  The unitary-center quotient leaves a two-dimensional relative-phase torus.  The derived matter template contains two multiplets with the same nonabelian label, $u^c$ and $d^c$, so admissibility completeness requires at least one relative phase.  The standard anomaly equations on the scanned template reduce that torus to one rational projective ray, hypercharge.  Once this ray is included, all five multiplets are separated; any second independent ray either violates the anomaly equations on the 45-Weyl template or separates no additional minimal-template distinction.  MD assigns that extra generator positive cost, so A2 removes it.  This is the whole no-redundancy proof, stated without APF shorthand.
\end{tcolorbox}

\noindent\textit{Why the theorem is not heuristic.}\quad
The selection of one $\U(1)$ from the $(K-1)$-torus involves three
independently verifiable steps, none of which is an aesthetic or
heuristic choice:
\begin{itemize}[nosep]
  \item \textit{Necessity} (Step~2) is a counting argument:
  $\delta > 0$ combined with EC forces at least one abelian factor.
  This is a finite combinatorial check on the five multiplets.
  \item \textit{Uniqueness of direction} (Step~4) is a
  number-theoretic result: the anomaly system is a determined system
  of polynomial equations with a unique projective solution.  This is
  verified by explicit computation (the discriminant of
  $z^2 - 2z - 8 = 0$ is $36$; the roots are $4$ and $-2$; they are
  related by the relabeling symmetry).
  \item \textit{Redundancy exclusion} (Step~5) is a capacity
  argument: the partition is already maximally fine (all singletons),
  so no new information can be added.  This is a logical tautology
  combined with the generator--channel bridge (each generator costs
  $\varepsilon^* > 0$).
\end{itemize}
The only conditional input is the matter content
$\mathcal{M} = \mathcal{M}_{\mathrm{SM}}$ (H3), which is derived in
$T_{\mathrm{field}}$ (\S\ref{sec:fermion_scan}).  Given that input,
the conclusion is a theorem, not a selection.


\smallskip\noindent\textit{Edge case: $\U(1)_{B-L}$ with sterile
$\nu_R$.}\quad
A frequently raised alternative is the anomaly-free combination
$\U(1)_{B-L}$ (baryon minus lepton number), which becomes anomaly-free
when the multiplet set is extended by three right-handed neutrinos
$\nu_R(\mathbf{1},\mathbf{1})$.  In this enlarged template
$\mathcal{M}' = \mathcal{M}_{\mathrm{SM}} \cup \{\nu_R\}$, both
$\U(1)_Y$ and $\U(1)_{B-L}$ are anomaly-free, and the question is
whether two independent abelian factors are admissible.

The APF excludes this scenario at two stages:
\begin{enumerate}[label=(\alph*),nosep]
  \item \textbf{$T_{\mathrm{field}}$ excludes $\nu_R$ from the minimal
  template.}\quad The fermion scan (\S\ref{sec:fermion_scan}) derives
  $\mathcal{M}_{\mathrm{SM}}$ as the unique minimum-DOF survivor at
  45 Weyl fermions.  The template
  $\mathcal{M}_{\mathrm{SM}} \cup \{\nu_R\}$ has $48$ DOF and survives
  all anomaly filters, but is non-minimal: it is the second-lightest
  post-CPT survivor (template \#3 in the waterfall table).  Since the
  multiplet set is $\mathcal{M}_{\mathrm{SM}}$, not $\mathcal{M}'$,
  the $B{-}L$ direction is not anomaly-free for the \emph{derived}
  content.  (Indeed, $B{-}L$ has a mixed gravitational anomaly
  $[\mathrm{grav}]^2[\U(1)_{B-L}] = 3 \ne 0$ unless $\nu_R$ is
  present.)
  \item \textbf{Capacity minimality excludes the second $\U(1)$ even
  if $\nu_R$ were present.}\quad Suppose counterfactually that
  $\mathcal{M}' = \mathcal{M}_{\mathrm{SM}} \cup \{\nu_R\}$ were the
  derived template.  Does Proposition~\ref{prop:unique_U1} still
  exclude a second $\U(1)$?  Yes: in $\mathcal{M}'$, the non-abelian
  labels conflate $u^c / d^c$ (both $(\bar{3},1)$) and
  $e^c / \nu_R$ (both $(1,1)$).  A single abelian grading with four
  distinct charges resolves all degeneracies: one $\U(1)$ suffices.
  $\U(1)_{B-L}$ is then a \emph{linear combination} of this grading
  and the non-abelian quantum numbers---not an independent gauge
  direction.  A second independent $\U(1)$ would carry capacity cost
  (one additional admissibility channel) with no additional
  multiplet-separation power, violating CM.
  \item \textbf{Observational cross-check (bridge-conditional,
  exited).}\quad Under the cosmological identification (BA6,
  \S\ref{sec:bridge_axioms}) the $\nu_R$-extended count
  $C_{\mathrm{total}} = 64$ reads $\Omega_\Lambda = 42/64$ against
  the 45-template's $42/61$; the comparison with the Planck
  measurement lives at the corpus homes (Papers~8 and~41, with the
  exited cosmology tail).  The core exclusion above is by minimality
  and consumes no bridge axiom.
\end{enumerate}
In summary, the ``exactly one $\U(1)$'' conclusion is robust against
the $B{-}L$ edge case: the extended multiplet set is non-minimal
($T_{\mathrm{field}}$), $B{-}L$ is anomalous for the minimal template,
and even if $\nu_R$ were included, $B{-}L$ would be a redundant linear
combination rather than an independent gauge factor.

\smallskip\noindent\textit{String-theoretic extra $\U(1)$'s and
St\"uckelberg masses.}\quad
Heterotic string compactifications and F-theory models generically
produce multiple abelian gauge factors beyond hypercharge.  In these
constructions, unwanted $\U(1)$'s acquire large masses via the
St\"uckelberg mechanism: the $\U(1)$ gauge boson ``eats'' a closed-string
axion, becoming massive at or near the compactification scale, while
the associated global selection rule (e.g., a $\mathbb{Z}_N$ discrete
symmetry) may persist at low energies.

In the APF, this scenario has a precise capacity interpretation:
\begin{enumerate}[label=(\alph*),nosep]
  \item A St\"uckelberg-massive $\U(1)$ is a gauge boson that no longer
  propagates as a long-range force.  In the admissibility language, the
  corresponding channel is ``saturated'': its admissibility cost is paid
  at the compactification scale, and at low energies it is absorbed into
  the vacuum sector ($L_{\mathrm{vac}}$; exited to Paper~45 at the v4.0 core split, BA9 in \S\ref{sec:bridge_axioms}).
  The channel is counted in $C_{\mathrm{total}}$ (as part of the UV
  structure) but does not contribute an independent low-energy gauge
  coupling.
  \item Capacity minimality (CM) does not merely select the gauge group
  with the fewest generators; it selects the group with the fewest
  \emph{active admissibility channels} at the IR saturation scale.
  A St\"uckelberg-massive $\U(1)$ is inactive at low energies and
  therefore excluded from the low-energy gauge group by the same
  principle that excludes $\nu_R$ from the low-energy fermion count:
  it is above the IR capacity ledger.
  \item The APF is therefore selecting the \emph{phase} in which only
  one $\U(1)$ combination remains both anomaly-free and massless at
  low energies.  This is precisely the phase that heterotic
  phenomenology identifies as the ``Standard Model embedding'': among
  the multiple $\U(1)$'s of the string construction, one linear
  combination is identified with hypercharge and remains massless, while
  the others acquire St\"uckelberg masses and decouple.  The APF's
  capacity-minimality principle provides the \emph{selection criterion}
  for this phase: the massless $\U(1)$ is the one that resolves all
  low-energy multiplet degeneracies at minimum capacity cost.
\end{enumerate}

\smallskip\noindent\textit{Kinetic mixing.}\quad
Kinetic mixing --- a gauge-invariant term
$\frac{\chi}{2} F_{\mu\nu} F'^{\mu\nu}$ coupling two abelian field
strengths --- needs a second propagating $\U(1)$, and
Proposition~\ref{prop:unique_U1} leaves none to mix with.  The
concern is moot for the low-energy gauge structure derived here, and
we say no more about it in the classification core: the capacity
accounting of hypothetical mixed pairs is program-register material
and belongs with the foundational program, not among the closure
theorems.

\begin{tcolorbox}[apfbox, title={Roadmap: from $\mathrm{PU}$ automorphisms to the physical SM group}]
\footnotesize
The gauge derivation has four logically distinct layers.
\begin{enumerate}[nosep,leftmargin=1.8em]
  \item \textbf{Algebra layer:} inner automorphisms of the finite
  semisimple admissibility algebra give $\prod_k\mathrm{PU}(n_k)$.
  This is the adjoint action on observables.
  \item \textbf{Field layer:} matter multiplets are vectors/modules in
  $\mathcal H_\omega$, so projective $\mathrm{PU}(n)$ actions must be
  lifted to linear $\SU(n)$ actions whenever fundamental matter exists.
  \item \textbf{Abelian layer:} $\U(1)_Y$ does not come from
  $\mathrm{Inn}(\mathcal A)$; it comes from the relative-phase torus
  $Z(U(\mathcal A))/\U(1)_{\mathrm{diag}}$, pruned by the no-redundancy
  theorem to a single anomaly-compatible direction.
  \item \textbf{Global layer:} central elements acting trivially on all
  derived matter multiplets are quotiented.  For the SM hypercharge ray
  this kernel is $\mathbb Z_6$.
\end{enumerate}
This box is the local answer to the common ambiguity ``why
$\mathrm{Inn}(\mathcal A)=\prod\mathrm{PU}(n)$ but the physical group is
$[\SU(3)\times\SU(2)\times\U(1)_Y]/\mathbb Z_6$.''  The first group acts
faithfully on observables; the second acts faithfully on the derived
matter content.
\end{tcolorbox}


\begin{theorem}[No-redundancy theorem for extra abelian factors]
\label{thm:no_extra_U1_v34}
In the minimal APF gauge/matter theorem, the only continuous abelian gauge
factor is the hypercharge ray $\U(1)_Y$.  Any additional continuous abelian
factor is excluded by one of three mutually exclusive mechanisms:
\begin{enumerate}[nosep,label=(\roman*)]
  \item it is a redundant central direction whose charges are an affine
  function of hypercharge on the derived template and therefore has strictly
  positive cost but no additional separating power;
  \item it is anomaly-incompatible on the 45-Weyl minimal template; or
  \item it separates distinctions not present in the minimal template, such as
  sterile-neutrino number or generation-resolved lepton number, and therefore
  belongs to a different enlarged admissibility problem.
\end{enumerate}
Consequently the theorem proves uniqueness of a single rational ray, not merely
existence of one convenient $\U(1)$ basis.
\end{theorem}

\begin{proof}
By Lemma~\ref{lem:abelian_source}, every continuous abelian gauge factor must
come from the relative-phase torus
$T_{\mathrm{rel}}=Z(\mathrm U(\mathcal A))/\mathrm U(1)_{\mathrm{diag}}$.
Let $q$ be any second abelian charge direction linearly independent of the
hypercharge assignment $Y$ on the minimal matter template.

If $q$ is constant on each non-abelian degeneracy class already separated by
$Y$, then replacing $(Y,q)$ by $Y$ preserves admissibility completeness.  Since
MD assigns a positive floor to each independent generator, A2 selects the
one-generator realization.  This is the capacity form of no redundancy.

If $q$ acts nontrivially on the minimal 45-Weyl template, it must also satisfy
the four anomaly equations.  The anomaly system in
Theorem~\ref{thm:hypercharge}, with the scanned representation content held
fixed and no sterile singlet added, has a one-dimensional solution space; after
minimal integral normalization this solution is precisely the hypercharge ray.
Thus an independent $q$ is not anomaly-compatible on that template.

The remaining familiar anomaly-free directions evade this conclusion only by
changing the template.  A $B-L$ direction becomes anomaly-compatible when a
right-handed neutrino or equivalent sterile record is included.  Family
non-universal directions such as $L_\mu-L_\tau$ require generation labels to
be independently gauge-addressable distinctions.  Both moves add new primitive
records, so they are outside the minimal theorem rather than counterexamples to
it.  The APF claim is therefore conditional and sharp: within the finite
minimal template fixed here, exactly one abelian ray survives; enlarged
record vocabularies define enlarged theorems with additional cost.

The normalization is fixed as follows.  The anomaly equations determine a ray
in $T_{\mathrm{rel}}$.  Minimal charge quantization fixes the primitive
integer representative by choosing $Y_{\mathrm{int}}=6Y$ with
$\gcd\{|Y_{\mathrm{int}}(r)|:r\in\mathcal M\}=1$.  Hence the continuous
freedom is only the conventional gauge-coupling normalization, not an
additional abelian factor.
\end{proof}

\begin{remark}[Right-handed neutrinos and enlarged abelian problems]
\label{rem:nuR_extended_v35}
If a sterile right-handed neutrino $\nu_R$ is included as a primitive matter record, the theorem's input template changes from the minimal 45-Weyl problem to a 48-Weyl problem.  In that enlarged problem the standard anomaly equations admit an additional $B-L$ direction.  APF therefore does not claim that $B-L$ is algebraically impossible.  The claim is narrower: $B-L$ is not forced by the minimal gauge/matter template and it adds an independently admissible sterile or lepton-number distinction at positive MD cost.  Hypercharge quantization on the minimal template remains a unique rational ray; in the $\nu_R$-extended template the abelian solution space is larger, and A2 must be re-run on the enlarged admissible set rather than imported from the 45-Weyl theorem.
\end{remark}


\newpage

\section{Scope of the three-generation input}
\label{sec:three_generation_scope}

This supplement proves the one-generation gauge/matter template and then
uses the three-copy capacity ledger when reporting the 45-Weyl-fermion result.
The necessity of exactly three generations is not re-derived here from first
principles.  It is imported from the capacity-counting branch of the program,
and the import has named theorems, each machine-verified.  The generation
count is over-determined by three independent capacity--beta identities, each
a linear equation with the unique solution $n = 3$: $6|b_3| = C_{\mathrm{vacuum}}$
($66 - 8n = 9n + 15$), $6|b_2| = C_{\mathrm{matter}}$ ($43 - 8n = 6n + 1$)
\coderef{check\_L\_beta\_capacity}{generations.py}, and the abelian member
$6|b_Y| = d_{\mathrm{eff}} - C_{\mathrm{total}}$ ($(40/3)n + 1 = 9n + 14$)
\coderef{check\_T\_gauge\_beta\_capacity\_tiling\_abelian}{gauge\_beta\_capacity\_tiling.py}.
Because these are equations rather than one-sided bounds, they carry both
directions: $n = 1, 2$, and $4$ fail all three.  Independently, the
max-saturated selection at the electroweak wall kills under-replicated rivals
outright ($n \in \{1,2\}$ subsaturated, $n \ge 4$ over the ceiling)
\coderef{check\_R\_Ngen\_neq\_3\_killed}{killed\_rivals.py}.
For the present supplement the statement is therefore conditional but explicit:
if the capacity ledger supplies three copies of the one-generation template,
then the fermion count is $3\times15=45$ Weyl fermions and the anomaly/hypercharge
closure repeats generationwise.  If one treats the generation number as a free
external input, the gauge/hypercharge derivation still holds per generation,
while the 45-count and cosmological denominator become bridge claims rather
than outputs of this supplement alone.

\section{Neutrino content of the minimum template}
\label{sec:neutrinos}
% ══════════════════════════════════════════════════════════════

The supported statement of this section is the scan's own: \emph{the
minimum template in the declared scan contains no light sterile
singlet.}  The 45-fermion survivor carries no right-handed neutrino
$\nu_R$; the nearest survivor that does (template \#3 of the
waterfall listing, at 48 DOF) sits strictly above the minimum and is
excluded at the minimality cut.  That is a classification fact about
the declared scan space, and it is all this section claims.

What the absence of $\nu_R$ means for neutrino \emph{masses} ---
the dimension-5 Weinberg-operator mechanism, the Majorana character
of the resulting masses, and the neutrinoless-double-beta-decay
observable --- is physics downstream of the classification, stated at
prediction register at the corpus homes (Papers~8 and~41, with the
exited cosmology tail).  So is the bridge-conditional observational
argument against a light $\nu_R$ (the $C_{\mathrm{total}} = 64$
reading under BA6).  Neither is consumed by any theorem of this
document.

\smallskip\noindent\textit{Why $\nu_R$ is not counted.}\quad A heavy
UV completion of neutrino mass (a Majorana state at a high scale
$\Lambda$) is above the IR saturation scale and is not a distinct
admissibility channel at low energies.  The capacity count
($C_{\mathrm{total}} = 61$) reflects the \emph{low-energy}
admissibility ledger; states above the saturation scale contribute to
the UV completion but not to the structural partition.


\begin{theorem}[$\nu_R$ extension theorem]
\label{thm:nuR_extension}
If a right-handed neutrino $\nu_R(\mathbf 1,\mathbf 1)$ is added as a
primitive matter distinction, the anomaly equations enlarge from the
45-Weyl minimal template to a 48-Weyl extended template.  In that extended
template, $B-L$ becomes anomaly-free, but it is not an additional abelian
factor in the minimal APF sense unless the theory also pays for an
independent sterile/generation-resolving admissibility record.  A single
abelian grading still separates all non-abelian degeneracies; a second
massless $\U(1)$ is capacity-redundant unless it carries new primitive
record content.  Thus $\nu_R$ is admissible as a non-minimal extension, but
including it changes the theorem from ``the unique minimum is 45 Weyl DOF
with one $\U(1)_Y$'' to ``a 48-Weyl extension admits additional anomaly-free
linear combinations only when extra record distinctions are supplied.''
\end{theorem}

\begin{proof}
Adding $\nu_R$ cancels the mixed gravitational anomaly of $B-L$ and opens a
larger abelian solution space in the standard anomaly equations.  However,
APF counts not only anomaly freedom but admissibility gain.  Hypercharge already
separates the non-abelian degeneracies $u^c/d^c$ and $e^c/\nu_R$ in the
extended set.  A second abelian factor therefore distinguishes no additional
minimal-template multiplet unless one declares the sterile or family label to
be a new primitive record.  Such a declaration is consistent, but it is an
extension of the carrier and raises the ledger from the 45-Weyl theorem to a
48-Weyl theorem.  The minimal result is therefore unaffected, while the
extension is explicitly classified.
\end{proof}

\smallskip\noindent\textit{Where the physics readings live.}\quad
The Majorana/$0\nu\beta\beta$ prediction register and the
bridge-conditional Planck comparison of the $\nu_R$-extended count
($C_{\mathrm{total}} = 64$ under BA6) are stated at the corpus homes
(Papers~8 and~41, with the exited cosmology tail), where they carry
their bridge conditioning.  What stays here is the scan fact: no
light sterile singlet in the minimum template.
\coderef{check\_T\_field}{gauge.py}


\newpage

% ╔═══════════════════════════════════════════════════════════════╗
% ║              PART III: CAPACITY COUNTING AND COSMOLOGY        ║
% ╚═══════════════════════════════════════════════════════════════╝

\part{Capacity Counting}
\label{part:cosmo}


% ══════════════════════════════════════════════════════════════
\section{Capacity counting and the BRST question}
\label{sec:capacity_count}
% ══════════════════════════════════════════════════════════════

% ── §9: Capacity counting and the BRST question ──

\subsection{Operational definition of capacity types}

A \emph{capacity type} is a structurally distinct admissibility role
in the low-energy ledger.  The count is:
\begin{equation}
  C_{\mathrm{total}} = \underbrace{45}_{\text{fermion types}}
    + \underbrace{4}_{\text{Higgs real components}}
    + \underbrace{12}_{\text{gauge generators}} = 61.
\end{equation}

\noindent The count is stated under an explicit hypothesis: the
channel-identification convention (BA5, \S\ref{sec:bridge_axioms};
front-matter cost-uniqueness ledger row) that each Weyl fermion
component, each real Higgs component, and each gauge generator
occupies exactly one capacity slot of equal weight --- the channel
unit fixed by the cost functional ($L_{\mathrm{Cauchy}}$:
$F(n) = n$).  The 61 is a theorem in that unit; sector-dependent slot
weights would dissolve it (BA5's stated failure mode).

\noindent\textit{Fermion types} (45): the 15 Weyl fermions per
generation ($Q = 6$ colored-doublet components $+$ $u^c = 3$
colored-singlet $+$ $d^c = 3$ $+$ $L = 2$ weak-doublet components
$+$ $e^c = 1$) times 3 generations.  Color and weak-isospin
multiplicity \emph{are} counted here as independent matter-sector
admissibility roles, because each distinct quantum-number configuration
is a separately admissible distinction in the capacity ledger.  The
arithmetic is $(6+3+3+2+1) \times 3 = 15 \times 3 = 45$, consistent
with \texttt{check\_L\_count} in \texttt{gauge.py} (\texttt{per\_gen =
6+3+3+2+1 = 15}).  Internal-structure gauge channels (the 27
fermion-internal channels counted in $L_{\mathrm{vac}}$, Paper~2
\S7.1) are a distinct bookkeeping allocated to the vacuum sector;
they are not the matter-sector fermion types counted here.  This
resolves an earlier prose formulation that read ``the 5 multiplet
types'' and counted $5 \times 3 = 15$ rather than 45; the correct
count is $(6+3+3+2+1) \times 3 = 45$.

\noindent\textit{Higgs components} (4): the 4 real degrees of freedom
of the Higgs doublet.  After SSB, 3 become Goldstone bosons (absorbed
into $W^\pm, Z$) and 1 is the physical Higgs.  All 4 are counted
because the capacity ledger operates at the structural level, before
SSB selects a vacuum.

\noindent\textit{Gauge generators} (12): $\dim(\SU(3)) + \dim(\SU(2))
+ \dim(\U(1)) = 8 + 3 + 1$.  Each generator is a structurally distinct
admissibility channel.

\subsection{Why gauge constraints are not subtracted (the BRST question)}
\label{sec:BRST}

In perturbative quantum field theory, the physical Hilbert space is the
BRST cohomology: gauge constraints subtract unphysical polarizations,
reducing the propagating degrees of freedom.  A natural question is: why
does the capacity count not perform a similar subtraction?

The answer is that capacity types count \emph{admissibility roles}, not
\emph{propagating degrees of freedom}.  The distinction is:

\smallskip\noindent\textit{Propagating DOF} (QFT accounting): count
on-shell polarization states.  A massless gauge boson has 2 polarizations;
BRST subtracts the longitudinal and timelike modes.  This gives
$2 \times 12 = 24$ propagating gauge DOF, not 12.

\smallskip\noindent\textit{Enforcement roles} (capacity accounting):
count structurally distinct channels in the admissibility ledger.  Each
gauge generator $T^a \in \mathfrak{g}$ is one independent admissibility
channel: it acts as a derivation on the algebra of observables, routing
admissibility between matter fields.  The number of generators is
$\dim(\mathfrak{g}) = 12$, regardless of how many polarization states
each generator supports.

The 12 gauge generators are assigned to the \emph{vacuum} sector
(Q1 = 0 in $P_{\mathrm{exhaust}}$) precisely because they are
admissibility \emph{infrastructure}, not admissibility \emph{consumers}.
A gauge generator cannot route through itself; it has no upstream
pathway to enter as a matter excitation.  BRST is irrelevant to this
accounting: it subtracts polarizations from the propagator, which is a
kinematic object; the capacity count is a structural object.

\begin{tcolorbox}[apfbox, title={Toy $\U(1)$ gauge model: constraint subtraction vs. channel accounting}]
\footnotesize\label{ex:brst_toy_v35}
A free Maxwell field has four components $A_\mu$.  Gauge fixing and BRST cohomology remove the longitudinal and timelike excitations, leaving two propagating photon polarizations.  The APF ledger does not count four minus two; it counts one gauge-routing generator, because there is one independent local redundancy direction acting on charged matter by $\psi\mapsto e^{i\alpha}\psi$.  For a nonabelian group the same rule gives one admissibility channel per Lie-algebra generator: $\dim\mathfrak g$.  Ghosts and gauge-fixing fields are bookkeeping devices for quantizing the quotient; they do not supply new charged matter distinctions and do not erase the existence of the generator that must be enforced.  This is why $\SU(3)$ contributes $8$ structural channels, $\SU(2)$ contributes $3$, and hypercharge contributes $1$, even though the propagating polarization count is different.
\end{tcolorbox}
\coderef{check\_L\_count}{gauge.py}

\subsection{$I2$ (gauge-cosmological identity): the $K = 61$ shared integer --- exit note}\label{sec:I2_at_L_count}

What this supplement owns is the integer $C_{\mathrm{total}} = 61$
and its structural partition $61 = 45 + 4 + 12$
($L_{\mathrm{count}}$, above).  The identity $I2$ --- that the same
integer reappears as the common denominator of the cosmological
projection, with the residual partition $3 + 16 + 42$ --- is stated,
proved against the ACC record, and compared with observation at its
corpus home (Paper~8, with the strictly finer subspace refinement in
Paper~6), where it carries its bridge conditioning (BA5--BA6,
\S\ref{sec:bridge_axioms}).  A reader who rejects a bridge axiom
loses exactly those readings; the integer and the structural
partition proved here are untouched.
\coderef{check\_I2\_gauge\_cosmological, check\_T\_ACC\_unification}{unification.py}


\smallskip\noindent\textit{Where the cosmological readings live.}\quad
What this supplement owns is the integer $C_{\mathrm{total}} = 61$,
its partition, and the identity $I2$.  The substrate-position-axis
architectural reading (the (61,\,102) per-slot lattice), the
sensitivity table, and the comparison of the bridge readings
$\Omega_\Lambda = 42/61$ and kin with the Planck measurements are
imported by Papers~8 and~41 at bridge grade (BA5--BA7,
\S\ref{sec:bridge_axioms}); that material is staged at those corpus
homes and carries the bridge conditioning with it.  A reader who
rejects a bridge axiom loses exactly those readings and nothing here.


\newpage

\appendix


\section{Comparison with ``SM from principles'' programs}
\label{app:literature}
% ══════════════════════════════════════════════════════════════

The comparisons kept in this document are the ones that bear on the
classification claim; the reconstruction-program and
operational-axiom positioning of the foundational chain lives in
\TSII's comparison appendix.

\subsection{SM-from-principles programs}
\label{sec:sm_programs}

A separate tradition attempts to derive the Standard Model within an
assumed quantum field theory framework.

\smallskip\noindent\textbf{Anomaly-based selection
(Minahan--Ramond--Warner, 1990).}\quad
Given $\SU(3) \times \SU(2) \times \U(1)$ and the requirement of
anomaly cancellation, the SM fermion content is strongly constrained
but not unique without additional assumptions (e.g., minimality,
family structure).  The APF's $T_{\mathrm{field}}$ performs a similar
scan, and the gauge template it consumes (input C2 of
Theorem~\ref{thm:conditional_classification}) is not an unexamined
assumption: it is the output of the graded APF chain (Theorem~R,
\TSII{}) together with this document's $L_{\mathrm{gauge,unique}}$
ranking, at the grades stated there.

\smallskip\noindent\textbf{Grand unification (Georgi--Glashow, 1974).}\quad
Embeds $\SU(3) \times \SU(2) \times \U(1)$ into $\SU(5)$.  Explains
charge quantization and predicts proton decay.  The APF derives the
product structure as capacity-minimal
(Proposition~\ref{prop:no_simple}: $\SU(5)$ costs $2\times$) and
obtains the charge-ratio quantization --- the unique rational ray of
hypercharges --- from anomaly cancellation
(Theorem~\ref{thm:hypercharge}) without grand unification; the
absolute charge unit is the compact-$\U(1)$/character-lattice
convention, not an anomaly output (\S\ref{sec:one_u1_closure}).

\smallskip\noindent\textbf{Connes NCG (1996).}\quad
The noncommutative geometry program derives the SM Lagrangian from a
spectral triple.  The input is a specific finite-dimensional algebra
(chosen to reproduce the SM).  The APF obtains the algebra from the foundational chain at its banked
grades ($T_{\mathrm{alg}}$, derived given the IJC occupancy; the
program document is \TSII{}), rather than postulating it; the two
programs converge on the same algebraic structure but arrive from
opposite directions.

\smallskip
\begin{center}
\small
\renewcommand{\arraystretch}{1.12}
\begin{tabular}{@{}p{3.2cm}p{4.5cm}p{5.0cm}@{}}
\toprule
\textbf{Question} & \textbf{NCG route} & \textbf{APF route} \\
\midrule
Finite algebra & Chosen as part of the spectral triple data. & Derived from finite admissibility, nonclosure, and Wedderburn decomposition. \\
Gauge fields & Inner fluctuations / inner automorphisms of the algebra. & Internal continuation-profile-preserving automorphisms, identified with $\mathrm{Inn}(\mathcal A)$. \\
$\U(1)$ source & Encoded in the chosen finite algebra and unimodularity conditions. & Relative-phase torus of $U(\mathcal A)$, pruned to one direction by capacity minimality and anomaly compatibility. \\
Global quotient & Determined by representation of the chosen algebra on the Hilbert space. & Derived from the kernel acting trivially on the scanned matter multiplets. \\
Anomaly handling & Standard consistency requirement on the represented fermions. & Standard anomaly algebra applied after APF fixes the template and excludes exotic survivors. \\
Novelty claim & Geometric/spectral encoding of known SM data. & Derivation claim for why this finite algebra/template is the minimum-cost admissible one. \\
\bottomrule
\end{tabular}
\end{center}


\begin{center}\footnotesize
\renewcommand{\arraystretch}{1.12}
\begin{tabular}{@{}p{3.0cm}p{5.2cm}p{5.2cm}@{}}
\toprule
\textbf{Subtle point} & \textbf{NCG mechanism} & \textbf{APF mechanism} \\
\midrule
Single $\U(1)$ & Unimodularity and central-extension choices restrict the unitary group of the finite spectral triple. & Relative-phase torus is enumerated first, then no-redundancy/A2 and anomaly compatibility select one primitive rational ray. \\
Hypercharge normalization & Fixed by representation of the finite algebra on the Hilbert space and unimodularity convention. & Anomaly equations determine a ray; primitive integer normalization $Y_{\rm int}=6Y$ fixes the charge lattice, while gauge coupling remains conventional. \\
Global quotient & Kernel of the represented finite-algebra unitary group on the fermion Hilbert space. & Kernel of the lifted $\SU(3)\times\SU(2)\times\U(1)$ action on the scanned matter modules; Proposition~2.26 of \TSII{} verifies it multiplet by multiplet. \\
Higgs/inner fluctuation & Higgs arises as an inner fluctuation of the finite Dirac operator. & Higgs appears as a capacity-minimal scalar bridge for Yukawa/admissible mass terms; its dynamical potential is downstream, not used to derive the gauge template here. \\
Finite algebra input & Chosen as spectral-triple data subject to NCG axioms. & Produced as the minimum-cost finite admissibility carrier, with NCG-like algebraic facts applied after selection. \\
\bottomrule
\end{tabular}
\end{center}

\smallskip\noindent\textit{Quaternionic branch and weak $\SU(2)$.}\quad
NCG often realizes the weak factor through an almost-commutative algebra
containing an $\mathbb H$ summand.  APF does not use that route.  The weak
factor appears as the simply connected lift of the $\mathrm{PU}(2)$ internal
automorphism factor acting on complex doublet matter modules.  Quaternionic
simple components are not part of the Paper~2 finite complex-sector theorem:
within the imported Paper~5 field-selection gate, they either collapse the
chirality/charged-module distinction needed for the scanned template or add an
anti-linear/reality structure not selected by the no-added-record criterion.
Thus APF replaces the quaternionic ingredient by complex doublet modules plus
the cover $\mathrm{PU}(2)\leftarrow\SU(2)$; the full real/complex/quaternionic
exclusion is imported from the field-selection supplement.

\smallskip\noindent\textit{Lattice chiral-gauge and symmetric-mass-generation routes.}\quad
Lattice chiral-gauge programs, mirror-fermion decoupling, symmetric mass
generation, and related RKD/disentangler constructions address a different
problem: whether a given anomaly-free chiral spectrum can be realized by a
local microscopic regulator without unwanted light mirror sectors.  APF uses
anomaly cancellation as a necessary filter, not as a sufficient proof of
lattice realizability.  The finite-capacity derivation therefore should be
read as an upstream selection theorem for the candidate low-energy spectrum;
a lattice construction would be an independent downstream regulator witness.
If an anomaly-free APF survivor failed all local-regulator tests, that would
be evidence against its physical admissibility, not a contradiction of the
anomaly algebra.

In the other direction, SMG/mirror-fermion successes support the APF split
between algebraic consistency and microscopic realizability: anomaly-free
spectra can sometimes be gapped without breaking symmetry, while anomalous
spectra cannot.  APF's additional claim is that the minimal finite-capacity
ledger selects the SM survivor before choosing a regulator.  The two questions
are therefore complementary: APF asks which template is minimum-cost; lattice
work asks whether the template admits a local UV completion.


\smallskip\noindent\textbf{Doplicher--Roberts reconstruction and
Pinzari's braided extensions.}\quad
The classical DR theorem (1989) reconstructs a compact group $G$ from
its representation category (a symmetric $C^*$-tensor category with
conjugates).  In the APF, T3 (Paper~1 Supplement) uses DR to recover
the gauge group from the admissibility algebra's representation
category; the gauge identification theorem
(Theorem~2.11 of \TSII{}) then identifies which automorphisms
constitute the gauge group.


\begin{proposition}[Finite APF admissibility net as a DR precursor]
\label{prop:apf_dr_mapping_v34}
A finite APF admissibility net determines the finite precursor of a
Doplicher--Roberts reconstruction datum.  The assignment sends an interface
$\Gamma$ to its field algebra $\mathcal A(\Gamma)$, sends an inclusion of
interfaces to the corresponding isotony inclusion, and sends each admissible
charged sector to a finite-dimensional endomorphism/module category of
$\mathcal A(\Gamma)$.  In this finite regime the following DR hypotheses are
present: a rigid semisimple tensor category of sectors, conjugates for charged
modules, a symmetric exchange rule for independent spacelike interfaces, and a
compact internal automorphism group reconstructed as
$\mathrm{Inn}(\mathcal A)$.
\end{proposition}

\begin{proof}[Proof sketch and hypothesis audit]
Semisimplicity and finite rigidity follow from the Wedderburn decomposition of
$\mathcal A(\Gamma)$.  Tensor product is the APF composition of independent
interfaces; conjugates are provided by dual modules.  The local-support lemma
gives symmetric interchange for independent finite interfaces, while the
continuum/type-III section records the extra hypotheses needed to pass to a
Haag--Kastler net: nuclearity/split estimates, Haag duality in the scaling
limit, and convergence of finite local algebras to type-III local algebras.
Thus the finite APF result is not claimed to be a full DHR theorem in the
continuum.  It is a finite, cost-equipped precursor whose robust output is the
compact internal group and its representation category.  The assumptions that
remain external to the finite theorem are exactly the continuum analytic
hypotheses, not the finite inner-automorphism calculation.
\end{proof}

\begin{center}\footnotesize
\renewcommand{\arraystretch}{1.12}
\begin{tabular}{@{}p{3.5cm}p{5.0cm}p{5.0cm}@{}}
\toprule
\textbf{DR/DHR datum} & \textbf{APF finite counterpart} & \textbf{Continuum status} \\
\midrule
Local observable/field net & Cost-equipped admissibility field algebra $\mathcal A(\Gamma)$ and invariant subalgebra $\mathcal A(\Gamma)^G$ & Requires scaling limit and Haag--Kastler regularity. \\
Selection sectors & Admissible charged modules/continuation sectors & Robust as representation category; analytic localization is continuum-dependent. \\
Tensor product & Independent-interface composition & Finite version exact; continuum version requires split/nuclearity hypotheses. \\
Compact gauge group & $\mathrm{Inn}(\mathcal A)=\prod_k \mathrm{PU}(n_k)$ plus central grading & Finite output; global topology refined by matter-module kernel. \\
Duality/reconstruction & T3 plus Theorem~2.11 of \TSII{} & Full DHR theorem imported only in the continuum bridge. \\
\bottomrule
\end{tabular}
\end{center}

Pinzari's program (2017--) extends DR to \emph{braided} tensor
categories and weak Hopf $C^*$-algebras.  In this setting, abelian
factors can emerge from the \emph{center} of the fusion category:
central objects (transparent to braiding) generate grading tori that
act as abelian gauge factors.  This is structurally parallel to the
APF's mechanism for $\U(1)_Y$: the relative-phase torus
$T_{\mathrm{rel}} = Z(\mathrm{U}(\mathcal{A})) /
\mathrm{U}(1)_{\mathrm{diag}}$ (Remark~2.18 of \TSII{}) plays the role
of Pinzari's grading torus, and the abelian factor arises from
central (= transparent) elements rather than from inner automorphisms.

The key structural parallel is: \emph{non-abelian gauge from inner
automorphisms, abelian gauge from the center}.  In DR/Pinzari, this
is ``gauge from the Tannakian group, grading from the M\"uger center.''
In APF, it is ``gauge from $\mathrm{Inn}(\mathcal{A})$, grading from
$T_{\mathrm{rel}}$ pruned by CM.''  The APF adds a \emph{selection
mechanism} (capacity minimality, Proposition~\ref{prop:unique_U1})
that is absent in the DR/Pinzari framework: the latter reconstructs
whatever symmetry the category has, while the APF derives the
\emph{minimal} symmetry compatible with admissibility completeness.


\smallskip\noindent\textit{Hopf, weak-Hopf, VOA, and quantum-group reconstructions.}\quad
Weak-Hopf, quantum-group, VOA, and categorical reconstruction programs allow
more general ``quantum symmetry'' than ordinary compact Lie-group inner
automorphisms.  APF intentionally restricts the finite Standard-Model core to
ordinary inner automorphisms of a semisimple admissibility algebra because the
objects being represented are classical-stable record distinctions at a finite
interface.  More general quantum symmetries would require additional braided
or fusion data in the carrier.  Such extensions are not forbidden by APF, but
they are not part of the minimum finite-interface theorem; they would be
higher-structure additions whose admissibility cost must be justified separately.

\smallskip\noindent\textbf{Finite tensor categories and symmetric
tensor categories.}\quad
Recent work on finite tensor categories (Etingof--Gelaki--Nikshych--Ostrik)
and symmetric tensor categories (Deligne's theorem and extensions by
Coulembier, Entova-Aizenbud, Ostrik) shows that structural constraints
on representation categories can have sharp classification
consequences.  Deligne's theorem, for example, classifies all
symmetric tensor categories of moderate growth as representation
categories of supergroups --- a result that constrains the possible
internal symmetry structures of any theory whose observables form a
symmetric tensor category.

The APF is not a tensor-categorical framework: its primary objects are
admissibility algebras and cost functionals, not abstract tensor
categories.  However, the representation category of the admissibility
algebra $\mathcal{A} \cong \bigoplus_k M_{n_k}(\mathbb{C})$
\emph{is} a finite semisimple tensor category (the category of
finite-dimensional $\mathcal{A}$-modules), and its symmetry is
classified by the theorems above.  The APF's contribution relative to
this literature is not a new classification of tensor categories but a
\emph{physical derivation} of which tensor category is realized: the
admissibility postulate A1 selects $\mathcal{A}$ and hence its
representation category, rather than postulating the category
directly.  The classification theorems then apply as \emph{imported
mathematics} (same epistemic status as Wedderburn--Artin or
Skolem--Noether in the present supplement).



\paragraph{Relation to $L_\infty/C_\infty$ and homotopy-algebraic gauge formalisms.}\label{par:homotopy_algebra}
Homotopy-algebraic approaches encode the consistency identities of gauge
theories---gauge closure, higher brackets, deformation theory, and the
organization of interactions---once the field complex and symmetry data are
specified.  The APF is aimed one step upstream.  The finite admissibility algebra itself
is not asserted to carry a primitive $L_\infty$ or $C_\infty$ structure.
Rather, once APF has fixed the field complex and gauge template, the usual
BV/BRST or homotopy-algebraic lift may be constructed to organize interactions
and gauge identities.  In this sense the two programs are complementary:
homotopy algebras organize consistent dynamics, while APF attempts to derive
the minimum-cost admissible input data.
\smallskip\noindent\textbf{Topological charge quantization and $\mathbb{CP}^k$ sigma-model routes.}\quad
Topological constructions can quantize charges by assigning matter to bundles or soliton sectors and imposing integrality of topological classes.  APF's hypercharge quantization is different in scope.  It does not use a background $\mathbb{CP}^k$ target space or a topological winding number; it first fixes the finite admissibility algebra and chiral template, then applies the standard anomaly equations to obtain a unique rational ray.  Consequently APF's additional claim is not merely charge quantization, but exclusion of redundant abelian directions within the minimal finite template.  If a topological route introduces extra sectors or sterile records, APF classifies that as an enlarged admissibility problem rather than a counterexample to the minimal theorem.

\smallskip\noindent\textbf{EFT anomaly routing and Goldstone contributions.}\quad
Effective-field-theory analyses remind us that anomaly terms can be redistributed among current definitions, Wess--Zumino terms, and Goldstone contributions.  The present supplement uses anomaly cancellation only at the renormalizable chiral-template level: F6 checks the UV gauge consistency of the candidate fermion spectrum before EFT operators or symmetry-breaking field redefinitions are introduced.  APF therefore does not infer new low-energy anomaly coefficients from the ledger; it uses the standard exact anomaly equations as a consistency filter on the APF-selected template.



\section{Code--theorem concordance}
\label{app:code}
% ══════════════════════════════════════════════════════════════

Every theorem in this supplement --- and every chain theorem it
consumes from \TSII\ --- is verified by a named check function in the
codebase.  The table below provides the mapping; rows whose section
column points into \TSII\ are consumed inputs, banked under the same
verification contract.

\smallskip
\begin{center}\small
\renewcommand{\arraystretch}{1.15}
\resizebox{\textwidth}{!}{\begin{tabular}{@{}llll@{}}
\toprule
\textbf{Theorem} & \textbf{Supplement \S} & \textbf{Function} & \textbf{File} \\
\midrule
$L_{\mathrm{Cauchy}}$ & \S1 of \TSII{} & \texttt{check\_L\_Cauchy\_uniqueness} & \texttt{L\_Cauchy\_uniqueness.py} \\
Gauge ident. & \S2.3 of \TSII{} & \texttt{check\_T3} & \texttt{core.py} \\
$T_{7B}$ & \S3.1 of \TSII{} & \texttt{check\_T7B} & \texttt{gravity.py} \\
$L_{\mathrm{irr,uniform}}$ & \S3.2 of \TSII{} & \texttt{check\_L\_irr\_uniform} & \texttt{core.py} \\
$B1'$ & \S3.3 of \TSII{} & \texttt{check\_B1\_prime} & \texttt{gauge.py} \\
$L_{\mathrm{gauge,unique}}$ & \S\ref{sec:gauge_class_supp} & \texttt{check\_L\_gauge\_template\_uniqueness} & \texttt{gauge.py} \\
$T_{\mathrm{field}}$ & \S\ref{sec:fermion_scan} & \texttt{check\_T\_field}; \texttt{fermion\_scan\_standalone.py}~\cite{BrookeFermionScan} & \texttt{gauge.py} \\
$L_{\mathrm{hypercharge}}$ & \S\ref{sec:hypercharge} & \texttt{check\_T\_field} (hypercharges) & \texttt{gauge.py} \\
$L_{\mathrm{count}}$ & \S\ref{sec:capacity_count} & \texttt{check\_L\_count} & \texttt{gauge.py} \\
$I2$ (gauge-cosmo) & \S\ref{sec:I2_at_L_count} & \texttt{check\_I2\_gauge\_cosmological}, \texttt{check\_T\_ACC\_unification} & \texttt{unification.py} \\
$P_{\mathrm{exhaust}}$ & exited (Paper~8; \S\ref{sec:bridge_axioms}) & \texttt{check\_P\_exhaust} & \texttt{core.py} \\
$L_{\mathrm{equip}}$ & exited (Paper~41; \S\ref{sec:bridge_axioms}) & \texttt{check\_L\_equip} & \texttt{cosmology.py} \\
$L_{\mathrm{sat,part}}$ & exited (Paper~8; \S\ref{sec:bridge_axioms}) & \texttt{check\_L\_saturation\_partition} & \texttt{cosmology.py} \\
$\Delta_{\mathrm{sig}}$ & exited (Paper~6; \S\ref{sec:bridge_axioms}) & \texttt{check\_Delta\_signature} & \texttt{spacetime.py} \\
$T_8$ ($d{=}4$) & exited (Paper~6; \S\ref{sec:bridge_axioms}) & \texttt{check\_T8} & \texttt{spacetime.py} \\
$T_{9,\mathrm{grav}}$ & exited (Paper~6; \S\ref{sec:bridge_axioms}) & \texttt{check\_T9\_grav} & \texttt{gravity.py} \\
$T_{\theta_{\mathrm{QCD}}}$ & exited (Paper~46; \S\ref{sec:bridge_axioms}) & \texttt{check\_T\_theta\_QCD} & \texttt{gauge.py} \\
$T_{\mathrm{vac,stab}}$ & exited (Paper~45; \S\ref{sec:bridge_axioms}) & \texttt{check\_T\_vacuum\_stability} & \texttt{gauge.py} \\
$T_{\mathrm{proton}}$ & exited (bank; BA8 in \S\ref{sec:bridge_axioms}) & \texttt{check\_T\_proton} & \texttt{gauge.py} \\
\bottomrule
\end{tabular}}
\end{center}

\smallskip\noindent The full codebase is available at
\texttt{\seqsplit{github.com/Ethan-Brooke/APF-Paper-2-The-Structure-of-Admissible-Physics}}.

\smallskip\noindent\textit{Reproducibility of the fermion scan.}\quad
The function \texttt{check\_T\_field} in \texttt{apf/gauge.py}
enumerates the template space in ordered-tuple form (4{,}680
iterations over the 1{,}680 distinct templates) and applies the seven
filters under the uniform-dimension F6 variant --- the 4/2 robustness
path of the predicate taxonomy (\S\ref{sec:seven_filters}) ---
verifying that the unique minimum-DOF survivor is the SM at 45 Weyl
fermions.  It records no per-stage waterfall counts.  The canonical
full-system waterfall (10/5; every count in \S\ref{sec:waterfall}) is
generated by \texttt{check\_L\_F6\_not\_from\_EC} in
\texttt{apf/ec\_inventory\_reading.py} and by the release's audit
files.  The exclusion-proof assertions (P1--P5) are verified within
\texttt{check\_T\_field}; hypercharge derivation and anomaly
cancellation are checked algebraically in exact rational arithmetic
(\texttt{fractions.Fraction}).  Running
\texttt{python verify\_all.py} from the repository root runs the full
check registry (both functions included) with no dependencies beyond
the Python standard library.

\smallskip\noindent\textit{Standalone verification script.}\quad
The file \texttt{fermion\_scan\_standalone.py}~\cite{BrookeFermionScan},
included in the GitHub repository and archived on Zenodo
(\url{https://doi.org/10.5281/zenodo.19154197}),
is a self-contained Python script ($\sim$550 lines, zero dependencies
beyond the standard library) that reproduces the complete scan from
scratch.  An interactive Colab walkthrough with waterfall visualisation
is available at
\url{https://colab.research.google.com/drive/1otO1Qwp_Lr3OsK6ykkeMRATanecPkKg6}.
Running \texttt{python3 fermion\_scan\_standalone.py} performs the following:
\begin{enumerate}[nosep,label=(\arabic*)]
  \item Enumerates all 1{,}680 distinct templates using
  \texttt{combinations\_with\_replacement}.
  \item Applies all 7 filters with exact rational arithmetic, including
  spectator-field reduction for $A(R) = 0$ representations.
  \item Prints the waterfall table with exact counts at each stage.
  \item Lists all 8 post-F6 survivors under its spectator-reduction
  predicate (the robustness form; the canonical full-system listing in
  \S\ref{sec:waterfall} carries 10).
  \item Lists all 4 post-CPT survivors (spectator-reduction form; the
  canonical count is 5).
  \item Verifies all 5 exclusion proofs (P1--P5) with explicit arithmetic.
  \item Derives both hypercharge solutions ($z = 1 \pm N_c$),
  verifies all 4 anomaly conditions exactly for each, and confirms
  their physical equivalence under $u \leftrightarrow d$ relabelling.
  \item Runs 4 built-in red-team challenges: (RT1) combinatorial
  cross-check of the enumeration count, (RT2) explicit spectator-field
  anomaly verification, (RT3) dead-code reachability audit, (RT4) the
  two-colored-doublet class-dominance enumeration (class minimum 54
  Weyl DOF, witness-verified; Proposition~\ref{prop:P4}).
  \item Asserts 12 final verification conditions.
\end{enumerate}
The script is designed so that any reader can verify the entire
fermion-content derivation on a standard laptop in under one second,
with no knowledge of the APF codebase required.

The Lie-algebra classification (\S\ref{sec:gauge_class_supp}) is verified by
\texttt{check\_L\_gauge\_template\_uniqueness}, which tests the
classification's 17 representatives (the families closed by the
monotonicity arguments of \S\ref{sec:gauge_class_supp}) against
R1--R3 and records the exclusion reason for each.  The Zenodo archive
(DOI: 10.5281/zenodo.18867119) provides a frozen snapshot of the
codebase corresponding to this supplement.




% ══════════════════════════════════════════════════════════════
\section{Fermion scan: algorithmic walkthrough}
\label{app:scan_walkthrough}
% ══════════════════════════════════════════════════════════════

This appendix makes the fermion-content derivation
(\S\ref{sec:fermion_scan}) verifiable without running code.  It
specifies the search space construction, gives the algorithm in
pseudocode, and traces the critical filter steps on explicit
examples.

\subsection{Search space construction}
\label{app:search_space}

A template is a multiset of $(r^{(3)}, r^{(2)})$ pairs.  The
search space is bounded by P1--P5 (\S\ref{sec:exclusion_proofs}):

\smallskip\noindent\textbf{Allowed $\SU(3)$ reps} (P1: $\dim < 10$):
$\{3, \bar{3}, 6, \bar{6}, 8\}$.
Dynkin indices: $T(3) = T(\bar{3}) = 1/2$, $T(6) = T(\bar{6}) = 5/2$,
$T(8) = 3$.  Cubic anomaly coefficients: $A(3) = 1/2$,
$A(\bar{3}) = -1/2$, $A(6) = 7/2$, $A(\bar{6}) = -7/2$, $A(8) = 0$.
(The quadratic \emph{index} $T(6) = 5/2$ and the cubic \emph{anomaly}
$A(6) = 7/2$ are distinct quantities; v3.9's appendix carried $5/2$
for both, a holdover from before the v3.7 correction.  The value is
computationally inert for the waterfall: sextet-containing templates
do reach F4, but no F4 verdict differs between $A(6) = 5/2$ and $7/2$
anywhere in the reachable space --- verified by direct recomputation
of all 144 sextet templates passing F1--F3, 2026-07-13.  The v3.7
changelog's stated reason, that sextets are AF-eliminated before F4,
was incorrect.)

\smallskip\noindent\textbf{Allowed $\SU(2)$ reps} (P2 for colored
fields; P3 dominance for colorless triplets): $\{1, 2\}$.  Dynkin indices: $T(1) = 0$, $T(2) = 1/2$.

\smallskip\noindent\textbf{Template structure} (P4--P5):
Each template contains exactly one colored doublet (P4: two colored
doublets force $\ge 54$ DOF, class dominance), $0$--$3$ colored
singlets, $0$--$1$
lepton doublets, and $0$--$2$ lepton singlets (P5: extra field types
beyond the anomaly-closed core cannot lower the minimum ---
dominance).

\smallskip\noindent\textbf{Enumeration count.}\quad
The colored doublet ranges over 5 reps.  The colored singlets form a
multiset of size $0$--$3$ drawn from 5 reps: $\binom{5+0-1}{0} +
\binom{5+1-1}{1} + \binom{5+2-1}{2} + \binom{5+3-1}{3} = 1 + 5 + 15
+ 35 = 56$ multisets.  The lepton doublet is present or absent: 2
choices.  The lepton singlets: $0$, $1$, or $2$ copies of
$(\mathbf{1},\mathbf{1})$: 3 choices.  Total:
$5 \times 56 \times 2 \times 3 = 1{,}680$ distinct templates.
The standalone verification script~\cite{BrookeFermionScan}
uses \texttt{combinations\_with\_replacement} to enumerate these
directly, and verifies the count against this closed-form formula
as a built-in red-team challenge (RT1).

\subsection{Algorithm in pseudocode}
\label{app:pseudocode}

\begin{small}
\begin{verbatim}
INPUT:  SU3_reps = {3, 3b, 6, 6b, 8}
        SU2_reps = {1, 2}
        N_gen = 3
        AF3_bound = 11,  AF2_bound = 22/3
        AF_coeff = 2/3

ENUMERATE all templates T:
  FOR each colored_doublet IN SU3_reps:           // 5
    FOR n_colored_singlets IN {0, 1, 2, 3}:
      FOR each combo IN C(SU3_reps, n_cs):        // combinations w/ replacement
        FOR has_lepton_doublet IN {True, False}:   // 2
          FOR n_lepton_singlets IN {0, 1, 2}:      // 3
            T = [(colored_doublet, 2)]
                + [(c, 1) for c in combo]
                + [(1, 2)] if has_lepton_doublet
                + [(1, 1)] * n_lepton_singlets
            total_tested += 1

            // F1+F2: Asymptotic freedom (exact rational)
            su3_cost = SUM over (a,b) in T:
                         T_3(a) * dim_2(b) * N_gen
            su2_cost = SUM over (a,b) in T:
                         T_2(b) * dim_3(a) * N_gen
            IF AF3_bound - AF_coeff * su3_cost <= 0: SKIP
            IF AF2_bound - AF_coeff * su2_cost <= 0: SKIP

            // F3: Doublet-singlet content
            IF no colored doublet or no colored singlet: SKIP

            // F4: [SU(3)]^3 anomaly cancellation
            IF SUM A_3(a) * dim_2(b) != 0: SKIP

            // F5: Witten anomaly
            IF SUM dim_3(a) [where b=2] is odd: SKIP

            // F6: Full anomaly system
            // (rational root test on the hypercharge
            //  quadratic -- details in §8)
            IF no rational hypercharge solution: SKIP

            // F7: CPT quotient
            canonical = min(sorted(T),
                           sorted(conjugate(T)))
            IF canonical already seen: SKIP
            Mark canonical as seen

            // Record survivor
            DOF = SUM dim_3(a) * dim_2(b) * N_gen
            survivors.append((DOF, T))

SELECT minimum-DOF survivor(s)
\end{verbatim}
\end{small}


\subsection{Filter-by-filter trace on critical examples}
\label{app:filter_trace}

We trace three templates through all seven filters: the SM (passes
all), a near-miss (fails at F6), and a common failure (fails at F4).

\smallskip\noindent\textbf{Example 1: SM template}
$\{(3,2),\, (\bar{3},1),\, (\bar{3},1),\, (1,2),\, (1,1)\}$.

\smallskip\noindent\textit{F1 (SU(3) AF):}\quad
$b_0^{(3)} = 11 - \tfrac{2}{3}\bigl[\tfrac{1}{2}\!\cdot\!2 +
\tfrac{1}{2}\!\cdot\!1 + \tfrac{1}{2}\!\cdot\!1\bigr]\!\cdot\!3
= 11 - \tfrac{2}{3}\!\cdot\!2\!\cdot\!3 = 11 - 4 = 7 > 0$.  \checkmark

\smallskip\noindent\textit{F2 (SU(2) AF):}\quad
$b_0^{(2)} = \tfrac{22}{3} - \tfrac{2}{3}\bigl[\tfrac{1}{2}\!\cdot\!3
+ \tfrac{1}{2}\!\cdot\!1\bigr]\!\cdot\!3
= \tfrac{22}{3} - \tfrac{2}{3}\!\cdot\!2\!\cdot\!3 = \tfrac{22}{3} - 4
= \tfrac{10}{3} > 0$.  \checkmark

\smallskip\noindent\textit{F3 (doublet--singlet content):}\quad
$(3,2)$ is a colored doublet; $(\bar{3},1)$ is a colored singlet.
Both present.  \checkmark

\smallskip\noindent\textit{F4 ($[\SU(3)]^3$):}\quad
$\tfrac{1}{2}\!\cdot\!2 + (-\tfrac{1}{2})\!\cdot\!1
+ (-\tfrac{1}{2})\!\cdot\!1 = 1 - \tfrac{1}{2} - \tfrac{1}{2} = 0$.
\checkmark

\smallskip\noindent\textit{F5 (Witten):}\quad
Doublets weighted by $\SU(3)$ dim: $3 + 1 = 4$ (even).  \checkmark

\smallskip\noindent\textit{F6 (full anomaly):}\quad
$N_c = 3$, $n_{\mathrm{cs}} = 2$, $n_{\mathrm{ls}} = 1$.
Discriminant $= 4 + 4(9-1) = 36$.  $\sqrt{36} = 6 \in \mathbb{Z}$.
\checkmark

\smallskip\noindent\textit{F7 (CPT):}\quad Canonical form
$= \min(\text{sorted}(T), \text{sorted}(\overline{T}))$.
$\overline{T} = \{(\bar{3},2),\, (3,1),\, (3,1),\, (1,2),\, (1,1)\}$.
Sorted: $(1,1)(1,2)(3,2)(\bar{3},1)(\bar{3},1)$ vs.\
$(1,1)(1,2)(3,1)(3,1)(\bar{3},2)$.  First is lexicographically
smaller.  Canonical = SM.  \checkmark

\smallskip\noindent DOF $= (6+3+3+2+1)\times 3 = 45$.
\textbf{Passes all filters.}

\bigskip
\noindent\textbf{Example 2: $\{(6,2),\, (\bar{6},1),\, (1,2),\, (1,1)\}$}
(a template with the sextet).

\smallskip\noindent\textit{F1:}\quad
$b_0^{(3)} = 11 - \tfrac{2}{3}[\tfrac{5}{2}\!\cdot\!2 +
\tfrac{5}{2}\!\cdot\!1]\!\cdot\!3 = 11 - \tfrac{2}{3}\!\cdot\!
\tfrac{15}{2}\!\cdot\!3 = 11 - 15 < 0$.  $\boldsymbol{\times}$ \textbf{Fails F1.}

\bigskip
\noindent\textbf{Example 3: $\{(3,2),\, (3,1),\, (1,2),\, (1,1)\}$}
(wrong-sign colored singlet).

\smallskip\noindent\textit{F1--F3:}\quad Pass (verify: $b_0^{(3)} = 11
- \tfrac{2}{3}\!\cdot\!\tfrac{3}{2}\!\cdot\!3 = 8 > 0$; both content
types present).

\smallskip\noindent\textit{F4:}\quad
$\tfrac{1}{2}\!\cdot\!2 + \tfrac{1}{2}\!\cdot\!1 = \tfrac{3}{2}
\ne 0$.  $\boldsymbol{\times}$ \textbf{Fails F4.}  The cubic anomaly is not
cancelled because both colored fields are in the $\mathbf{3}$ (same
sign of $A$).

\bigskip
\noindent\textbf{Example 4: $\{(3,3)_{0},\, (\bar{3},1)_{+1},\,
(\bar{3},1)_{-1},\, (\bar{3},1)_{0}\}$} (the colored weak-triplet content
--- the sharpest known rival, 18 Weyl fermions as a single copy).

\smallskip\noindent This content is chiral, satisfies every gauge and
gravitational anomaly condition exactly, satisfies Witten parity (no
doublets), and as a \emph{single copy} is asymptotically free in both
non-abelian factors --- it is the strongest stress test of the scan's
completeness, and it is excluded at every multiplicity by three independent
banked routes.  At $n = 1$ or $2$ it is an under-replicated rival, killed by
the max-saturated selection \coderef{check\_R\_Ngen\_neq\_3\_killed}{killed\_rivals.py}
and by the capacity--beta identities, which no $n \ne 3$ satisfies
\coderef{check\_L\_beta\_capacity}{generations.py}.  At $n \ge 2$ it
destroys $\SU(2)$ asymptotic freedom outright ($(2/3)(2)(3)(2) = 8 > 22/3$;
at the forced $n = 3$, $12 > 22/3$ --- Proposition~P2's arithmetic).  And its
own capacity--beta system has \emph{no} solution at any $n$: the non-abelian
identities give $n = 27/13$ and $n = 23/10$, non-integer, with the
convention-free $6|b_3|$ equation failing every $n$ by a margin of at least
$18$.  The Standard-Model template satisfies all three identities
simultaneously at $n = 3$; this rival satisfies none at any $n$.

\subsection{Completeness argument}
\label{app:completeness}

The scan's boundedness rests on three facts:

\begin{enumerate}[nosep]
\item \textbf{Representation bound (P1--P2, theorems).}\quad Every
$\SU(3)$ rep with $\dim \ge 10$ and every colored $\SU(2)$ rep with
$\dim \ge 3$ exceeds the AF budget as a \emph{single field} (the cost
of one such field, summed over 3 generations, exceeds $b_0$).  These
are closed-form inequalities verified in \S\ref{sec:exclusion_proofs}.
Therefore the allowed reps are exactly
$\SU(3) \in \{3,\bar{3},6,\bar{6},8\}$ and $\SU(2) \in \{1,2\}$.

\item \textbf{Multiplicity caps (declared inputs; P3--P5 dominance).}\quad
At most 1 colored doublet (P4: two colored doublets force $\ge 54$
DOF --- the class-dominance enumeration of
Proposition~\ref{prop:P4}), at most 3 colored singlets,
at most one lepton doublet, at most 2 lepton singlets --- template size
$|\mathcal{T}| \le 7$.  These caps are declared enumeration inputs
(\S\ref{sec:exclusion_proofs}).  P3 and P5 are dominance results:
within the declared space, every survivor carrying types beyond the
five SM roles extends an anomaly-closed core and sits strictly above
it in DOF, so the enumerated enlargement directions cannot lower the
minimum.

\item \textbf{Complete enumeration within the caps.}\quad Within the
bounded set ($5 \times 56 \times 2 \times 3 = 1{,}680$ distinct
templates), every combinatorial possibility is generated and tested.
No branch is skipped.  The enumeration is deterministic and
order-independent (the final result depends only on the filter
outcomes, not on the order of testing).
\end{enumerate}

\noindent Together, (1)--(3) establish the licensed claim: every
chiral fermion template within the declared representation bounds and
multiplicity caps is either in the 1{,}680-template set or is excluded
by a closed-form proof, and the SM at 45 Weyl DOF is the unique
minimum \emph{among the enumerated candidates}.  Completeness of the
caps themselves --- that no admissible sub-45-DOF template exists
outside them --- is the open proof obligation of the reproducibility
contract (\S\ref{sec:fermion_scan}), carried by P3/P5-style dominance
for the enumerated directions and otherwise open.

\subsection{Robustness checks}
\label{app:robustness}

\begin{enumerate}[nosep]
\item \textbf{Exact arithmetic.}\quad All filter computations use
Python's \texttt{fractions.Fraction} (exact rational arithmetic).  No
floating-point approximation is made at any stage.  The anomaly
conditions are checked as exact equalities, not as
$|\text{residual}| < \epsilon$.

\item \textbf{Filter-order independence.}\quad The final survivor set
is determined by the conjunction of all seven filters.  The sequential
application (F1 before F2 before F3, etc.)\ is an optimization
(early termination), not a logical dependency.  Reordering the filters
produces the same final set.  The standalone script
(\texttt{fermion\_scan\_standalone.py}) applies all filters and records
intermediate counts; permuting the filter order and comparing outputs
verifies independence.

\item \textbf{CPT symmetry.}\quad The CPT quotient (F7) identifies
each template with its charge conjugate.  The conjugation map
$3 \leftrightarrow \bar{3}$, $6 \leftrightarrow \bar{6}$,
$8 \to 8$, $1 \to 1$ is an involution.  The canonical representative
is the lexicographic minimum of the pair $\{T, \overline{T}\}$.  The
10 post-F6 survivors under the canonical full-system predicate form 5
CPT pairs (listed in \S\ref{sec:waterfall}); the standalone script's
spectator-reduction predicate yields 8, forming 4 pairs.

\item \textbf{Sensitivity to generation count.}\quad The scan assumes
$N_{\mathrm{gen}} = 3$ (derived in Paper~2, \S9).  With
$N_{\mathrm{gen}} = 1$ or $2$: the AF bounds loosen ($b_0$ terms
scale as $N_{\mathrm{gen}}$), admitting additional survivors.  With
$N_{\mathrm{gen}} \ge 4$: the AF bounds tighten, reducing or
eliminating survivors.  The $N_{\mathrm{gen}} = 3$ scan is therefore
not vacuous (it admits templates) and not trivial (it excludes most).  The looser $N_{\mathrm{gen}} \in \{1,2\}$ survivor sets are not an open door: those generation counts are themselves excluded (the capacity--beta identities and the rival kill, \S\ref{sec:three_generation_scope}), so their would-be survivors are killed at the multiplicity level before content filtering arises.

\item \textbf{Second survivor.}\quad The second CPT-inequivalent
survivor (DOF $= 48$, SM $+$ extra singlet) provides a
falsifiability test: if the SM contained 48 Weyl fermions (e.g., with
a right-handed neutrino), the scan's minimum would move off the
Standard-Model template, and the count would move to $C = 64$ ---
whose bridge-conditional comparison with observation (BA6) lives at
the corpus homes (\S\ref{sec:neutrinos}).
\end{enumerate}


% ══════════════════════════════════════════════════════════════
% BIBLIOGRAPHY (placeholder)
% ══════════════════════════════════════════════════════════════

\begin{thebibliography}{99}

\bibitem{BrookePaper1}
E.\,S.~Brooke, ``Paper~1: The Enforceability of Distinction,''
Admissibility Physics, 2025.
Submitted to \textit{Foundations of Physics}.

\bibitem{BrookeSupplement}
E.\,S.~Brooke, ``Technical Supplement: End-to-End Derivation from Finite
Enforcement Capacity to Quantum Structure,'' companion to Paper~1, 2025.

\bibitem{Humphreys1972}
J.\,E.~Humphreys, \textit{Introduction to Lie Algebras and Representation
Theory}, Springer, 1972.

\bibitem{FultonHarris1991}
W.~Fulton and J.~Harris, \textit{Representation Theory: A First Course},
Springer, 1991.

\bibitem{Witten1982}
E.~Witten, ``An SU(2) Anomaly,'' \textit{Phys.\ Lett.\ B} \textbf{117},
324--328 (1982).

\bibitem{PeskinSchroeder1995}
M.\,E.~Peskin and D.\,V.~Schroeder, \textit{An Introduction to Quantum
Field Theory}, Westview Press, 1995.

\bibitem{GrossWilczek1973}
D.\,J.~Gross and F.~Wilczek, ``Ultraviolet Behavior of Non-Abelian
Gauge Theories,'' \textit{Phys.\ Rev.\ Lett.} \textbf{30}, 1343 (1973).

\bibitem{Politzer1973}
H.\,D.~Politzer, ``Reliable Perturbative Results for Strong Interactions?''
\textit{Phys.\ Rev.\ Lett.} \textbf{30}, 1346 (1973).

\bibitem{Planck2020}
Planck Collaboration, ``Planck 2018 results.\ VI.\ Cosmological
parameters,'' \textit{Astron.\ Astrophys.} \textbf{641}, A6 (2020).

\bibitem{PDG2024}
Particle Data Group, ``Review of Particle Physics,''
\textit{Phys.\ Rev.\ D} \textbf{110}, 030001 (2024).

\bibitem{Hardy2001}
L.~Hardy, ``Quantum theory from five reasonable axioms,''
arXiv:quant-ph/0101012 (2001).

\bibitem{CDP2011}
G.~Chiribella, G.\,M.~D'Ariano, and P.~Perinotti,
``Informational derivation of quantum theory,''
\textit{Phys.\ Rev.\ A} \textbf{84}, 012311 (2011).

\bibitem{Hohn2017}
P.\,A.~H\"ohn, ``Quantum theory from questions,''
\textit{Phys.\ Rev.\ A} \textbf{95}, 012102 (2017).

\bibitem{Connes1996}
A.~Connes, ``Gravity coupled with matter and the foundation of
noncommutative geometry,''
\textit{Commun.\ Math.\ Phys.} \textbf{182}, 155--176 (1996).

\bibitem{Bousso1999}
R.~Bousso, ``A covariant entropy conjecture,''
\textit{JHEP} \textbf{07}, 004 (1999).

\bibitem{HKM1976}
S.\,W.~Hawking, A.\,R.~King, and P.\,J.~McCarthy,
``A new topology for curved space-time which incorporates the causal,
differential, and conformal structures,''
\textit{J.\ Math.\ Phys.} \textbf{17}, 174--181 (1976).

\bibitem{Malament1977}
D.\,B.~Malament, ``The class of continuous timelike curves determines
the topology of spacetime,''
\textit{J.\ Math.\ Phys.} \textbf{18}, 1399--1404 (1977).

\bibitem{Kolmogorov1933}
A.\,N.~Kolmogorov, \textit{Grundbegriffe der Wahrscheinlichkeitsrechnung},
Springer, 1933.

\bibitem{Lovelock1971}
D.~Lovelock, ``The Einstein tensor and its generalizations,''
\textit{J.\ Math.\ Phys.} \textbf{12}, 498--501 (1971).

\bibitem{BrookeFermionScan}
E.\,S.~Brooke, ``APF Fermion Template Scan---Standalone Verification
Script (v2),'' 2026.
Zenodo: \url{https://doi.org/10.5281/zenodo.19154197}.
GitHub: \url{https://github.com/Ethan-Brooke/APF-Paper-2-The-Structure-of-Admissible-Physics}.
Interactive walkthrough:
\url{https://colab.research.google.com/drive/1otO1Qwp_Lr3OsK6ykkeMRATanecPkKg6}.

\bibitem{DonnellyGiddings2016PRD93}
W.~Donnelly and S.\,B.~Giddings, ``Diffeomorphism-invariant observables
and their nonlocal algebra,'' \textit{Phys. Rev. D} \textbf{93}, 024030
(2016).

\bibitem{DonnellyGiddings2016PRD94}
W.~Donnelly and S.\,B.~Giddings, ``Observables, gravitational dressing,
and obstructions to locality and subsystems,'' \textit{Phys. Rev. D}
\textbf{94}, 104038 (2016).

\bibitem{SanchezQuevedo2023}
L.\,M.~S\'anchez and H.~Quevedo, ``Thermodynamics of the FLRW apparent
horizon,'' \textit{Phys. Lett. B} \textbf{839}, 137778 (2023).

\end{thebibliography}


\end{document}
